% SAME-FAMILY GENERAL COMPACT SOURCE DRAFT, 25 October2026.
% Exact19 whole A/B/D consumers; no new paper, canonical or main edit.
% No arithmeticity, field/integrality, enumeration, imports, tests or builds.
% DRAFT pending complete different-author/root review and shared-tail closure.
\section{The whole classical forests at exact dimension nineteen}
\label{app:exact19:wholeclassical}

These are three actual-rank consumers for
the general compact argument.
Every point and pressure below belongs
to one original farthest vertex \(p\),
chosen from one strict interior point.
We keep its single decomposition
\(z=E(p-x)\), \(x\in p^\perp\),
\(\|x\|<1\).
No arithmetic restriction is used.

\subsection{A whole \(A_{19}\): the window and the three reference shapes}
\label{app:exact19:Awindow}

Use the zero-sum subspace of
\(\mathbb R^{20}\), with actual old roots
and duals
\[
 n_i=(e_i-e_{i+1})/\sqrt2,\quad
 \Omega_i=\sqrt2(P_i-iP_{20}/20),\quad
 D_i=\Omega_i^2=i(20-i)/10 .
\]
Write \(h_i=x\cdot\Omega_i\), \(h_0=h_{20}=0\).
The sequence is strictly concave and
\[
                 \|x\|^2=\frac12\sum_{i=0}^{19}(h_{i+1}-h_i)^2 .
\]
A central site \(5\le i\le15\) is weak
when \(h_i\le\sqrt{D_i}-.61\), equality
included.

We use the shared complete \(.61\) window
of Lemma~\ref{app:tail:lem:Awindow2021}
at \(N=20\), before any contrary
component-energy bound: every weak cap
is pure, with no spectator projection,
and is either
\[
 v_i=g\Omega_i+\Omega_e/2,\quad e=1\text{ or }19,
 \quad \cos(\pi/5)\le g<27/25 ,
\]
or a near-four cap at \(i=5,6\) or
a reflected site.
The first row has own native five,
finite native at least seven, or dotted;
the smaller own labels and parallel case
have their actual affine or feasible-null
obstructions.
The literal helpers are nonadjacent.
These are the actual finite-white,
own-zero, foreign-clearing and mixed
two-cap conclusions of the shared
window proof, with retained paths of
rank at least four and full connected
rank nineteen; they do not follow by
applying a higher-rank whole-forest
exclusion here.
Its full internal proof precedes
the higher-tail energy consumers;
it is not a published classification
or an additional assumed field condition.

There is a useful stronger removal of
near four at this rank.
For a left near-four cap at \(i=5,6\),
its positive finite replacement at
\(j=i-1\) is an actual endpoint.
Its white forest is
\(B_{21-i}+A_{i-2}\); positive inverse
projection proves feasibility against
every original wall, including the
released \(j\).
The exact edge cosine is
\[
 \cosh L_j=\sqrt{2(i-1)/20}
                 \left\{(21-i)/2+(20-i)/\sqrt2\right\}>11 .
\]
For \(i=5\), use
\(.632(8+7.5(1.414))>11\);
for \(i=6\), use
\(.707(7.5+7(1.414))>11\).
Its normalized dual demand is
\(>\sqrt{5/6}>.91\).
At the opposite end, either a second
near-four mate has that same demand,
or helper nonadjacency gives a strong
\(k=14\) or fifteen, with demand
\(>1-.61/2.7>.774\).
The correct normalized correlation is
\[
 \rho=\sqrt{\frac{j(20-k)}{(20-j)k}}
                       \le1/\sqrt7<.38 .
\]
Cauchy would give
\((.91+.774)^2<2+2(.38)=2.76\);
instead \(1.684^2=2.835856\).
Thus near four is impossible under
\(\|x\|<1\), without a stronger
component assumption.

All remaining weak sites are therefore
nonadjacent pure helpers.
Their edge-ratio target is
\[
 R(q)=100/61-\frac{q-.61}{q(2q-.61)} .
\]
It increases for \(q\ge5/2\), while
\(q\le\sqrt{10}<16/5\) and
\(R(16/5)=3/2-251/565104<3/2\).
For side order \(m\ge5\),
\(\sigma=g+1/(2m)\), \(\delta=(m+1)/(2m)\),
and
\[
 \frac{\sigma^2}{\delta}-(4g-2)
                   =\frac{(g-1-1/(2m))^2}{\delta}\ge0 .
\]
Weakness consequently sharpens the
entire helper range to \(g<17/16\).
Native five has \(.809<g<.81\);
every other finite own has \(.9<g<1\).

Take rational symmetric heights
\[
 H_5=2.128,\ H_6=2.288,\ H_7=2.406,\
 H_8=2.488,\ H_9=2.536,\ H_{10}=2.552,\qquad
 H_{20-i}=H_i .
\]
They lie strictly below the strong
floors by positive rational squaring.
Linear interpolation on the tails
gives nonnegative dual coefficients.
Omit every actual internal weak knot
by replacing its height by the mean
of its neighbours.
The holes are nonadjacent, so their
open affected edge pairs are disjoint
and their exact quarter-square losses
add.
If an end site five or fifteen is
weak, flatten that entire end through
six or fourteen instead.
The complete reference data are
\[
\begin{array}{c|c|c|c}
 \text{weak end count}&E=\|w\|^2\text{ at least}
          &\text{end block lengths}&\text{individual capacities}\\ \hline
 0&.9544848-5(.000441)=.9522798&5,5&
             \beta_5,\beta_{15}\ge83/625\\
 1&3462479/3750000&6,5&
             \beta_6\ge79/600,\ \beta_{15}\ge83/625\\
 2&1341604/1500000&6,6&
             \beta_6,\beta_{14}\ge79/600 .
\end{array}
\tag{A19refs}\label{app:exact19:eq:Arefs}
\]
For count one orient the weak end to
the left by relabelling the old path.
This assumes no symmetry of the polytope.
There are at most five internal holes
in the first row, four in the other
rows.
The one-end row loses at most
\(4(.000441)\); the two-end row loses
at most \(4(.000324)\).
Their exact undeleted energies follow
from
\[
 a^2/5+(b-a)^2+T,\quad
 \tfrac12[b^2/6+a^2/5+(b-a)^2+2T],\quad
 b^2/6+T,
\]
where \(a=2.128\), \(b=2.288\) and
\(T=.118^2+.082^2+.048^2+.016^2\).
The boundary coefficients are
\[
 (a/5-(b-a))/2=83/625,\qquad
                    ((7/6)b-H_7)/2=79/600 .
\]
An adjacent hole increases these
capacities.
All positive coefficients are paid
by actual strong floors, so
\(x\cdot w\ge E\) in every allowed
hole pattern.
The constant signed means obey
\[
                  .3<b_5<.302,\qquad1/4<b_6<.27 .
\]

\subsection{The actual helper feet and their complete continuous bounds}
\label{app:exact19:Ahelperfeet}

For a weak site five, write
\(c=Ap-v\), \(v=g\Omega_5+\Omega_e/2\).
The fresh exact times are
\[
 A_{\rm far}^2=(300g^2+20g-21)/40,\qquad
 A_{\rm near}^2=(300g^2+60g-21)/40 .
\tag{A19times}\label{app:exact19:eq:Atimes}
\]
Both formulas concern \(N=20\).

The following constructions give actual
pure spatial pressures \(x\cdot U\ge H\).
At native five use the finite \(H_4\)
correction
\[
 k=g n_5+\tfrac23n_e+\tfrac13n_{\rm adj},\qquad
 \delta=2/3-g^2>0,\quad U=v-k .
\]
The normalized finite projected point
is obtained from \(p+(A/\delta)(c+k)\).
Its future time is positive, its square
is \(-1-A^2/\delta\), and all selected
products vanish.
Every omitted original-wall product
is nonpositive by the positive
normal coefficients and acuteness.
The pressure is
\[
 H=A+\delta/A-\sqrt{\delta(A^2+\delta)}/A
                   >A-\sqrt\delta>A-.111 .
\]
The spacelike correction \(c+k\)
alone is not a hyperbolic point.

For finite own at least seven,
put \(d=1-g^2\) and use
\(\xi=c+g n_5+2d n_e\).
Its selected products are zero,
zero and \(-1/2+2d<0\);
its square is \(-d(1-4d)<0\).
It is future and satisfies every
other original wall termwise.
Thus \(U=v-g n_5-2d n_e\),
\(H=A-\sqrt{d(1-4d)}\), with the
last radical at most \(1/4\).
For dotted own use
\(\xi=c+g^{-1}n_5\), of square
\(-(1-g^{-2})\), zero cap product
and negative own-root product.
It gives \(U=v-g^{-1}n_5\),
\(H=A-\sqrt{1-g^{-2}}\).
All vectors \(U\) are spatial in
the original \(p^\perp\).

A near helper forces \(x_1>.64\).
Indeed weakness gives \(h_5<2.129\),
and dropping the positive old-root
corrections from the pressure yields
\(x_1\ge\sqrt2(H-gh_5)\).
At native five, \(A>2.365\), giving
\[
              x_1>1.414(2.365-.111-.81(2.129))>.64 .
\]
For finite \(g\ge.9\),
\(A_{\rm near}'(g)>\sqrt{7.5}>2.7\);
thus \(A_{\rm near}-2.129g\) increases.
At \(.9\), \(A>2.626\), giving
\(1.414(2.626-.25-.9(2.129))=.6502986\).
For dotted \(1<g<17/16\),
\(A_{\rm near}'<3\), whereas
\[
                  \frac{d}{dg}\sqrt{1-g^{-2}}
                  =\frac1{g^3\sqrt{1-g^{-2}}}>2 .
\]
Here \(g^3<5/4\), the radical \(<2/5\).
The difference \(A_{\rm near}-\sqrt{1-g^{-2}}-2.129g\)
decreases.
At \(17/16\), its time is \(>3.0879\)
and radical \(<.338\), which again
gives \(x_1>.64\).
All time/derivative comparisons
follow by positive squaring in
\eqref{app:exact19:eq:Atimes}.

If site five is weak, its first-six
reference block has mean \(b_6<.27\).
Grafting \(.64\) into its first
coordinate, preserving the sum,
raises its energy by
\((6/5)(.64-.27)^2=.16428\).
The other five coordinates remain
\(>(1.5-.64)/5=.172\), above the
next old coordinate \(<.118\).
This is a legal ordered dual
reference, and even the weakest
row of \eqref{app:exact19:eq:Arefs}
then exceeds one.
Every weak end helper is therefore
FAR, in both orientations.

For a far helper on six/six blocks,
block averaging gives
\[
 \Pi U=(5g/6)\Omega_6+\Omega_{14}/12,\quad
 Y=U-\Pi U,\quad w\cdot Y=0,\quad
              \Delta=H-2.288(5g/6+1/12).
\]
The three residual squares are
\[
\begin{array}{c|c}
 \text{own}&Q=\|Y\|^2\\ \hline
 5&(2/3)g^2+1/12\\
 \text{finite }\ge7&4g^4-(16/3)g^2+29/12\\
 \text{dotted}&(5/3)g^2-19/12+g^{-2}.
\end{array}
\tag{A19Q}\label{app:exact19:eq:AhelperQ}
\]
At native five, \(A_{\rm far}>2.188\),
so \(\Delta>.34>.30\), \(Q<.521<13/12\).
For all other finite owns,
\(A_{\rm far}-2.288(5g/6+1/12)\)
increases, since
\(A_{\rm far}'>\sqrt{7.5}>5(2.288)/6\).
The complete radical/time seals are
\[
\begin{array}{c|c|c|c}
 g\text{ interval}&A_{\rm far}\text{ at left}>
      &\sqrt{(1-g^2)(4g^2-3)}\le&\Delta>\\ \hline
 (.9,.91]&2.449&.2318&.3105333\ldots\\
 [.91,.92]&2.478&.2434&.3088666\ldots\\
 [.92,1)&2.5065&.25&.3117 .
\end{array}
\]
The first two radical bounds use
the increasing quartic on those
short intervals; its global maximum
is \(1/16\).
The tight middle lower bound is
\(23165/75000\).
The finite residual polynomial
increases for \(g^2>.81\) and
tends to \(13/12\).
Thus every finite case has
\(\Delta>.30\), \(Q<13/12\).

In the dotted interval the same
derivative argument makes \(\Delta\)
decrease.
At \(17/16\), \(A_{\rm far}>2.9108\),
the radical is \(<.338\), and
\(2.288(5g/6+1/12)=2.2165\).
Thus \(\Delta>.3563>.35\).
Its \(Q\) increases with \(g^2\ge1\)
and at \(17/16\) equals
\(229/768+256/289<148/125\);
the cleared margin is \(271>0\).
These bounds cover the whole
continuous dotted interval.

\subsection{Two weak ends: the exact nonzero cross and both prices}
\label{app:exact19:Atwoweak}

Suppose both five and fifteen
are weak.
Use the six/six reference in
\eqref{app:exact19:eq:Arefs}.
For a left far-helper foot,
its scaled first and last residual
blocks are
\[
\begin{split}
 \sqrt2Y_{\rm own}&=(g/3,g/3,g/3,g/3,g/3-k,-5g/3+k),\\
 \sqrt2Y_{\rm helper}&=(1/6,1/6,1/6,1/6-m,1/6-l+m,-5/6+l).
\end{split}
\]
Here \((k,l,m)\) equals
\((g,2/3,1/3)\) at native five,
\((g,2(1-g^2),0)\) for other finite
owns, and \((1/g,0,0)\) for dotted.
The blocks have zero sum and
direct expansion gives
\eqref{app:exact19:eq:AhelperQ}.
For the right foot of parameter
\(d\), reverse and negate the two
blocks.
In the product of an own block
with a reversed helper block only
its final two deviations survive;
their scaled product is \(-g/3\),
independently of \(k,l,m\).
Consequently for EVERY type pair,
\[
 Y_L\cdot Y_R=(g+d)/6,\qquad
          \|Y_L+Y_R\|^2=Q_L+Q_R+(g+d)/3 .
\]
We have not assumed orthogonality
of the two actual feet.

Choose \(\mu=13/100\) and form
one vector \(W=w+\mu(Y_L+Y_R)\).
At each of the two individual
anchors, the combined debit is
\[
           \mu(5g/6+1/12)
               <403/3200<79/600 .
\]
The remaining coefficient is
nonnegative before using its
original floor.
Both actual pressures and one
completed square give the following
complete three-pair table:
\[
\begin{array}{c|c|c}
 \text{types}&\|Y_L+Y_R\|^2<&\text{gain over }E>\\ \hline
 \text{finite/finite}&17/6&6487/60000\\
 \text{finite/dotted}&17729/6000&.11906\\
 \text{dotted/dotted}&9229/3000&.13000 .
\end{array}
\]
Native five belongs to the finite
column.
For example the first gain is
\(.26(.60)-.0169(17/6)\).
All three gains follow from the
full cross above and the uniform
deficit/norm bounds.
Thus
\[
            \|x\|^2>
               {1341604\over1500000}+{6487\over60000}
                     =1+{3779\over1500000}>1 .
\]
This excludes the two-weak-end case
with every internal hole retained.

\subsection{Zero or one weak end: the actual endpoints and their feet}
\label{app:exact19:Aactualends}

For either of the other reference
rows, if \(x_1\ge.55\) graft \(.55\)
into its constant first block,
preserving the block sum.
For length five the gain exceeds
\((5/4)(.55-.302)^2=.07688\);
for length six it exceeds
\((6/5)(.55-.27)^2=.09408\).
Both exceed \(1-E\) in their
respective rows.
The remaining block coordinates
stay \(>.2375\) or \(>.19\),
above the next coordinate \(<.16\);
a neighbouring hole only improves
that junction.
Reflection gives
\(x_1,-x_{20}<.55\).
The actual endpoint comparison is
\[
               x_1^2\ge(19/20)
                       \frac{\cosh L-1}{\cosh L+1}.
\]
Since \(19/60>.55^2\), both actual
endpoint cosines are less than two.

Deleting one retains \(A_{18}\).
Finite classification gives exactly
isolation; one simple-three at
\(2,3,18,19\); or one simple-four
at two or nineteen.
Two contacts make a nonpositive
cycle; exceptional finite types
cannot contain this rank-eighteen
retained path.
For \(v=g\Omega_1+\beta\Omega_j\)
the complete nonisolated data are
\[
\begin{array}{c|c|l}
 \text{white contact}&\cosh L&\text{small-own obstruction}\\ \hline
 3@2&(38g+18)/20&g=0:\ A^2<0;\quad g=1/2:\ \text{affine triangle}\\
 3@3&(38g+17)/(2\sqrt{20})&g=0:\ \text{affine }E_8\\
 3@18&(38g+2)/(2\sqrt{20})&g=0:\ A^2<0\\
 3@19&(38g+1)/20&g=0:\ A^2<0;\quad g=1/2:\ A^2=0\\
 4@2&(38g+18\sqrt2)/\sqrt{40}&\text{already }>2\text{ at }g=0\\
 4@19&(38g+\sqrt2)/\sqrt{40}&g=0:\ A^2<0 .
\end{array}
\]
For the affine \(E_8\) at old three,
the actual arms have lengths
two (old two/one), one (the cap),
and five (old four through eight).
The positive affine combination has
future time through the cap and
nonpositive omitted-wall products.
Every other positive native contact
is at least \(1/2\); in the first
row the next allowed own is at
least \(1/\sqrt2\).
Substitution makes every row
except \(3@19\) exceed two.

A remaining nonisolated left endpoint
is therefore the far-three cap
\[
 v_L=g\Omega_1+\Omega_{19}/2,\quad
 A_L^2=(19/10)g^2+g/10-21/40,\quad
                    1/\sqrt2\le g<39/38<1.1 .
\]
If both endpoints had this form,
their time squares would be
\(<(8/5)g^2,(8/5)d^2\), since
\(.3g^2+.1g-.525<0\) for \(g<1.1\).
But
\[
 v_L\cdot v_R=(gd+1/4)/10+(19/20)(g+d)
                        >(19/11)gd>(8/5)gd>A_LA_R .
\]
The actual walls are distinct:
coincidence would force \(g=d=1/2\).
Their positive mutual contradicts
acuteness.
At least one endpoint is isolated.

For an isolated endpoint,
\(v=\sigma\Omega_1\),
\(A^2=(19/10)\sigma^2-1>0\).
Own zero/three/four has nonpositive
time, six has its actual affine
\(6,3\) principal, and a parallel
own has the nonzero feasible future
null \(c+n_1\).
The complete remaining own list is
five, finite at least seven, or
dotted.
For dotted own, the actual edge
bound also gives \(19/40<u=\sigma^{-2}<1\).

Normalize each actual isolated foot
by \(\sqrt2\sigma\).
On the left write
\(Z=\zeta-P_{20}/20\), \(\sum\zeta_i=1\);
the right formula is negative reversal.
The complete profiles, pressure
\(x\cdot Z\ge t\) and square \(R=\zeta^2\)
are
\[
\begin{array}{c|c|c|c}
 \text{own}&\zeta&t&R\\ \hline
 5&(1/4,1/4,1/4,1/4)&>.34&1/4\\
 \text{finite }\ge7&(1-u/2,1-u/2,u-1)&
 \{\sqrt{1.9-u}-\sqrt{7u-4-3u^2}\}/\sqrt2&
                      3-4u+3u^2/2\\
 \text{dotted}&(1-u/2,u/2)&
 \{\sqrt{1.9-u}-\sqrt{u(1-u)}\}/\sqrt2&
                      1-u+u^2/2 .
\end{array}
\tag{A19endfeet}\label{app:exact19:eq:Aendfeet}
\]
For native five, use the finite
\(H_4\) correction
\((3/2)\sigma n_1+\sigma n_2+(\sigma/2)n_3\).
Its residual
\(\delta=1-(3/2)\sigma^2>0\)
gives the actual finite projected
point \(p+(A/\delta)(c+k)\),
of square \(-1-A^2/\delta\).
It is future and satisfies every
original wall by the selected
product calculation and positivity.
Its raw pressure is
\(H=AL/(L+\sqrt\delta)\),
\(L^2=A^2+\delta=.4\sigma^2\).
The positive-square bounds
\(A>.4934,L>.5112,\sqrt\delta<.1352\)
give \(H>.390\) and \(t>.34\).

For finite own at least seven,
use
\(c+\sigma^{-1}n_1+2(\sigma^{-1}-\sigma)n_2\).
Its cap/own products vanish, the
next product is
\((3-4\sigma^2)/(2\sigma)<0\),
and its square is
\(-(7-4\sigma^2-3/\sigma^2)<0\).
The dotted core is
\(c+\sigma^{-1}n_1\), square
\(-(1-u)\), zero cap product
and negative own product.
Both have positive future time,
positive coefficients on actual
normals, and every omitted-wall
product nonpositive.
They prove the remaining two
profiles, for this same \(x\).

The following scalar seals cover
all their parameters.
Native seven satisfies
\(.9009<\sigma<.901\), hence
\(1.23<u<1.233\),
\[
 \sqrt{1.9-u}>.8167,\quad
 \sqrt{7u-4-3u^2}<.2671,\quad
                             t>.388,\quad R<.34935 .
\]
The cosine bracket follows by signs
of \(8\sigma^3-4\sigma^2-4\sigma+1\);
it asserts no field restriction.
For finite own at least eight,
\(u<4-2\sqrt2<1.1716\),
\(\sqrt{1.9-u}>.85346\), and
\[
 7u-4-3u^2=1/12-3(u-7/6)^2<1/12 .
\]
Thus \(t>.399\), \(R\le.5\).
At \(u\ge1.1\), \(R\le.415\);
at \(u<1.1\), the radicals improve
to \(>.8944\) and \(<.265\),
giving \(t>.4448\).

For dotted own put \(q=\sqrt{1.9}\),
\(a=q-\sqrt{.9}=1/(q+\sqrt{.9})<.43\).
The lower chord and circle maximum
give
\[
 \sqrt{1.9-u}-\sqrt{u(1-u)}
       \ge q-au-\sqrt{u(1-u)}>.6185,
\]
using \(q>1.378\), \(\sqrt{.9}>.948\),
\(\sqrt{1+a^2}<1.089\).
Hence \(t>.437\).
On \(u\ge.5\) the sharper chord
\(\sqrt{1.9-u}\ge1.4178-.4692u\)
and \(\sqrt{1+.4692^2}<1.105\)
give \(t>.445\).
On \(u<.6\), use \(1.14-.5=.64\)
to get \(t>.4525\).
Also \(R<.6378125\) throughout,
and \(R\le.58\) when \(u\ge.6\).
All square-root and chord comparisons
are strict rational square checks.

On a constant end block of length
\(m\), the projection of \(Z\) is
its boundary dual divided by
\(m\sqrt2\).
Thus for \(V=Z-\Pi Z\),
\[
 w\cdot V=0,\quad \|V\|^2=R-1/m,\quad
                         \Delta_h=t-b_m .
\]
A paid coefficient \(\nu\ge0\)
costs \(\nu/(m\sqrt2)\) at that
individual boundary.
Combining the actual pressure with
the remaining original floors
BEFORE its square gives
\[
                 \|x\|^2\ge E+2\nu\Delta_h-\nu^2(R-1/m).
\tag{A19endpay}\label{app:exact19:eq:Aendpay}
\]

\subsection{No weak end: one isolated endpoint already suffices except at five}
\label{app:exact19:Azeroweak}

For zero weak ends, take \(\nu=1/2\)
in the isolated finite case.
Its price is \(<1/14<83/625\).
At native seven the gain is
\(>.086-.15/4=.0485\).
For all finite labels at least eight,
the complete subdivision is
\[
\begin{array}{c|c|c|c}
 u&t>&R-1/5<&\Delta_h-(R-1/5)/4>\\ \hline
 (1,1.1)&.4448&.3&.0678\\
 [1.1,1.14)&.4122&.215&.05645\\
 [1.14,1.1716)&.399&.1894&.04965 .
\end{array}
\]
The middle interval uses radicals
\(>.8717,<.28868\), with
\(.4122(1.4143)<.8717-.28868\).
The last uses \(>.85346,<.28868\);
the polynomial \(R\) decreases here,
with \(R(1.1)=.415,R(1.14)=.3894\).
Thus every finite isolated own
at least seven gives
\[
                         \|x\|^2>.9522798+.0485>1 .
\]

For an isolated dotted own, take
\(\nu=1/3\), price
\(<1/21<83/625\).
At \(u<.6\) the gain is
\(>(2/3)(.1505)-.438/9=31/600\).
At \(u\ge.6\) it is
\(>(2/3)(.143)-.38/9=239/4500\).
Both gains exceed \(1-.9522798\).

Only isolated native five remains.
It cannot coexist with a nonisolated
far-three opposite: its time is
\(<.497\), while the latter has
time \(<1.257\).
Their spatial product is
\(\sigma(d/10+19/20)>.95(.809)=.76855\),
whereas the time product is
\(<.497(1.257)=.624729\).
Their actual mutual would be
positive.
Thus the other endpoint is isolated
as well, and the preceding cases
force it too to have native five.

For each normalized native-five foot,
\(\Delta_h>.038\), \(\|V\|^2=.05\).
Use \(\nu=3/4\) at BOTH ends;
each price is \(<3/28<83/625\).
The two residuals have disjoint
first/last five-block support and
zero block sum, so their cross
is exactly zero.
The single vector
\(w+(3/4)(V_L+V_R)\) gives
\[
 \|x\|^2>.9522798+
       2\{(3/2)(.038)-(9/16)(.05)\}
                       =1.0100298>1 .
\]
This exhausts zero weak ends.

\subsection{One weak end: the helper and endpoint must be paid jointly}
\label{app:exact19:Aoneweak}

Orient the weak site to five.
It is far by the helper argument.
On the six/five reference its
block projection and residual are
\[
 \Pi U=(5g/6)\Omega_6+\Omega_{15}/10,\quad
 Y=U-\Pi U,\quad w\cdot Y=0,
 \qquad\Delta=H-2.288(5g/6)-2.128/10 .
\]
Relative to six/six,
the square decreases by \(1/60\)
and the deficit decreases by
\(83/3750\).
The exact finite middle-row bound
\(23165/75000\), and the stronger
native-five/dotted bounds, give
\[
                       \Delta>.2867,\qquad Q=\|Y\|^2<1.168 .
\]
The dotted upper square is
\(148/125-1/60=1751/1500<1.168\).

We must retain the full cross with
the isolated endpoint.
Write the far-helper correction
at old nineteen/eighteen as
\((l,m)=(2/3,1/3)\),
\((2(1-g^2),0)\), or \((0,0)\).
Its scaled last-five residual is
\[
             \sqrt2Y_{\rm helper}
               =(1/5,1/5,1/5-m,1/5-l+m,-4/5+l).
\]
For an isolated STRONG-side endpoint,
its residual there is \(1/5\)
minus the reversed \(\zeta\).
Direct multiplication gives
\[
\begin{array}{c|c}
 \text{isolated own}&Y\cdot V\\ \hline
 5&1/(20\sqrt2)\\
 \text{finite }\ge7&
       \{.8-u/2+m(1.5u-2)\}/\sqrt2\le(.8-u/2)/\sqrt2\\
 \text{dotted}&
       \{.8-u/2-l(1-u)-mu/2\}/\sqrt2\le(.8-u/2)/\sqrt2 .
\end{array}
\]
The finite sign uses \(u<1.233<4/3\).
For an isolated WEAK-side endpoint,
its \(\zeta\) ends before coordinate
five; multiplication with the own
six-block row printed earlier gives
\[
                             Y\cdot V=g/(3\sqrt2)<85/336 .
\]
These formulas retain every correction
coefficient, including its sign.

For paid nonnegative \(\mu,\nu\),
one vector \(W=w+\mu Y+\nu V\)
therefore yields
\[
 \|x\|^2\ge E+2\mu\Delta+2\nu\Delta_h
           -\mu^2Q-2\mu\nu(Y\cdot V)-\nu^2(R-1/m).
\tag{A19joint}\label{app:exact19:eq:Ajoint}
\]
Both individual anchor debits must
fit before either old floor is used.

If an isolated endpoint is on the
strong side, choose \(\mu=.148\).
Its own-six debit is
\[
                       \mu(5g/6)<3145/24000<79/600 .
\]
At fifteen, the helper AND endpoint
together cost
\[
       \mu/10+\nu/(5\sqrt2)
              <37/2500+4/35<83/625
\]
for all choices \(\nu\le.8\) below.
The common helper gain exceeds
\(.059279328\).
Using the full cross, the complete
table is
\[
\begin{array}{c|c|c|c|c|c}
 \text{isolated strong-side own}&\nu&
 \Delta_h>&R-1/5<&Y\cdot V<&\text{total gain over }E>\\ \hline
 5&.8&.038&.05&1/28&.07962218\\
 7&.3&.086&.15&.133&.085568928\\
 \text{finite }\ge8,\ u\ge1.1&.3&.097&.215&.179&.082234128\\
 \text{finite }\ge8,\ u<1.1&.3&.1428&.3&.215&.098867328\\
 \text{dotted},\ u\ge.6&.25&.143&.38&.358&.080537328\\
 \text{dotted},\ u<.6&.25&.1505&.438&.402&.077406328 .
\end{array}
\]
For instance the last row is exactly
\(.059279328+.5(.1505)-.0625(.438)-.074(.402)\).
All cross bounds use \(\sqrt2>7/5\)
and the stated intervals.
Every gain is \(>.0774\).

If the strong endpoint is not
isolated, the weak endpoint is
isolated by the actual endpoint
classification.
It cannot have native five:
the same time/spatial comparison
with a far-three opposite, just
proved, would force that opposite
to be isolated.
The remaining finite/dotted
weak-side cases use
\(\mu=.1,\nu=.35\).
Now BOTH feet spend the SAME
six coefficient; the full debit is
\[
 \mu(5g/6)+\nu/(6\sqrt2)
                   <17/192+1/24=25/192<79/600 .
\]
At fifteen the cost is only \(.01\).
The common helper gain is \(>.04566\)
and the full cross penalty is
\(<.07(85/336)=119/6720\).
The complete table is
\[
\begin{array}{c|c|c|c}
 \text{isolated weak-side own}&\Delta_h>&R-1/6<&
                                    \text{total gain over }E>\\ \hline
 7&.118&.183&.08813416\\
 \text{finite }\ge8&.129&1/3&.07741833\\
 \text{dotted}&.167&.472&.08703166 .
\end{array}
\]
The tight finite row is
\(.04566+.7(.129)-.1225/3-119/6720
=.077418333\ldots\).
All gains again exceed \(.0774\).
Thus in every one-weak-end case,
\[
              \|x\|^2>{3462479\over3750000}+.0774
                       =1+{2729\over3750000}>1 .
\]

\begin{proposition}\label{app:exact19:prop:wholeA19}
A whole \(A_{19}\) forest cannot occur
at the original farthest vertex of
a compact Coxeter nineteen-polytope.
\end{proposition}
\begin{proof}
The number of weak sites among the
two central end sites five/fifteen
is exactly zero, one or two.
The three preceding arguments
exclude these cases, with equality
assigned consistently to weak.
Their shared local window, endpoint
and foot lemmas precede all three
consumers and do not use any
end-case conclusion.
Every remaining internal weak
site was retained through its
actual omission, and every new
pressure was paid inside one
original-vector reference.
\end{proof}


\subsection{Whole \(B_{19}\): the real short layers before any terminal conclusion}
\label{app:exact19:Bwindow}

Use actual unit roots
\[
 n_i=(e_i-e_{i+1})/\sqrt2\ (i\le18),\qquad n_{19}=e_{19},
 \quad \omega_i=\sqrt2P_i\ (i\le18),\quad\omega_{19}=P_{19}.
\]
Now \(x_1>\cdots>x_{19}>0\).
Put \(S_i=x\cdot P_i\),
\(f_i=\sqrt i-.58/\sqrt2\),
\(d_i=\sqrt{i+1}-\sqrt i\), and
\[
                         E_K=f_5^2/5+\sum_{i=5}^{K-1}d_i^2 .
\]
Weak includes equality \(S_i\le f_i\).
The local whole-diagram lemma
\ref{app:tail:lem:whole58local} applies
literally through fifteen:
both retained old components have
rank at least four.
Its complete weak list is
\[
 v=g\omega_i+\omega_1/2,\quad
 A^2=2ig^2+2g-1/2,\quad
 \cos(\pi/5)\le g<6/5,
\tag{B19packet}\label{app:exact19:eq:Bpacket}
\]
with own five, finite at least seven,
or dotted, together with isolated
dotted five.
Pure helpers are nonadjacent.
We must extend this at the actual
\(B_3,B_2\) boundary.

At sixteen the right \(B_3\) has
a long-three extension or the
short-three \(F_4\) extension.
The latter is foreign-sealed;
own zero has the actual affine
\(F_4\) principal, whereas every
positive own has ratio
\(\ge1+1/\sqrt2>83/50\).
A full bridge is \(B_{19}\) and
directly strong.
Thus the same list holds through
sixteen, also at the old \(.61\)
threshold.
In particular some new strong
\(K=15\) or sixteen exists.
The full conditional early-reference
Lemma~\ref{app:tail:lem:whole58reference}
therefore gives a paid vector
supported through \(K\), of sum
\(f_K\), remaining individual
\(P_K\) coefficient \(>1/9\), and
effective energy
\[
               L_{15}>.94006,\qquad L_{16}>.95618,
               \quad f_{15}<3.463,\ f_{16}<3.590 .
\tag{B19early}\label{app:exact19:eq:Bearly}
\]
Effective energy means \(2D_W-\|W\|^2\),
where \(D_W\) is its actual pressure
demand; it need not equal \(\|W\|^2\).

At seventeen the right block is
\(B_2\), with simple-three at either
end.
The long-three row is strong or
affine.
For short three, own zero/four/six
have actual affine principals,
parallel has its positive future
null, and finite own at least seven
has directly strong ratio
\(\sqrt2g+1/2>83/50\).
A full bridge has residual \(1/34\)
and numerator
\(\ge g+1/34+1/(2\sqrt2)\), strong
even at own zero.
Only own three and five remain.

For own five, let \(g=\varphi/2\),
\(\alpha=\sqrt2g+1/2\).
The LEFT actual \(H_4\), using
seventeen/sixteen/fifteen, has
\[
 \delta=(5-3\varphi)/8,\quad A^2=17\alpha^2-1/2,
 \quad V_1=\cdots=V_{14}=1,\quad V^2<17,
                    \quad x\cdot V>G>4 .
\]
The point is the positive finite
projection, not its spacelike
correction; every omitted-wall
product is nonpositive.
Use \(\alpha>1.643\),
\(\sqrt\delta/\alpha<.083\)
for the last bound.
An OLD strong eleven or twelve
exists by the original bulk.
Their positive references, in both
old initial-five branches, have
\[
 q_{11}>.8472225,\ q_{12}>.8689525,\quad
 f_{11}^{\rm old}<2.886,\ f_{12}^{\rm old}<3.034,
                    \quad\text{individual endpoint price}>1/7 .
\]
These are the actual prefix
construction with its holes retained,
not an assertion that old five is
strong.
Their support lies in the flat
first fourteen.
One paid square gives
\[
 .8472225+2(4-2.886)/7-6/49>1.043,\qquad
 .8689525+2(4-3.034)/7-5/49>1.042 .
\]

For own three, the SAME cap has a
finite \(D_{19}\) replacement after
releasing the short old nineteen.
Positive inverse projection gives
an actual reciprocal short-edge
endpoint, with
\[
 \cosh L_{19}=17/2+19\sqrt2/4>15,\qquad
               S_{19}^2>133/8>(4.07)^2 .
\]
Use \eqref{app:exact19:eq:Bearly},
and the pressure \(Z=P_{19}\).
The legal \(1/9\) endpoint payment
gives
\[
 .94006+2(4.07-3.463)/9-4/81>1.025,\qquad
 .95618+2(4.07-3.590)/9-3/81>1.025 .
\]
Thus coverage reaches seventeen.

Nonadjacency is also actual at this
boundary: replacing two adjacent
helper caps at sixteen/seventeen
retains the long \(A_{15}\) and
the independent old \(B_2\).
Adding their actual far-one leaves
gives the positive \(D_{17}+B_2\).
The later seventeen/eighteen pair
instead gives \(D_{18}+A_1\).
The positive two-cap inverse and
all-wall projection therefore apply.
No long nineteenth old normal is
invented.
Every weak site five through
seventeen is now
\eqref{app:exact19:eq:Bpacket},
except isolated dotted five.

\subsection{An isolated long-end foot is a local geometric kernel}
\label{app:exact19:longfeet}

We record the point formulas once.
They apply to any actual isolated
long-end cap with initial long roots
one/two/three as above, in \(B\) or
\(D\), whenever the listed own occurs.
Write
\[
                  c_L=Ap-\gamma\omega_1,\qquad A^2=2\gamma^2-1 .
\]
For each remaining own type the
actual normalized pressure is
\(x\cdot Z\ge h\), \(\sum Z_j=1\):
\[
\begin{array}{c|c|c|c}
 \text{own}&Z&h&R=Z^2\\ \hline
 5&P_4/4&\sqrt{(\varphi-1)/4}&1/4\\
 \text{finite }\ge7&
 (1-u/2)(e_1+e_2)+(u-1)e_3&
 \{\sqrt{2-u}-\sqrt{7u-4-3u^2}\}/\sqrt2&
                   3-4u+3u^2/2\\
 \text{dotted}&(1-u/2)e_1+(u/2)e_2&
 \{\sqrt{2-u}-\sqrt{u(1-u)}\}/\sqrt2&
                   1-u+u^2/2 .
\end{array}
\tag{B19headfeet}\label{app:exact19:eq:longfeet}
\]
Here \(u=\gamma^{-2}\), in
\((1,\cos(\pi/7)^{-2}]\) for finite
own, and \((0,1)\) for dotted.
All coordinates of \(Z\) are
nonnegative.

At native five, the actual \(H_4\)
correction is
\(\gamma(3n_1/2+n_2+n_3/2)\),
with residual
\(\delta=1-3\gamma^2/2>0\).
The future finite projected point
\(p+(A/\delta)(c_L+k)\) has square
\(-1-A^2/\delta\) and every original
wall product nonpositive.
Its displacement is
\((\gamma/4)\omega_4\), and its
pressure is \(H=AB/(B+\sqrt\delta)\),
\(B=\gamma/\sqrt2\).
The real native-five identity
\(4\gamma^2=2\gamma+1\) gives
\(\sqrt\delta=(1-\gamma)/\sqrt2\),
hence \(H=\gamma A\) and the first
row exactly.
In particular \(h>.393075\),
by positive rational squaring.

For finite own at least seven,
the actual negative core is
\[
 c_L+\gamma^{-1}n_1+2(\gamma^{-1}-\gamma)n_2 .
\]
Its cap/own products vanish, the
next product is
\((3-4\gamma^2)/(2\gamma)<0\),
and its square is
\(-(7-4\gamma^2-3/\gamma^2)<0\).
For dotted own use
\(c_L+\gamma^{-1}n_1\), square
\(-(1-u)\), zero cap product and
negative own product.
Both have positive future time
and positive coefficients on actual
normals.
Every omitted original-wall product
is nonpositive termwise.
Their spatial displacements, divided
by \(\sqrt2\gamma\), give the other
two rows on the SAME original \(x\).
This proves the kernel without an
\(I\) spectator or an arithmetic
hypothesis.

For every finite own at least seven
and every dotted own,
\[
                              .8h-.16R>.2893 .
\tag{longF}\label{app:exact19:eq:longF}
\]
At native seven,
\(.9009<\gamma<.901\) gives
\(A>.7894\),
\(\sqrt{7-4\gamma^2-3/\gamma^2}<.2395\),
\(\sqrt2\gamma<1.2743\);
thus \(h>.4315\).
Also \(1.2318<u<1.233\) and
\(R<.349\), giving \(>.28936\).

For finite labels at least eight,
\(1<u<1.172\).
Put \(a=.8/\sqrt2\),
\(D=(u-1)(4-3u)\).
The derivative of the left side is
\[
 a\left\{-\frac1{2\sqrt{2-u}}-\frac{7-6u}{2\sqrt D}\right\}
                         +.64-.48u .
\]
On \((1,7/6]\), the first term is
\(<-.28\), the second is nonpositive,
and the last is at most \(.16\).
On \([7/6,1.172)\), \(\sqrt D>.284\),
the positive second term is
\(<.032\), and the other terms
are \(<-.28+.08\).
Thus the scalar decreases with
\(u\); its minimum is at native
eight.
There the radical bounds are
\(>.910,<.289\), and \(R<.373\).
Hence \(h>.4388\) and the scalar
is \(>.29136\).

For dotted own,
\(.5656<a<.5657\).
On \(u\le.5\), use
\(\sqrt{2-u}\ge1.4142-.379u\),
\(R\le1-.75u\), giving
\[
 .8h-.16R>.63987152-.0944u-.5657\sqrt{u(1-u)}
                                      >.30587152 .
\]
The circle bound uses
\(\sqrt{.0944^2+.5657^2}<.5736\).
On \(u\ge.5\), use
\(\sqrt{2-u}\ge1.4494-.4494u\),
\(R\le.75-.25u\), giving
\[
 .8h-.16R>.69978064-.2142u-.5657\sqrt{u(1-u)}
                                      >.29023064 ,
\]
with \(\sqrt{.2142^2+.5657^2}<.6049\).
The chords lie below the concave
root by their endpoint square
checks.
This proves the entire continuous
dotted range.

\subsection{The \(B_{19}\) long head and its paid gains}
\label{app:exact19:Bhead}

The actual long head is isolated.
Indeed deleting one retains
\(B_{18}\), whose complete finite
one-node list is isolation or
simple-three at its long leaf two.
A second contact makes a cycle;
another four or an additional fork
is not finite; exceptional finite
ranks are too small.
The nonisolated cap would have
\[
 v=\gamma\omega_1+\omega_2/2,\quad
 A^2=2\gamma^2+2\gamma,\quad\cosh L_1=2\gamma+1 .
\]
Own zero has zero time and own
three is the actual affine
triangle, so \(\gamma\ge1/\sqrt2\)
and \(x_1>\sqrt{\sqrt2-1}>.643\).

The OLD actual bulk through fifteen
has some strong \(K=14\) or fifteen.
If old five is strong, its paid
reference has energy \(>.90715\),
with every actual weak hole omitted.
Graft
\(.643\) into its first-five block
of mean \(<.361\).
The other four coordinates remain
\(>.2904\), above the outgoing
slope \(<.214\), and the gain is
\((5/4)(.643-.361)^2=.099405\).
If old five is weak, six is strong.
Begin only with the FORMAL
strong-five curve, whose scalar
energy is \(E_K^{\rm old}\).
Here \(f_i^{\rm old}=\sqrt i-.61/\sqrt2\)
and \(E_K^{\rm old}=(f_5^{\rm old})^2/5
+\sum_{i=5}^{K-1}d_i^2\).
Before assigning any pressure to
it, replace its ENTIRE first-six
block by \(.643\) and five equal
slopes of sum \(f_6^{\rm old}-.643\).
Their energy exceeds
\[
 .643^2+(2.01815-.643)^2/5=.7916565045 ,
\]
whereas the formal first-six
energy is \(E_6^{\rm old}<.6971\).
The new slopes are ordered:
\(.643>(f_6^{\rm old}-.643)/5>.275\),
above every outgoing slope
\(\le d_6<.197\).
There is NO coefficient at five.
All later actual weak knots lie
in \(7,\ldots,K-1\); at most four
are nonadjacent, each costing
\(<1/3200\).
Their affected open intervals
are disjoint from the first-six
replacement, even for a hole
at seven.
Only now combine the actual
pressures at one, six and the
remaining strong knots.
One square gives
\[
 \|x\|^2>
 E_K^{\rm old}-E_6^{\rm old}
      +.7916565045-4/3200
 >.9084-.6971+.7916565045-.00125
       =1.0017065045>1 .
\]
Both references are ordered and
paid by the original \(x\).
Thus the nonisolated head is
impossible.

For the isolated cap,
\(A^2=2\gamma^2-1>0\) excludes
own zero/three/four.
Own six has the affine \(6,3\)
principal on the actual first two
roots; parallel has its feasible
positive future null.
The remaining owns are exactly
native five, finite at least seven,
or dotted.
Their actual feet are
\eqref{app:exact19:eq:longfeet}.

A native-five head forces true
new strong five/six/seven.
For a helper at \(i\le7\),
\[
 c_L\cdot c_i=\gamma_0(2g+1)-A_0A_i,\qquad
       \gamma_0=\varphi/2,\quad A_0^2=(\varphi-1)/2 .
\]
Since \(\gamma_0^2/A_0^2=\varphi^3/2>21/10\),
\[
 A_i^2-\frac{\gamma_0^2}{A_0^2}(2g+1)^2
                   <(28g^2-32g-13)/5<0
                  \qquad(.8<g<1.2).
\]
The convex quadratic is negative
at both endpoints.
The actual mutual would be
positive.
For the remaining weak isolated
dotted five,
\[
 c_L\cdot c_5=\varphi g-A_0\sqrt{10g^2-1}.
\]
Its sign is strictly negative,
since the squared difference in
the opposite order is
\((4\varphi-6)g^2-(\varphi-1)/2
\ge(7\varphi-11)/2>0\).
Weakness gives \(g<1/.58<2\);
\(\varphi-\sqrt{5(\varphi-1)}>-3/20\)
then bounds the mutual in
\((-3/10,0)\), a forbidden native
interval.
This proves the stated three
strong sites without a later
terminal floor.

For any strong-five \(.58\)
reference with first-five mean
\(b=f_5/5<.3652\) and individual
\(P_5\) coefficient \(>.151\),
a head correction
\(\mu(Z-P_5/5)\) costs \(\mu/5\)
before its floor is used.
It is orthogonal to that constant
block and has square \(R-1/5\).
Its effective gain is
\[
                             G_L=2\mu(h-b)-\mu^2(R-1/5).
\]
The complete paid gains are
\[
\begin{array}{c|c|c}
 \text{head}&\mu&G_L>\\ \hline
 5&.56&.015456\\
 \text{finite }\ge7&.4&.02914\\
 \text{dotted}&\sqrt2/5&.03384 .
\end{array}
\tag{B19headgains}\label{app:exact19:eq:Bheadgains}
\]
The first uses \(h>.393\);
the second is
\eqref{app:exact19:eq:longF}
minus \(.8b\) plus \(.032\).
For the dotted row, expand its
gain as
\[
 .4\{\sqrt{2-u}-\sqrt{u(1-u)}-\sqrt2b\}
                                -.04(1.6-2u+u^2),
\]
where \(\sqrt2b<.5166\).
For \(u\le.5\), its lower bound is
\(.29504-.0916u-.4\sqrt{u(1-u)}>.04404\),
using \(\sqrt{.0916^2+.4^2}<.4104\).
For \(u\ge.5\), it is
\(.32912-.15976u-.4\sqrt{u(1-u)}>.03384\),
using \(\sqrt{.15976^2+.4^2}<.4308\).
The same two concave root chords
prove these inequalities.
All three individual prices fit.

\subsection{Curved early references with the whole initial-five debt}
\label{app:exact19:Bcurvedearly}

Some subsequent pressures need a
stronger early reference than
\eqref{app:exact19:eq:Bearly}.
We give it by the actual initial
five alternatives, retaining all
holes and full cross terms.
Each resulting vector is supported
through \(K=15\) or sixteen,
has sum \(f_K\), and keeps endpoint
coefficient \(>1/9\).

First record a useful exact scalar:
\[
                              E_{17}>.9726645 .
\tag{B19E17}\label{app:exact19:eq:BE17}
\]
Its first block is
\(f_5^2/5=1.03364-.116\sqrt{10}>.6668157\),
using \(\sqrt{10}<3.162278\).
For \(a=4i+2\),
\(u=d_i^2\) satisfies \(u(a-u)=1\).
Three positive reciprocal
substitutions give
\(u>(a^2-1)/(a(a^2-2))\).
At denominator \(10^7\), the full
list of lower numerators and
cleared positive margins is
\[
\begin{array}{c|r|r@{\qquad}c|r|r}
 i&L_i&\text{margin}&i&L_i&\text{margin}\\ \hline
5&455488&5248&11&217494&13464\\
6&385186&536&12&200080&8000\\
7&333704&14240&13&185248&115712\\
8&294372&20208&14&172465&14860\\
9&263340&21360&15&161332&72272\\
10&238230&27080&16&151549&273164 .
\end{array}
\]
The margin is
\(10^7(a^2-1)-L_i a(a^2-2)\).
The numerators sum to \(3058488\),
proving \eqref{app:exact19:eq:BE17}.
This is a finite rational hand
certificate, not a script-output
premise.
Also \(d_{17}^2>1/70\), giving
\(E_{18}>.98659\).

If five is strong, its curved
reference loses \(<1/1024\) over
all nonadjacent internal holes.
Adding the actual head pressure
inside that reference gives
\[
                   L_{15}>.955516,\qquad L_{16}>.971636 .
\]
These conservative values use the
smallest gain in
\eqref{app:exact19:eq:Bheadgains};
the appropriate finite/dotted head
is stronger.

If five is a weak helper, the
actual packet gives \(x_1>.5455\)
and six is strong.
Use the first-six reference through
these actual heights, followed by
the ENTIRE curved retained window.
Its first-six energy is \(>.74385\),
rather than the ordinary \(<.7126\).
Since \(d_{15}^2<.016134\),
\(d_{16}^2<.015155\),
\eqref{app:exact19:eq:BE17} gives
\(E_{15}>.9413755,E_{16}>.9575095\).
After all holes,
\[
                          L_{15}>.9716,\qquad L_{16}>.9877 .
\]
No long-head energy was added to
this initial-coordinate gain.

If five is isolated dotted,
the native-five head has already
been excluded by its compatibility.
For the other head types keep its
actual foot \(x\cdot Z\ge h\)
and the initial-five foot
\[
 Z_5=P_5+(v/2)(e_6-e_5),\quad
                x\cdot Z_5>T_5>1.8185,\quad .4<v<1 .
\]
This is the actual dotted-five
negative core; the weak condition
forces \(v>.4\) by the full
continuous loss bound in the
shared early-reference lemma.
Choose a genuinely strong
\(J=6\) or seven.
The reference is constant
\(m=f_J/J\) on the first \(J\)
coordinates and follows every
retained curved knot through \(K\).
The complete conservative data are
\[
\begin{array}{c|c|c|c|c|c}
 J&m<&w_{15}^2>&w_{16}^2>&
 \text{individual }P_J\text{ coefficient}>&Q_5<\\ \hline
 6&.3399&.9217&.9378&.1435&77/150\\
 7&.31938&.9041&.9202&.1366&194/175 .
\end{array}
\]
For the energies, \(f_6^2/6>.69317\),
\(f_7^2/7>.71400\).
The bound \(d_i^2>1/(4i+2)\) gives
at denominator \(10^6\) the lower
numerators
\[
38461,33333,29411,26315,23809,21739,
20000,18518,17241,16129
\]
for \(i=6,\ldots,15\).
Those from six through fourteen sum to \(228827\),
and from seven to fourteen to
\(190366\); \(d_{15}^2>.016129\).
The maximal omission losses are
\(<.000267\) after six and
\(<.000181\) after seven, by the
decreasing sums \(1/[32(j-1)^3]\)
over their earliest permitted
nonadjacent sets.
These prove every energy entry.
The \(P_J\) price is \(m-d_J\);
the end-\(K\) price stays \(>1/9\).
Finally
\[
 Q_5=\|Z_5-(5/J)P_J\|^2
          =5-25/J-v+v^2/2<4.68-25/J .
\]

Let \(Y=Z-P_J/J\),
\(Y_5=Z_5-(5/J)P_J\).
Their first-\(J\) sums vanish,
so they are orthogonal to \(w\).
Because \(Z\) is supported on the
first three coordinates and
\(Z_5=1\) there,
\[
             Y^2=R-1/J,\qquad Y\cdot Y_5=1-5/J .
\]
The one vector \(w+.4Y+.09Y_5\)
spends the COMBINED
\(P_J\) price \(.85/J\), which
fits both rows.
Its gain is
\[
 .8h-.16R-.8m+.16/J
 +.18(T_5-5m)-.0081Q_5-.072(1-5/J).
\]
The nonzero last cross is retained.
Using \eqref{app:exact19:eq:longF},
the gain is \(>.04930\) at six
and \(>.06699\) at seven.
Thus both initial branches give
\[
                         L_{15}>.9710,\qquad L_{16}>.9871 .
\]
These are effective energies of
the same paid vector, not new
independent norm summands.

The identical curved construction
at a genuinely strong seventeen
has initial reference squares
\(>.9530,.9354\) in the last case,
so its gains give totals
\(>1.00230,1.00239\).
For helper five it gives
\(>.97231+.03125-1/1024>1.0025\).
At a later actual strong eighteen,
the extra endpoint increment
dominates the possible additional
hole, so the conclusions remain
valid.
These are conditional uses of
an ACTUAL strong endpoint;
none is used to assume that
endpoint beforehand.


\subsection{The zero-own short-only terminal walls are impossible}
\label{app:exact19:Bzeroown}

Suppose an actual last-long cap is
\[
 c_T=Bp-\delta P_{19},\quad B^2=19\delta^2-1,\qquad
                 \delta=\cos(\pi/m),\quad m=5,6,7 .
\]
We do not yet assume that this
exhausts the terminal possibilities.
If seventeen were weak, its actual
helper would have
\(A_H^2=34g^2+2g-1/2\).
The two walls are distinct by their
old-one contacts.
Put
\[
 M=A_HB-\delta(34g+1)/\sqrt2=-c_H\cdot c_T .
\]
For \(b=\delta^2>.65,g>.8\), the
difference of the two positive
squared terms is
\[
                 (2b-1)(34g^2+2g)+1/2-10b .
\]
It increases in both variables
and at \((.65,.8)\) equals \(1.008\).
Thus \(M>0\).
Positive squaring of its derivatives
also proves monotonicity in both:
the respective positive differences
are
\[
 (68b-34)A_H^2+18(19b-1),\qquad
                 (38b+17)A_H^2-171b+9 .
\]
For the latter use \(A_H^2>22.86\),
\(.65<b<.812\).
At \(g=6/5,\delta=.901\), rational
square bounds give
\(A_H<7.132,B<3.798\), and
\(\delta(34g+1)/\sqrt2>26.6\).
Consequently
\[
                     0<M<7.132(3.798)-26.6
                             =.487336<1/2 ,
\]
a forbidden actual-wall product.
Seventeen is therefore genuinely
strong before any head split.

If the isolated long head were
finite at least seven or dotted,
the three initial-five constructions
of Section~\ref{app:exact19:Bcurvedearly}
at that actual strong seventeen
would contradict the norm.
For strong five, the reference
square is \(>.97133\) and its
respective head gain is
\(>.02914\) or \(>.03384\).
The helper-five and dotted-five
branches already exceed one by
their full curved/joint references.
Thus only a native-five long
head remains.

Write it as
\[
 c_L=Ap-qe_1,\quad
 A^2=(\varphi-1)/2,\quad q=\varphi/\sqrt2,\quad
          x\cdot(P_4/4)\ge h=\sqrt{(\varphi-1)/4}>.393075 .
\]
The true strong five/six/seven
compatibility was proved before
this terminal assumption.
Use the ordinary reference through
strong seventeen, omitting all
actual internal weak helpers.
The largest possible curvature
loss is bounded by holes
\(8,10,12,14,16\):
\[
 \sum_{j=8,10,12,14,16}\frac1{32(j-1)^3}
 <.000091108+.000042867+.000023479+.000014224+.000009260
 <.000181 .
\]
The decreasing weights and
nonadjacency make this earliest
permitted independent set maximal.
Hence the actual reference has
\[
 w^2>.9724835,\quad b=f_5/5<.365190,\quad
 f_{17}<3.712985,\qquad
                  \text{prices }P_5>.151,\ P_{17}>.12 .
\tag{B19tightref}\label{app:exact19:eq:Btightref}
\]
No unproved strong eighteen
was used.

At terminal short five, the actual
mutual magnitude is exactly
\[
 -c_L\cdot c_T=
       \frac{\sqrt{4+15\varphi}-(\varphi+1)}{2\sqrt2}
                                 \in(.953,.956).
\]
Use \(1.618<\varphi<1.619\),
\(1.414<\sqrt2<1.415\) and positive
squaring for this interval.
But \(\cos(\pi/10)<.952\), while
\(\cos(\pi/11)>47/49>.956\):
the latter follows from
\(\pi<22/7\) and
\(\cos t\ge1-t^2/2\).
There is no finite native magnitude
in this gap, and the magnitude is
neither zero nor at least one.
The two actual walls cannot have
this product.

At short six put
\(\delta=\sqrt3/2\), \(B=\sqrt{53}/2\)
and \(k=AB-q\delta\).
The complete rational seals
\[
 .5558<A<.556,\quad1.1440<q<1.1442,\quad
 .8660<\delta<.8661,\quad3.6400<B<3.6401,\quad
                                      1.032<k<1.034
\]
follow from positive squaring.
The ACTUAL positive combination
\(\xi=c_L+c_T\) is future timelike,
square \(-2(k-1)\), with both
selected cap products \(1-k<0\).
Every old and unknown wall product
is nonpositive termwise.
It gives the original pressure
\[
 U=qe_1+\delta P_{19},\qquad
              x\cdot U\ge H=A+B-\sqrt{2(k-1)}>3.9348 .
\]
The reference pairing is
\(w\cdot U=qb+\delta f_{17}<3.634\);
its deficit exceeds \(.3\).
After paying BOTH anchors, the
actual residual is
\[
 V=U-\delta P_{17}-(q/5)P_5
          =q(e_1-P_5/5)+\delta(e_{18}+e_{19}),\qquad
       w\cdot V=0,\quad V^2=(4/5)q^2+2\delta^2<2.548 .
\]
With \(\nu=.1\), the individual
prices are \(<.023\) at five and
\(<.087\) at seventeen.
A single square gives
\[
                      \|x\|^2>.9724835+.06-.02548>1.007 .
\]
No separate head-energy estimate
was added here.

At short seven we need the full
head/core cross.
Let \(t=27/25\), \(k=AB-q\delta\).
The exact polynomial
\(8\delta^3-4\delta^2-4\delta+1=0\)
gives the rational seals
\[
\begin{gathered}
 .9009688<\delta<.9009690,\quad
 .5558929<A<.5558931,\quad
 1.1441227<q<1.1441231,\\
 3.79778<B<3.79779,\qquad1.08033<k<1.08035 .
\end{gathered}
\]
For the first bracket the cubic
increases on \((.9,.901)\);
at \(.901-e\) it equals
\(.000257608-8.275224e+17.624e^2-8e^3\),
whose signs at \(e=.0000312,.0000310\)
are negative/positive.
The other seals are positive
square/product comparisons.

Now \(\xi=c_L+t c_T\) is future,
with selected cap products
\(1-tk<0,t-k<0\), and square
\(-s\), \(s=2tk-1-t^2>0\).
All other original-wall products
are nonpositive.
The same original vector satisfies
\[
 U=qe_1+t\delta P_{19},\qquad
 x\cdot U\ge H=A+tB-\sqrt s>4.2486453 ,
\]
because \(\sqrt s<.40885\).
Also \(w\cdot U=qb+t\delta f_{17}<4.03074\),
so the actual deficit is \(>.2179\).

Use the independent actual native-five
\(H_4\) pressure through
\(Y=P_4/4-P_5/5\).
Its square is \(1/20\), its cross
with \(w\) is zero, and its deficit
is \(h-b>.027885\).
For the two-cap point pay the
residual
\[
 V=U-t\delta P_{17}-(q/5)P_5
        =q(e_1-P_5/5)+t\delta(e_{18}+e_{19}).
\]
Its full products are
\[
 w\cdot V=0,\quad
 V^2=(4/5)q^2+2t^2\delta^2<2.941,\qquad
                         Y\cdot V=q/20<.057207 .
\]
Choose \(\mu=.48,\nu=.065\)
and use ONE vector \(w+\mu Y+\nu V\).
Its complete individual debits are
\[
        P_5:\ (\mu+\nu q)/5<.111<.151,\qquad
        P_{17}:\ \nu t\delta<.064<.12 .
\]
The actual primitive pressures,
with these coefficients removed
before their floors are used,
give the single completed-square
gain
\[
\begin{split}
 &.96(.027885)-.2304(.05)
  +.13(.2179)-.004225(2.941)-.0624(.057207)\\
 &\hspace{4cm}=.0275811582 .
\end{split}
\]
The last term is the NONZERO
head/core cross.
Together with
\eqref{app:exact19:eq:Btightref},
the total is \(1.0000646582>1\).
Thus all three zero-own short-only
native terminal walls are impossible.


\subsection{The complete actual \(B_{19}\) last-long classification}
\label{app:exact19:Blastlong}

Release eighteen.
The retained old diagram is
\(A_{17}\) on one through seventeen
and the short \(A_1\) at nineteen.
A finite completion has at most
one contact to \(A_{17}\):
simple-three at \(1,2,16,17\),
or endpoint four at one/seventeen.
Exceptional finite diagrams cannot
contain \(A_{17}\).
A short-only contact can have ANY
finite native label.
If both retained components are
joined, full rank-nineteen
classification permits only a
simple-three at left endpoint
one/seventeen and short-three/four.
An internal left attachment
introduces a second inadmissible
fork; an internal four is not a
large finite \(B\) diagram.
This exhausts the actual finite
white list before any own cutoff.

Write the own product magnitude
as \(g\ge0\).
With no short contact, a white
three at \(j\) has
\[
 \Delta=1-j(18-j)/36,\quad\sigma=g+j/36;
\]
a white endpoint four has
\(\Delta=1/18\),
\(\sigma=g+j/(18\sqrt2)\).
The strong ratio target at
\(q=\sqrt{36}=6\) is \(<33/20\).

For completeness keep the entire
small-own list.
Own-zero far three has negative
time; far next-to-end three and
far four have zero time.
Near-three/next-to-end zero-own
has the actual affine two-leaf
fork ending in four, and near
four has the \(4,3,4\) principal.
Own three with a white three
has its actual affine cycle;
far-four/own-three has the affine
\(3,3,4,3\) path, while near
four/own-three is directly strong.
Own four without a short contact
has the \(4,4\) principal through
old eighteen/nineteen.
Own six has the left \(6,3\)
principal if its selected neighbour
is uncontacted; the other near
roles already have a directly
strong ratio.
Parallel has the positive future
null \(c+n_{18}\).
Every affine null has a positive
actual-cap coefficient, future
time and nonpositive products
with all omitted walls.

For every remaining \(g\ge\varphi/2\),
near three, next-to-end three
and endpoint four are directly
strong by the displayed residuals.
Only the pure far-three helper
\eqref{app:exact19:eq:Bpacket}
can remain weak on the left.
Its changed-loss cutoff is
\(g<6/5\), reproved by the actual
edge ratio, not inferred from an
old weak set.

An isolated cap at this final
index is also strong or impossible.
Use the LEFT old roots
eighteen/seventeen/sixteen.
Native five has divided finite
\(H_4\) loss
\[
        6-\sqrt{36-u}+\sqrt{u-3/2}<.13+.168<.58 .
\]
Native seven has the negative
left-\(A_2\) loss
\(6-\sqrt{36-u}+\sqrt{4/3-u}<.105+.319\).
Finite labels at least eight
have singleton loss
\(<.102+.415\).
For dotted own,
\[
 6-\sqrt{36-u}+\sqrt{u(1-u)}
       <.084u+\sqrt{u(1-u)}
       <(.084+1.004)/2<.58 .
\]
The positive finite points or
negative cores are the same actual
constructions proved above, now
at dual square thirty-six.
Their products with short nineteen
are nonpositive through the old
eighteen term.
The spatial \(\sqrt2P_{19}\) used
in a displacement is not an old
long normal or its short dual.

If both a left endpoint \(j=1,17\)
and the short component are
contacted, with short magnitude
\(\delta=1/2\) or \(1/\sqrt2\),
the exact data are
\[
             \Delta=19/36-\delta^2,\qquad
             \sigma=g+j/36+\delta/\sqrt2 .
\]
Short four is strong even at
own zero.
For short three and every
positive own,
\[
 \sigma/\sqrt\Delta\ge
 \frac{1/2+1/36+1/(2\sqrt2)}{\sqrt{5/18}}>33/20 .
\]
At own zero with \(j=1\), the
actual selected sixteen/seventeen/
eighteen/nineteen/cap path is
affine \(F_4\).
At \(j=17\), the cap is exactly
\((\omega_{17}+\omega_{19})/2\),
already excluded by its ACTUAL
short-edge mate in the boundary
proof.
No mixed-left/short weak cap
remains.

With no left contact write
\(v=g\omega_{18}+\delta\omega_{19}\).
Its outgoing ratio is
\[
                         \frac{g+\delta/\sqrt2}{\sqrt{1-\delta^2}} .
\]
For own zero this is exactly
\(\cot(\pi/m)/\sqrt2\).
Short three has the affine \(F_4\)
path and short four the double-four
principal.
At label eight the ratio is
\(1+1/\sqrt2>33/20\), increasing
thereafter.
The remaining short-five/six/seven
rows were proved impossible in
Section~\ref{app:exact19:Bzeroown}.

For positive own and short at
least five, the ratio is
\(>53/30>33/20\).
For short four and own at least
four it is \(\ge1+1/\sqrt2\);
only own-three/short-four remains.
For short three, own-three remains.
Own four has a finite \(B_{19}\)
replacement after releasing short
nineteen, with actual reciprocal
edge cosine \(55/2>27\).
Positive finite projection makes
that endpoint actual, giving
\(S_{19}^2>247/14>(4.2)^2\).
The early references
\eqref{app:exact19:eq:Bearly}
with the \(1/9\) price exclude
this pressure in both \(K\) rows.
Parallel again has its actual
positive feasible null.
We next exclude every finite
own at least five and dotted
own in the remaining short-three
row.

Put \(\alpha=\sqrt2g+1/2\);
its actual time is
\(A^2=18\alpha^2-3/4\).
For finite \(\varphi/2\le g<1\),
put \(d=1-g^2\), \(D=g/\sqrt2+1/2\),
\(\eta=d/D\), \(\lambda=g+\eta/\sqrt2\).
The actual positive combination
\[
                             \xi=c+\lambda n_{18}+\eta n_{19}
\]
has products zero with the cap
and eighteen, and
\(-D+\eta/2<0\) with nineteen.
Its square is \(-k\),
\(k=d-\eta^2/2>0\), because
\(2D^2-d=\alpha^2-3/4>0\).
It is future and satisfies every
omitted wall termwise.
Its normalized spatial displacement
\(Z\) is one on the first seventeen
coordinates, has square \(<18\),
and gives
\[
                 x\cdot Z\ge G=(A-\sqrt k)/\alpha>3.87 .
\]
Here is the full interval seal.
For \(\alpha>1.643\),
\[
 \frac{k}{\alpha^2}
   =\frac{d(2-1.5/\alpha^2)}{(\alpha+.5)^2}<.11,
 \qquad d=.875+\alpha/2-\alpha^2/2 .
\]
The numerator derivative is at
most \(-1.143(1.44)+.347(.677)<0\)
and its denominator increases.
At \(1.643\) the value is
\(<.347(1.445)/4.59<.11\).
Also \(A/\alpha>4.209\), while
\(\sqrt{.11}<.332\).
Thus the stated bound is strict
on every finite own in the range.

For ALL dotted \(g>1\), use instead
\(\xi=c+g^{-1}n_{18}\).
Its cap product is zero, both
old eighteen/nineteen products
are negative, its square is
\(-(1-g^{-2})\), and its time
is positive.
The normalized displacement again
has first seventeen coordinates
one and square \(<18\).
Since \(\alpha>1.9\),
\(A/\alpha>4.21\).
For \(t=1/g\in(0,1)\),
\[
 \frac{\sqrt{1-g^{-2}}}{\alpha}
       =\frac{t\sqrt{1-t^2}}{\sqrt2+t/2}<1/3 .
\]
For \(t\le1/2\), its square is
\(\le3/32\); for \(t\ge1/2\)
it is \(<100/1089\), both \(<1/9\).
Hence \(G>3.876>3.87\) throughout
the unbounded dotted interval.

Use the refined early references
of Section~\ref{app:exact19:Bcurvedearly}.
Every earlier residual is supported
through \(K\), so \(Z-P_K\)
vanishes on its ENTIRE support.
It is exactly orthogonal to all
early head/dotted-five corrections,
including their mutual nonzero
cross already paid there.
The remaining \(1/9\) endpoint
price is available.
The new pure tail gain exceeds
\[
 \begin{cases}
 2(3.87-3.463)/9-3/81>.0534,&K=15,\\
 2(3.87-3.590)/9-2/81>.0375,&K=16 .
 \end{cases}
\]
Each exceeds one when added to
EVERY strong/helper/isolated-dotted
five effective-energy row above.
This closes all short-three
finite-own-at-least-five and
dotted-own cases.

Thus the full classification
leaves either strong eighteen,
a pure far helper there, or
precisely the two weak walls
\[
       (b)\quad v_T=(\omega_{18}+\omega_{19})/2,\qquad
       (c)\quad v_T=\omega_{18}/2+\omega_{19}/\sqrt2 .
\tag{B19survivors}\label{app:exact19:eq:Bsurvivors}
\]

\subsection{Strong eighteen and a weak far helper are also excluded}
\label{app:exact19:Bendstates}

If eighteen is strong, the actual
window through seventeen justifies
the full curved references there.
A helper at five gives
\(E_{18}+.03125-1/1024>1\).
An isolated dotted five gives
the two-early-foot \(J=6,7\)
contradictions at that ACTUAL
strong endpoint.
For strong five,
\[
                      w^2>E_{18}-1/1024>.985613 .
\]
The smallest paid head gain
\(>.015456\) already contradicts
the norm.

If eighteen is a weak far helper,
seventeen is strong.
At this boundary, nonadjacency
uses the positive finite
\(D_{18}\) replacement from
the retained \(A_{16}\) and two
simple leaves, together with the
independent old short \(A_1\).
Both helper caps have zero
contact to the short root.
This checks the actual mixed
two-cap hypotheses; no virtual
long nineteenth root is used.

The full curved/early-foot
constructions at strong seventeen
exclude helper-five and isolated
dotted-five.
Hence five is strong.
Its reference through seventeen
has
\(w^2>.97133\), sum \(f_{17}\),
first-five mean \(b<.3652\),
\(P_5\) price \(>.151\) and
\(P_{17}\) price \(>1/9\).

For the actual weak eighteen
helper \(A^2=36g^2+2g-1/2\).
Let \(\omega_{19}^*=\sqrt2P_{19}\)
denote a SPATIAL vector only.
With \(\alpha=\sqrt2g\), the
complete actual pressure feet
are
\[
\begin{array}{c|c|c|c}
 \text{own}&U&x\cdot U\ge T&V^2,\ V=U-\alpha P_{17}\\ \hline
 5&g(\omega_{17}+\omega_{19}^*)/2+\omega_3/6&
          H(A,2/3-g^2)&g^2+1/6\\
 7,8&g(\omega_{17}+\omega_{19}^*)/2+\omega_2/4&
          A-\sqrt{g^2-3/4}&g^2+1/4\\
 \text{finite }\ge9&
       g(\omega_{17}+\omega_{19}^*)/2+\omega_1/2&
          H(A,1-g^2)&g^2+1/2\\
 \text{dotted}&
       (g-g^{-1})\omega_{18}
          +(\omega_{17}+\omega_{19}^*)/(2g)+\omega_1/2&
          A-\sqrt{1-g^{-2}}&2g^2-3/2+g^{-2}.
\end{array}
\]
Here \(H(A,\delta)\) is the actual
positive finite-projection pressure.
The normal corrections are respectively
\[
 g n_{18}+\tfrac23n_1+\tfrac13n_2,\quad
 g n_{18}+n_1/2,\quad
 g n_{18},\quad
 g^{-1}n_{18}.
\]
They use ONLY actual long roots
one/two/eighteen.
At native five their finite
\(H_4\) residual is \(2/3-g^2\);
at seven/eight the core has
negative square \(-(g^2-3/4)\);
the next finite pair has residual
\(1-g^2>0\), and the dotted core
has square \(-(1-g^{-2})\).
Every selected point is future,
and every omitted-wall product,
including that with actual SHORT
nineteen, is nonpositive.
The finite pressure comes from
a point, never from a spacelike
correction by itself.

All four rows satisfy
\[
          T-w\cdot U>.21,\qquad V^2<2.1,\qquad
                w\cdot U=\alpha f_{17}+b/\sqrt2 .
\tag{B19helper18}\label{app:exact19:eq:Bhelper18}
\]
Use \(\sqrt2f_{17}<5.251\),
\(b/\sqrt2<.259\).
Native five has
\(A>4.9678,\sqrt\delta<.111,g<.81\),
giving deficit \(>.3444\).
Native seven has
\(A>5.5245\), radical \(<.249\),
\(g<.901\), giving \(>.2853\).
Native eight has
\(A>5.663\), radical \(<.322\),
\(g<.9239\), giving \(>.2306\).
For finite labels at least nine,
at \(g=.939\),
\(A>5.75,\sqrt{A^2+\delta}>5.765\),
\(\sqrt\delta<.344\), so
\(H>5.42\) and deficit \(>.2303\).
The difference \(H-5.251g\)
increases thereafter:
\(A'>6,H_A>47/50,H_\delta<0\),
and \(\delta=1-g^2\) decreases.
For dotted \(1<g<6/5\),
the difference decreases since
\(A'<6.1\) and the radical
derivative \(>1\).
At \(6/5\),
\(A^2=2687/50,A>7.3307\),
radical \(<.5528\), giving
deficit \(>.2177\).
The finite \(V^2\) bounds are
immediate; the dotted square
increases to \(1867/900<2.1\).
These positive-square/derivative
seals cover every own interval.

Insert BOTH primitive head and
helper pressures in one vector
\[
                   W=w+\mu(Z-P_5/5)+\nu V,\qquad\nu=1/16 ,
\]
with the head-specific \(\mu\)
in \eqref{app:exact19:eq:Bheadgains}.
The individual prices are
\(\mu/5\) at five and
\(\nu\alpha<.107<1/9\) at seventeen.
The residual \(V\) has an initial
piece \(P_s/(s\sqrt2)\),
\(s=3,2,1\).
If \(\ell_s\) is the mean of the
first \(s\) coordinates of \(Z\),
then EXACTLY
\[
                         (Z-P_5/5)\cdot V=(\ell_s-1/5)/\sqrt2 .
\]
The total subtracted cross is
\(<.002476\) for native-five head,
\(<.010608\) for other finite head,
and \(\le.02\) for dotted head.
The helper gain from
\eqref{app:exact19:eq:Bhelper18}
is \(>.018046875\).
Thus the single squared-reference
totals exceed respectively
\[
\begin{split}
 .97133+.015456+.018046875-.002476&>1.0023,\\
 .97133+.02914+.018046875-.010608&>1.0079,\\
 .97133+.03384+.018046875-.020000&>1.0032 .
\end{split}
\]
No independent head optimum was
added to an already optimized
helper energy.
Strong eighteen and its pure-helper
weak case are now both impossible.

\subsection{Each terminal survivor supplies its own genuine strong seventeen}
\label{app:exact19:Bsurvivorhead}

Only the two actual walls in
\eqref{app:exact19:eq:Bsurvivors}
remain.
Their first eighteen coordinates
are \(a\), last coordinate
\(\delta\), with
\[
       (b):\ a=(\sqrt2+1)/2,\ \delta=1/2;\qquad
       (c):\ a=\sqrt2,\ \delta=1/\sqrt2 .
\]
Their time square is
\(A_T^2=18a^2+\delta^2-1\).

If seventeen were weak, let its
actual helper have parameter \(g\),
\(A_H^2=34g^2+2g-1/2\).
The mutual magnitude is
\[
                  M=A_HA_T-a(34g+1)/\sqrt2 .
\]
The actual walls are distinct by
their old-one products.
With \(\kappa=a^2+\delta^2-1\),
the positive squared difference
is \(A_H^2\kappa-9a^2\).
For \(g>.8\), it is positive in
both rows: in (b), \(\kappa=1/\sqrt2\)
and it exceeds
\(22.86(.707)-9(1.458)>3\);
in (c), it exceeds
\(22.86(1.5)-18>16\).
So \(M>0\).
The squared derivative comparison
\(34\kappa A_H^2+18A_T^2>0\)
shows that \(M\) increases with \(g\).

For (b), \(g<1.2\) gives
\[
                M<7.132(5.048)-.8533(41.8)
                           =.334396<1/2 ,
\]
using \(A_T<5.048,a/\sqrt2>.8533\).
This violates the native gap.
For (c), \(A_T=\sqrt{71/2}<5.959\)
and \(a/\sqrt2=1\), so
\[
                        0<M<7.132(5.959)-41.8=.699588<.7 .
\]
The only possible native magnitude
is \(1/2\).
But \(M(1)=1/2\) exactly, so strict
monotonicity forces \(g=1\).
The actual helper then has the
nonzero positive future null
\(c_H+n_{17}\), impossible in a
compact polytope.
Thus seventeen is strong in BOTH
survivors.

With a finite head at least seven
or a dotted head, the actual
strong-seventeen references now
close every initial-five branch:
helper five and isolated dotted
five by their complete constructions,
strong five by
\(.97133+.02914>1\) or
\(.97133+.03384>1\).
Therefore both survivors force
the isolated native-five long head,
its \(H_4\) pressure
\(x\cdot P_4/4>.393075\),
and the genuine strong
five/six/seven/seventeen sites.
These inputs precede the fresh
changed-loss arguments below.


\subsection{A common complete white list at the newly tested index seventeen}
\label{app:exact19:Bchangedlist}

Deleting seventeen retains \(A_{16}\)
on one through sixteen and \(B_2\)
on eighteen/nineteen.
The full finite-white list is
isolation; left simple-three at
\(1,2,15,16\); left endpoint-four
at one/sixteen; right simple-three
at eighteen or nineteen; or a
left endpoint-three joined to
either right simple-three.
Two contacts to one old tree
make a cycle.
A right-four or larger extension
of \(B_2\) is nonpositive.
Full bridges have rank nineteen,
so exceptional finite types,
an internal four and a second
fork supply no further alternatives.
There are no spectator components.

For a unit actual cap \(c=Ap-v\)
write \(g=-c\cdot n_{17}\ge0\).
The left-only exact data are
\[
\begin{array}{c|c|c}
 \text{white}&\Delta&\sigma\\ \hline
 3@j,\ j=1,16&9/17&g+j/34\\
 3@j,\ j=2,15&2/17&g+j/34\\
 4@j,\ j=1,16&1/17&g+j/(17\sqrt2).
\end{array}
\]
A full bridge has residual \(1/34\);
even at own zero its numerator
is \(.5+1/34\) for right-long
three, or \(1/(2\sqrt2)+1/34\)
for right-short three.
Both resulting ratios exceed two.
This complete list will be tested
AFRESH under each terminal wall;
its old \(.58\)-strong labels
do not imply the new floors.

The common actual small-own
reductions, when not replaced by
a direct mutual calculation below,
are as follows.
Zero-own far-three has negative
time, far next-to-end three and
far four have zero time.
The near next-to-end zero-own
has the affine two-leaf fork
ending in four; near-four
zero-own has the affine chain
ending in four at both ends.
Own-three with a white-three
has the actual all-three cycle.
Own-four has the selected
\(4,3,4\) path through
\(c,17,18,19\); own-six has the
\(6,3\) path through \(c,17,18\);
parallel has the positive feasible
future null \(c+n_{17}\).
Each null has positive time
through the actual cap and every
omitted-wall product nonpositive.
The endpoint-four positive-own
rows are directly tested below,
including own three.

\subsection{Short-four terminal: a new \(.56\) floor at the one actual index}
\label{app:exact19:Bshortfourfloor}

Assume survivor (c),
\(v_T=\sqrt2P_{18}+e_{19}/\sqrt2\),
\(A_T^2=71/2\).
We prove
\[
                        S_{17}>F=\sqrt{17}-.56/\sqrt2 .
\tag{B19F56}\label{app:exact19:eq:BF56}
\]
At \(q=\sqrt{34}<6\), its sufficient
ratio target satisfies
\[
 T_{.56}(q)=25/14-\frac{q-.56}{q(2q-.56)}
                      <25/14-34/429<1.707 .
\]
The first function increases on
the relevant range.

For isolated caps, the complete
small-own/time/affine/null
reductions leave five, seven,
finite at least eight, or dotted.
Use the actual LEFT corrections
at seventeen/sixteen/fifteen.
For native five, the divided
finite-\(H_4\) loss is
\[
 q-\sqrt{34-u}+\sqrt{u-3/2}<.133+.168<.56,
                     \quad u=8-4\varphi<1.528 .
\]
For native seven the left-\(A_2\)
negative core gives loss
\[
 q-\sqrt{34-u}+\sqrt{4/3-u}<.107+.319<.56,
                         \quad1.2318<u<1.233 .
\]
For finite at least eight the
actual pair projection gives
\(<.102+.415<.56\).
For dotted own the actual core
\(c+g^{-1}n_{17}\), together with
concavity, gives loss
\[
 q-\sqrt{34-u}+\sqrt{u(1-u)}
  <.087u+\sqrt{u(1-u)}
  <(.087+1.004)/2=.5455<.56 .
\]
The core/finite-point signs are
exactly those proved earlier;
all corrections use actual long
roots and have nonpositive
product with short nineteen.
Thus every isolated cap satisfies
the new floor, including native
eight.

For left-only rows, the small
owns in the common list are
removed first.
For all remaining \(g\ge\varphi/2>.809\),
near endpoint-three has ratio
\[
       (g+8/17)/\sqrt{9/17}>1.2795/.728>1.757 .
\]
The next-to-end rows exceed \(2.5\)
and endpoint-four exceeds \(3.5\).
Only far-three at one needs its
actual terminal interaction.

For that actual cap,
\(v_H=g\omega_{17}+\omega_1/2\),
\(A_H^2=34g^2+2g-1/2\), and
\[
                    M(g)=A_H\sqrt{71/2}-(34g+1)=-c_H\cdot c_T .
\]
The two walls are distinct by their
old-one products.
For \(g>.8\), its squared-product
difference is \(3A_H^2/2-18>0\),
so \(M>0\).
Its squared derivative comparison
\(51A_H^2+18(71/2)>0\) proves strict
increase.
We have \(M(1)=1/2\).
At \(g=243/200=1.215\),
\[
 A_H^2=52.12165<(7.21954)^2,\quad
 \sqrt{71/2}<5.9582,\qquad
 M<7.21954(5.9582)-42.31=.705463228<.707<\cos(\pi/4).
\]
Thus for \(g\le1.215\) the only
possible native magnitude is \(1/2\);
strict increase forces the forbidden
parallel helper \(g=1\).
Every remaining actual far helper
has \(g>1.215\), hence
\[
                      \frac{34g+1}{\sqrt{612}}
                         >42.31/24.739>1.71>T_{.56}(q).
\]
This tests the entire new cap
domain; the earlier \(.58\) cutoff
\(g<1.2\) has not been imposed.

For right-long three, own zero
has its affine star and own three
its affine triangle.
Every remaining own has
\(g\ge1/\sqrt2\), ratio
\(\ge1+1/\sqrt2>1.707\).
For right-short three, own
zero/four/six/parallel has the
actual affine/null obstruction.
Own-three and own-five are
impossible by the previous
ACTUAL short-edge mate and left
\(H_4\) boundary arguments, which
did not require the new seventeen
floor.
Every other finite own at least
seven has ratio
\(\sqrt2g+1/2>1.77\); dotted
owns are stronger.
Finally all full bridges have
ratio \(>2\), including own zero.
The complete new finite-white/
own domain is exhausted, proving
\eqref{app:exact19:eq:BF56}.

\subsection{Short-four terminal: pool the tail before raising its endpoint}
\label{app:exact19:Bshortfourfinish}

Use the old \(.58\) reference
through seventeen, but deliberately
omit fifteen/sixteen as well as
every actual weak knot.
Because helpers are nonadjacent,
the last retained \(K\) before
seventeen is fourteen or thirteen.
If fourteen is weak, it merges
with the forced terminal cluster;
all earlier omissions stay isolated.

All earlier holes cost \(<.000181\).
The terminal cluster averages
at most four slopes
\(d_{13},\ldots,d_{16}\), whose
successive gaps are
\[
                         <1/(4\cdot13\sqrt{13})<.0054 .
\]
For four numbers with successive
gaps at most \(\kappa\), their
variance about the mean is
\(\le5\kappa^2\), by
\((1/4)\sum_{i<j}(a_i-a_j)^2\)
and distances one/two/three.
The shorter three-slope cluster
costs at most \(2\kappa^2\).
Thus the terminal loss is
\(<.000146\), and FULL loss
\(<.00035\).
The affected open intervals
are disjoint even when they
share a retained endpoint.

Now raise only seventeen by
\(\Delta=.02/\sqrt2\),
\(.01414<\Delta<.01415\).
Both new terminal slopes are
\(<.133\):
\[
 \frac{\sqrt{17}-\sqrt{14}+\Delta}{3}<.133,\qquad
 \frac{\sqrt{17}-\sqrt{13}+\Delta}{4}<.133 .
\]
Their incoming slopes are at
least \(d_{13}>.136\) and
\(d_{12}>.14\), including any
earlier isolated omission.
The reference remains concave,
with every coefficient nonnegative.
Its endpoint coefficient is
\(>d_{16}>.123\); its first-five
block and \(P_5\) coefficient
\(>.151\) are unchanged.

The old terminal chord has slope
\(\ge d_{16}>.123\).
Lifting its endpoint adds
\[
                  2(\text{old slope})\Delta+
                         \Delta^2/(17-K)>.246(.01414)=.00347844 .
\]
Together with the scalar
\eqref{app:exact19:eq:BE17},
\[
             w^2>.9726645-.00035+.00347844=.97579294,
                       \qquad\sum_jw_j=F .
\]
Only seventeen uses the new floor;
all other active floors are
original strict \(.58\) floors.

The SAME actual terminal cap is
also the short-nineteen mate.
Deleting old nineteen and adjoining
it to the retained \(A_{18}\) by
its actual simple-three at eighteen
gives positive finite \(A_{19}\).
Its positive inverse projection
is future, strict against every
unknown original wall.
The omitted old nineteen starts
at zero and becomes strict through
the positive cap coefficient and
cap-short product \(-1/\sqrt2\).
Thus this is an actual endpoint
of that original edge, with
\[
                    \cosh L_{19}=37/\sqrt{20},\qquad S_{19}>3.860 .
\]
For the last bound, using
\(\sqrt{20}<4.47214\), the exact
cleared comparison is
\[
 (19-14.8996)37=151.7148
              >(19+14.8996)4.47214=151.603757144 .
\]
We do not assume this neighbouring
\(A_{19}\) vertex is farthest.

Insert the native-five head
pressure by
\(w_{\rm head}=w+.56(P_4/4-P_5/5)\).
The individual \(P_5\) price
is \(.112<.151\).
Its residual has square \(1/20\),
zero sum on the first-five block
and zero cross with \(w\).
Its gain is
\[
                         1.12(.027885)-.56^2/20=.0155512 .
\]
The modified reference still has
support through seventeen, sum
\(F\), and \(P_{17}\) coefficient
\(>.123\).

Extend it by the last-two slope
\(r=(3.860-F)/2\).
The bounds \(3.720<F<3.72713\)
give \(0<r<.07\), so its
\(P_{17}\) debit fits.
The head residual and new tail
residual have disjoint support
and cross zero.
One final square yields gain
\[
             .0155512+(3.860-F)^2/2
                     >.0155512+(.13287)^2/2 .
\]
Consequently
\[
                         \|x\|^2>
 .97579294+.0155512+.00882721845=1.00017135845>1 .
\]
This excludes the entire
short-four terminal survivor.

\subsection{Short-three terminal: retain the orthogonal dotted exception}
\label{app:exact19:Bshortthreefloor}

Assume survivor (b).
Put \(s=\sqrt2\), \(a=(1+s)/2\),
\(k=as=1+s/2\).
Then
\[
 v_T=aP_{18}+\tfrac12e_{19},\qquad
 A_T^2=51/4+9s,\qquad
                         34A_T^2-289k^2=17s .
\]
The actual index-seventeen cap
is distinct from \(c_T\): the
latter contacts both retained
\(B_2\) roots with inverse cost
\(1+s/2>1\), impossible for that
cap's positive finite white block.
Write
\[
          L=v\cdot v_T>0,\quad M=AA_T-L=-c\cdot c_T,\qquad
                    D=A^2A_T^2-L^2 .
\]
Actual acuteness/native products
give \(M=0\) or \(M\ge1/2\).
Since \(A,A_T,L>0\), \(D\) has
the sign of \(M\).
In particular \(D<0\) means a
POSITIVE actual-wall product.

We prove the fresh floor
\[
                          S_{17}>F=\sqrt{17}-.50/\sqrt2 .
\tag{B19F50}\label{app:exact19:eq:BF50}
\]
Its sufficient ratio target is
\(T_{.50}(q)=2-(q-.50)/(q(2q-.50))<2\),
\(q=\sqrt{34}\).

For isolated caps,
\[
 v=g\omega_{17},\quad A^2=34g^2-1,\quad
             L=17kg,\qquad D=17sg^2-A_T^2 .
\]
Own zero has negative time.
Every positive finite own and
parallel has \(g\le1\), giving
\(D\le8s-51/4<0\); all violate
acuteness.
Thus an isolated cap is dotted.

The function
\(M(g)=A_T\sqrt{34g^2-1}-17kg\)
strictly increases:
its squared derivative comparison is
\(578sg^2+289k^2>0\).
If \(M=0\), the exact parameter is
\[
                          u=g^{-2}=17s/A_T^2,\qquad .94<u<.95 .
\]
Orthogonality is RETAINED.
The actual positive core
\(c+g^{-1}n_{17}\) has future time,
square \(-(1-u)\), zero cap
product, negative own product
and every omitted-wall product
nonpositive.
Its spatial displacement is
\[
 Y=g\{(1-u)\omega_{17}
                  +(u/2)(\omega_{16}+\omega_{18})\}.
\]
Original concavity gives
\(x\cdot Y\le g h_{17}\).
Farthest comparison consequently
gives
\[
                 h_{17}\ge\sqrt{34-u}-\sqrt{u(1-u)} .
\]
The loss below \(\sqrt{34}\) is
\(<.087u+\sqrt{u(1-u)}
<.083+.238=.321<.50\),
so the orthogonal case satisfies
the new floor.

If \(M>0\), then \(M\ge1/2\).
At \(g=7/4\),
\[
 (119k/4+1/2)^2-A_T^2(825/8)
                         =(684-451s)/16>0 ,
\]
so \(M(7/4)<1/2\).
Strict increase forces \(g>7/4\)
and \(u<16/49\).
The same actual foot has loss
\[
          <.087(16/49)+4\sqrt{33}/49<.029+.469=.498<.50 .
\]
This covers every isolated dotted
parameter and all mutual labels,
including zero.

For left-only roles use the
complete residual table of
Section~\ref{app:exact19:Bchangedlist}.
Keep its actual zero-own and
affine/null reductions, except
near endpoint-three at sixteen,
which will be tested directly.
Every positive-own endpoint-four
has ratio
\(\ge\sqrt{17}/2+1/\sqrt{34}>2\).
Every remaining next-to-end
three has \(g\ge\varphi/2>.809\)
and ratio \(>2.5\).

For near endpoint-three at sixteen,
\[
\begin{gathered}
 v=g\omega_{17}+\omega_{16}/2,\quad
 A^2=34g^2+32g+7,\quad L=k(17g+8),\\
                   D=17sg^2+16sg-27/4-s .
\end{gathered}
\]
At \(g=0\), \(D<0\).
For finite \(g\ge1/2\), \(D>0\),
starting with \((45s-27)/4>0\).
The identity
\((34g+16)^2=34A^2+18\)
gives the positive derivative
comparison \(17sA^2+18A_T^2>0\)
for \(M\).
At \(g=1\),
\[
                         (25k+1/2)^2-73A_T^2=32-39s/2>0 .
\]
Thus \(0<M<1/2\) for all finite
owns and parallel, impossible.
Every remaining dotted \(g>1\)
has ratio \(>25/\sqrt{153}>2\),
since \(625>612\).

For far endpoint-three at one,
the common small-own/null
reductions leave \(g\ge\varphi/2>.8\).
The exact data are
\[
\begin{gathered}
 A^2=34g^2+2g-1/2,\quad L=k(17g+1/2),\\
                   D=17sg^2+sg-27/4-19s/4 .
\end{gathered}
\]
This polynomial is positive at
\(.8\) and increases thereafter,
so \(M>0\).
The identity
\((34g+1)^2=34A^2+18\)
again gives the positive derivative
comparison.
At \(g=3/2\),
\[
                         (26k+1/2)^2-79A_T^2=33-22s>0 .
\]
Thus every admissible remaining
helper must have \(g>3/2\),
whose ratio exceeds
\(52/\sqrt{612}>2\).
This tests the full role, not
only its former \(.58\)-weak
parameters.

For right-long three at eighteen,
the exact data give
\[
\begin{gathered}
 A^2=34g^2+34g+8,\quad L=k(17g+9),\\
                   D=17sg^2-(51/2)g-39/2-9s .
\end{gathered}
\]
The convex quadratic is negative
at both zero and one; at one
it is \(8s-45<0\).
It is therefore negative on the
WHOLE interval \([0,1]\).
Own zero, all finite owns and
parallel violate acuteness.
A remaining dotted \(g>1\)
has ratio \(\sqrt2g+1/\sqrt2>2\).

For right-short three at nineteen,
\[
\begin{gathered}
 A^2=34g^2+17sg+15/4,\quad
 L=17kg+(19+18s)/4,\\
                   D=17sg^2-(17/2+17s)g-61/4-9s .
\end{gathered}
\]
This convex quadratic is negative
at zero and \(5/4\), the latter
value being \(-207/8-59s/16\).
Thus only \(g>5/4\) can survive
actual acuteness, with ratio
\(\sqrt2g+1/2>2\).
Finally all full bridges are
directly strong by their ratio
\(>2\), including own zero.
This exhausts every actual white
and own parameter, and proves
\eqref{app:exact19:eq:BF50}.
The \(.56\) floor from the other
terminal case was not used.

\subsection{Short-three terminal: pool every possible final weak cluster}
\label{app:exact19:Bshortthreefinish}

Keep the old \(.58\) floors,
omit every actual weak knot,
and deliberately omit thirteen
through sixteen as well.
Since at least one of eleven/
twelve is strong, the last
retained \(K\) before seventeen
is twelve if twelve is strong,
otherwise eleven.
This includes every old weak set.

Earlier isolated omissions cost
\(<.000181\).
The terminal cluster averages
at most six slopes
\(d_{11},\ldots,d_{16}\).
Their consecutive gaps are
\[
                        <\kappa=1/(4\cdot11\sqrt{11})<.00686 .
\]
For \(n\) such numbers the exact
variance identity gives
\[
 \sum_i(a_i-\overline a)^2
 =\frac1n\sum_{i<j}(a_i-a_j)^2
                         \le n(n^2-1)\kappa^2/12 .
\]
At \(n=6\), this is
\((35/2)\kappa^2<.000824\);
the five-slope case is smaller.
The clusters have disjoint open
edge intervals, so TOTAL loss
is \(<.00101\).

Lift seventeen by
\(\Delta=.08/\sqrt2\),
\(.05656<\Delta<.05657\).
Both outgoing terminal slopes
are \(<.144\):
\[
 \frac{\sqrt{17}-\sqrt{11}+\Delta}{6}<.144,\qquad
 \frac{\sqrt{17}-\sqrt{12}+\Delta}{5}<.144 .
\]
Their incoming slopes are
\(\ge d_{10}>.154\), respectively
\(\ge d_{11}>.147\), including
any earlier isolated omission.
Thus the final interpolation is
concave, paid and supported on
old strong floors, with only
seventeen using the new \(.50\)
floor.
No assertion is made about the
unknown thirteen--sixteen floors.

The old last chord has slope
\(\ge d_{16}>.123\), so its
endpoint lift adds
\(>.246(.05656)=.01391376\).
The new reference satisfies
\[
                          w^2>.9726645-.00101+.01391376=.98556826 .
\]
Its first-five mean is still
\(b<.365190\), with individual
\(P_5\) coefficient \(>.151\).
Insert the actual native-five
\(H_4\) pressure in the SAME
vector
\[
                               W=w+.56(P_4/4-P_5/5).
\]
Its price \(.112\) fits.
Its residual has square \(1/20\),
zero block sum and zero cross
with \(w\).
The one original-vector square
gives
\[
 \|x\|^2>.98556826+1.12(.027885)-.56^2/20
                              =1.00111946>1 .
\]
No short-edge energy or second
farthest vector is added.
This excludes survivor (b).

\begin{proposition}\label{app:exact19:prop:wholeB19}
A whole \(B_{19}\) forest cannot
occur at the original farthest
vertex of a compact Coxeter
nineteen-polytope.
\end{proposition}
\begin{proof}
The actual window and full
last-long finite-white classification
give strong eighteen, a weak pure
far helper there, or the two walls
\eqref{app:exact19:eq:Bsurvivors}.
Every zero-own short-only
exception was excluded using its
own actual head/terminal points.
The strong and helper end states
are impossible.
Each survivor supplies its genuine
strong seventeen and isolated
native-five head before the fresh
single-index arguments use them.
The separate short-three and
short-four proofs then close both
remaining walls.
Every own, white, finite, parallel
and dotted role has been kept,
and every point is feasible against
the entire original polytope.
\end{proof}


\subsection{Whole \(D_{19}\): the original signed frame and a fresh head reduction}
\label{app:exact19:Dhead}

Use unit actual roots and duals
\[
\begin{gathered}
 n_i=(e_i-e_{i+1})/\sqrt2\ (i\le18),\quad
 n_{19}=(e_{18}+e_{19})/\sqrt2,\quad
 \omega_i=\sqrt2P_i\ (i\le17),\\
 \omega_{18}=(P_{18}-e_{19})/\sqrt2,\qquad
 \omega_{19}=(P_{18}+e_{19})/\sqrt2 .
\end{gathered}
The original vector satisfies
\(x_1>\cdots>x_{18}>|x_{19}|\),
\(\|x\|<1\).
The last coordinate remains signed.
The centre is seventeen and the
two actual forks are eighteen/
nineteen; no artificial long
fork normal is introduced.

At loss \(.58\), the complete local
whole-diagram lemma
\ref{app:tail:lem:whole58local}
applies literally at five through
fourteen: the retained right
\(D_{19-i}\) has rank at least five.
It retains every finite fork-to-
\(E_6,E_7,E_8\) extension;
a surviving positive-own fork
ratio is
\[
                >(.8+.25)/\sqrt{3/8}>83/50 .
\]
Thus every weak cap there is the
pure far-three helper
\[
        v=g\omega_i+\omega_1/2,\quad
        A^2=2ig^2+2g-1/2,\quad
        \varphi/2\le g<6/5 ,
\]
with own five, finite at least
seven or dotted, except for
isolated dotted five.
Actual helper nonadjacency uses
the retained positive right block
including BOTH fork normals.
Hence some \(K=13\) or fourteen
and some \(J=6\) or seven is strong,
before any head or late-layer
conclusion.
No fifteen/sixteen/centre floor
has yet been used.

Let \(f_i=\sqrt i-.58/\sqrt2\),
\(E_J=f_5^2/5+\sum_{i=5}^{J-1}d_i^2\).
The old \(E_{13}>.88989\), and
changing \(.61\) to \(.58\)
adds \(-.00357+.006\sqrt{10}>.0154\).
Thus \(E_{13}>.90529\).
The elementary bounds needed below
are
\[
\begin{gathered}
 E_6<.7126,\quad E_7<.7510,\quad
 2.0393<f_6<2.0394,\quad2.2356<f_7<2.2357,\\
 f_5/5<.3652,\qquad d_6<.197,\quad d_7<.183 .
\end{gathered}
\]
They follow, for example, from
\(f_5^2/5=1.03364-.116\sqrt{10}<.6668248\),
\(d_5^2<.0456,d_6^2<.03852\).
All remaining internal holes cost
\[
 \sum\frac1{32(j-1)^3}
 <1/4000+1/10976+1/23328+1/42592
                               <.000408<1/1024 .
\]
Only actual weak knots are omitted.

For a true first-coordinate floor
\(x_1>X\), use the reference
through \((1,X),(J,f_J)\), then
every retained strong knot through
the actual \(K=13,14\).
Its first slopes are \(X\),
then \(m=(f_J-X)/(J-1)\).
For \(X=.643,.668\), they are
positive decreasing and \(m>d_J\).
The single-reference energy is
\[
                       >E_{13}-1/1024+C_J(X)-E_J,
             \qquad C_J(X)=X^2+(f_J-X)^2/(J-1).
\]
At \(X=.643\) this gives
\[
\begin{array}{c|c|c|c}
 J&C_J(X)>&q_J=w^2>&\text{individual }(P_1,P_J)\text{ prices}>\\ \hline
 6&.803379738&.99509&(.3637,.08226)\\
 7&.8361781266&.98949&(.3775,.08243).
\end{array}
\tag{D19headrefs}\label{app:exact19:eq:Dheadrefs}
\]
Indeed the prices are \(X-m\)
and at least \(m-d_J\);
omitting the next knot increases
the latter.
At \(X=.668\), the corresponding
\(C_6>.822316738,C_7>.8557856266\)
already give reference energies
\(>1.01403,1.00909\).
These are actual original-vector
references, not a second energy
budget.

Deleting old one retains \(D_{18}\).
Its complete finite one-node list
is isolation or simple-three at
its long leaf two.
Two contacts make a cycle;
an additional fork/internal double
is not a large finite \(D\)
extension; a short-fork attachment
would require an exceptional
finite rank nineteen.
For a nonisolated head,
\[
 v_L=g_L\omega_1+\omega_2/2,\quad
 A_L^2=2g_L^2+2g_L,\quad
                   \cosh L_1=2g_L+1 .
\]
Own zero has zero time, three
has the actual affine triangle,
and parallel has its feasible
positive null.
Every other own has
\(g_L\ge1/\sqrt2\), hence
\(x_1>\sqrt{\sqrt2-1}>.643\).
For finite own at least five
and dotted own,
\[
                     x_1^2\ge g_L/(g_L+1)>.668^2 .
\]
If new five is weak, choose the
independently available strong
\(J=6,7\) and use the preceding
\(X=.668\) reference, contradiction.
If new five is strong, even
\(X=.643\) suffices: graft it
into the first-five block of mean
\(<.3652\), preserving its sum.
The remaining coordinates stay
\(>.2957>d_5\), and the gain is
\[
                       (5/4)(.643-.3652)^2=.09646605 .
\]
The full reference then gives
\(.90529-1/1024+.09646605>1.00077\).
Thus only the nonisolated native-four
head could remain:
\[
                  g_L=1/\sqrt2,\quad A_L^2=1+\sqrt2,\quad x_1>.643 .
\]

If five is a pure weak helper,
let its own be \(r\), with
\(A_5^2=10r^2+2r-1/2\).
The actual distinct-wall magnitude is
\[
             M(r)=A_LA_5-(2+\sqrt2)r-(1+\sqrt2)/2 .
\]
Its squared-product difference is
\[
 (4+6\sqrt2)r^2-(2+\sqrt2)r-5/4-\sqrt2 ,
\]
positive already at \(.8\) and
increasing.
Thus \(M>0\).
The identity
\((10r+1)^2=10A_5^2+6\)
gives a positive squared derivative
comparison
\((4+6\sqrt2)A_5^2+6(1+\sqrt2)>0\).
So \(M\) strictly increases.

At native seven,
\[
 M(r_7)<1.5538(3.0693)
             -(3.4142(.9009)+1.2071)=.48612556<1/2 .
\]
At native eight, positive squaring
shows \(M(r_8)>1/2\):
the required difference is equivalent
to
\(3+(5/2)\sqrt2>2(2+\sqrt2)r_8\),
whose square difference is
\(3/2+\sqrt2>0\).
Finally \(M(1)<\sqrt2/2\), since
the positive square difference is
\(11/4-3\sqrt2/2>0\).
There is no own label between
seven and eight.
Every finite helper own therefore
puts the mutual strictly in
\((0,1/2)\) or
\((1/2,\cos(\pi/4))\), impossible.
This is a discrete native-angle
argument over the real numbers.

For dotted helper own \(r>1\),
its exact weak ratio at five is
\((10r+1)/\sqrt{60}\).
Since
\[
 T(\sqrt{10})<T(16/5)=50/29-655/4656<317/200
\]
and the ratio at \(r=1.13\)
has square \(15129/6000>(317/200)^2\),
weakness forces \(r<1.13\).
Put \(t=1/r\), so \(t>.8849\).

The ACTUAL positive future core
\(c_5+r^{-1}n_5\) has square
\(-(1-t^2)\), zero cap product,
negative own product and every
other actual-wall product nonpositive,
including that with the head.
Its normalized pressure is
\[
 \begin{gathered}
 Z=P_5+(t^2/2)(e_6-e_5)+(t/2)P_1,\\
 x\cdot Z\ge
 T(t)=\{\sqrt{10+2t-t^2/2}-t\sqrt{1-t^2}\}/\sqrt2 .
 \end{gathered}
\]
Use the \(J=6\) reference of
\eqref{app:exact19:eq:Dheadrefs};
six is strong by helper
nonadjacency.
Projection to its first coordinate
and constant next-five block gives
\[
 \Pi Z=(1/5+t/2)P_1+(4/5)P_6,\qquad
 Q=\|Z-\Pi Z\|^2=4/5-t^2+t^4/2<.35 .
\]
Its true reference pairing is
\(w\cdot Z=.8f_6+(.2+t/2)(.643)
<1.76012+.3215t\).
The complete continuous deficit
satisfies \(T(t)-w\cdot Z>.0482\).
Indeed \(T(t)-.3215t\) increases
for \(t\ge7/8\): the derivative
contains
\((2t^2-1)/(\sqrt2\sqrt{1-t^2})>1/\sqrt2\)
and a positive further root
derivative.
At \(t_0=.8849\), the two radicals
give \(T(t_0)>2.0929\).
Thus the deficit is
\(>2.0929-.3215(.8849)-1.76012
=.04828465\).
The square checks use
\(.783<t_0^2<.784\),
\(10+2t_0-t_0^2/2>3.373^2\)
and \(.784(.217)<.413^2\).

Insert \(w+.1(Z-\Pi Z)\).
It spends at most \(.07\) at
\(P_1\) and exactly \(.08\) at
\(P_6\), both below the separate
available prices.
The residual is orthogonal
to this NONCONSTANT first reference.
One square gives
\[
                         \|x\|^2>.99509+.2(.0482)-.01(.35)
                                           =1.00123>1 .
\]

If five is instead isolated dotted,
its actual first-six core supplies
\[
 Z_0=P_5+(u/2)(e_6-e_5),\qquad
                   x\cdot Z_0>T_0>1.8185,\quad2/5<u<1 .
\]
These are the original local
roots; no \(B_{20}\) conclusion
has been transferred.
At a strong \(J=6\),
\[
 \Pi_6Z_0=P_1/5+4P_6/5,\quad
 Q_6<.48,\quad w\cdot Z_0<1.76012,\quad
                            T_0-w\cdot Z_0>.05838 .
\]
The coefficient \(.1\) costs
\(.02\) at one and \(.08\) at
six, giving
\[
                         \|x\|^2>.99509+.2(.05838)-.01(.48)>1.001966 .
\]
At strong \(J=7\),
\[
 \Pi_7Z_0=P_1/3+2P_7/3,\quad Q_7<76/75,\quad
             w\cdot Z_0<1.7048,\quad T_0-w\cdot Z_0>.1137 .
\]
Coefficient \(.12\) costs \(.04\)
at one and \(.08\) at seven,
giving
\[
                    \|x\|^2>.98949+.24(.1137)-.0144(76/75)>1.002186 .
\]
The \(P_1\) projections and debits
are essential; a constant-block
projection \(Z_0-(5/J)P_J\) would
be wrong for these references.

The nonisolated head is now
completely excluded.
The head is isolated with time
\(A_L^2=2g_L^2-1>0\).
Own zero/three/four has nonpositive
time; six has the actual affine
\(6,3\) path; parallel the
nonzero positive future null.
Its complete remaining own list
is native five, finite at least
seven, or dotted.
No centre, tail or outgoing-fork
floor was assumed in this proof.

\subsection{An independent old strong-five floor in whole \(D_{19}\)}
\label{app:exact19:Doldfive}

We next prove
\[
                              S_5>\sqrt5-.61/\sqrt2
\tag{D19old5}\label{app:exact19:eq:Doldfive}
\]
without any head hypothesis.
It is the OLD \(.61\) floor,
not the new \(.58\) floor.

Assume contrariwise old five is
weak.
The rank-valid generic inventory
at \(M=11\), \(19\ge M+6\),
gives a pure helper
\[
 c_5=Ap-g\omega_5-\omega_1/2,\quad
 A^2=10g^2+2g-1/2,\quad
                 \varphi/2\le g<6/5 ,
\]
own five, finite at least seven,
or dotted.
Its actual first-six points give
\[
                         x_1>a_0=.5965,\qquad S_6>B_0=2.01815 .
\]
There is an independent old strong
\(K=11\) or twelve.
Let \(f_i=\sqrt i-.61/\sqrt2\),
\(F_K=B_0+\sqrt K-\sqrt6<f_K\).
The interpolation through
\((1,a_0),(6,B_0),(K,F_K)\)
has positive decreasing slopes
\(a_0,.28433,\ldots,.28433,
1/(\sqrt K+\sqrt6)\).
Its precise paid data are
\[
\begin{array}{c|c|c|c}
 K&q_K>&F_K<&\text{individual endpoint price}>\\ \hline
 11&.9104&2.886&1/7\\
 12&.9300&3.034&1/7 .
\end{array}
\]
The first-six energy is
\(.7600299945\).
The eleven tail exceeds
\(.8671^2/5\); the twelve tail
exceeds \(.17\).
A pressure \(V=1\) through \(K\),
\(x\cdot V\ge G,V^2<R\), therefore
adds legally
\[
                            2(G-F_K)/7-(R-K)/49 .
\tag{D19weak5ref}\label{app:exact19:eq:Dweak5ref}
\]
This is payment at the actual
anchor, not subtraction of two
independent lower bounds.

Under this assumed weak five
we can classify the two late
layers afresh.
Releasing fifteen retains
\(A_{14}+D_4\); releasing sixteen
retains \(A_{15}+D_3\).
The right \(D_4\) permits simple
three at any leaf.
The right \(D_3=A_3\) permits
middle-three, end-three/four/five.
Two contacts to one right piece
make a cycle.
The long left path admits only
the usual endpoint/next-to-end
extensions, and the full rank-
nineteen bridges are \(D_{19}\)
or the end-three \(A_{19}\)
bridge.
No exceptional full diagram,
internal four or second fork
is missing.
There is no foreign component.

Here are explicit one-sided
seals before the new right rows.
For \(i=15,16\), an isolated
own has \(q^2=2i\ge30\).
The actual left \(H_4\), left
\(A_2\), singleton and dotted
cores give losses respectively
\[
 <.15+.168,\qquad <.12+.319,\qquad
 <.11+.415,\qquad
 <.093u+\sqrt{u(1-u)}
                 <(.093+1.005)/2 .
\]
Each is below \(.61\).
The finite ranges are
\(u_5<1.528\), \(1.2318<u_7<1.233\),
and \(1<u_{\ge8}<1.172\);
positive squaring at \(q^2=30\)
proves the first loss terms.
The native-five concavity correction
uses \(i+1\le17\), an ACTUAL
long dual here.

For left-only contact at \(j\),
the exact data are
\(\sigma=g+j/(2i)\) and
\(\Delta=1-j(i-j)/(2i)\)
for simple-three, or
\(\sigma=g+j/(\sqrt2i)\),
\(\Delta=1/i\) for endpoint-four.
After the actual own-zero,
three/four/six affine and
parallel-null reductions, near
endpoint-three has
\(\sigma>1.266,\sqrt\Delta<.731\);
next-to-end three has
\(\sigma>.8,\sqrt\Delta<.366\);
endpoint-four has a larger ratio.
The positive-own-three near-four
row is directly strong, while
far-four/own-three retains its
affine principal.
These ratios exceed \(25/16\).
Only the pure far LEFT helper
is retained.
Own parallel has its feasible
positive future null.

At fifteen the right LONG-leaf
three has numerator \(g+1/2\),
residual \(1/2\).
Own zero has the actual affine
double fork on
\(c,15,16,17,18,19\), kernel
\((1,1,2,2,1,1)\).
Own three is an actual triangle;
every own at least four has
ratio \(\ge1+1/\sqrt2\).
At sixteen the right middle-three
has its analogous own-zero
four-leaf affine star, own-three
triangle and same strong ratio
for larger owns.
Right end-four at own zero has
the weighted affine principal
on \(c,f,17,\overline f,16\),
kernel \((\sqrt2,2,2,1,1)\);
positive own has ratio
\(\ge1+1/\sqrt2\).
End-five is the sealed finite
\(H_4\), directly strong EVEN
at own zero.
Full bridges have their own-zero
cycles or middle-three star.
For positive own, the worst
end-three bridge has residual
\(5/32\), numerator \(>3/4\);
the other bridge residuals are
smaller.
The old target is \(<25/16\)
through these indices, because
\(T(\sqrt{32})<T(6)<25/16\).
All these rows are strong.

The remaining one-sided right
fork-three at fifteen or end-three
at sixteen has own-zero affine
\(E_8\) with arms five/two/one.
Own three gives the actual FIVE-
cycle at fifteen, FOUR-cycle
at sixteen.
Own four has the affine star
centred at the released root,
with squared leaf weights
\(1/2,1/4,1/4\).
Own six has its actual \(6,3\)
principal, and parallel its
feasible positive null.

For native five use the LEFT
actual \(H_4\) projection on
\(c,n_i,n_{i-1},n_{i-2}\).
Its positive correction gives
\[
 \alpha^2=\varphi+5/8>56/25,\quad
 \delta_5=(5-3\varphi)/8,\quad
 V_1=\cdots=V_{i-3}=1,\quad V^2<i,
\]
\[
 x\cdot V>G>
       \sqrt{i-\kappa/\alpha^2}-.091,\qquad
       \kappa=1/2\ (i=15),\quad5/8\ (i=16).
\]
The actual finite point is future
and all omitted-wall products
are nonpositive.
The displacement is not presumed
to be a point by itself.
The demands are \(>3.753,3.873\),
from
\(3.844^2<15-25/112\),
\(3.964^2<16-125/448\).
They are flat on BOTH early
reference supports.
Even the conservative
\[
                      .91+2(3.75-3.034)/7-5/49>1.0125
\]
contradicts the norm.

At the \(D_4\) row every finite
own at least seven or dotted
own is directly strong:
\(\sqrt2(g+1/4)>1.626>25/16\).
At the \(D_3\) row its ratio is
\((g+1/4)/\sqrt{5/8}\).
Weakness forces \(g<1\), so
dotted own is directly strong.
For the remaining \(.9<g<1\),
use the actual core
\[
                          \xi=c+g^{-1}n_{16}+(2g)^{-1}n_{15}.
\]
Its cap and fifteen products
vanish, its sixteen product is
\(-g+3/(4g)<0\), and its square
is \(-k\), \(k=1-3/(4g^2)>0\).
It is future and satisfies BOTH
original forks and every other
wall termwise.
At this actual rank,
\[
 \alpha=\sqrt2(g+1/4),\quad A^2=16\alpha^2-5/8,\quad
          V_1=\cdots=V_{14}=1,\quad V^2<16 .
\]
Since \(\alpha^2>529/200\),
\(\sqrt{16-125/529}>3.97\);
also \(\sqrt k/\alpha<.308\).
Thus \(G>3.662\).
The two early-reference bounds are
\[
 .9104+2(3.662-2.886)/7-5/49>1.0300,\qquad
 .9300+2(3.662-3.034)/7-4/49>1.0277 .
\]
Every non-left-helper possibility
at fifteen/sixteen is excluded.

Those two remaining left helpers
cannot both be weak.
Their actual mutual satisfies
\[
 c_{15}\cdot c_{16}\ge
 1/2-ab+(a/b+b/a)/4\ge1-ab>-11/25>-1/2 ,
\]
so distinct actual walls are
orthogonal by the native gap.
After retaining the left
\(A_{14}\) and ENTIRE right
\(D_3\), their common helper
cost is \(14/30\), and their
two-cap Schur matrix
\(I-(14/30)\mathbf1\mathbf1^T\)
is positive with positive inverse.
Its acute positive projection
is an actual feasible point,
with the original two signed
forks in a positive retained
block.
The already proved mixed-packet
formula, with \(C_{ij}=2\min(i,j)\),
times \(>2\) and own parameters
\(>.8\), gives weighted loss
\(<7/16+1/(4\sqrt{10})<.61\).
This contradicts two weak sites.
Thus fifteen or sixteen is
genuinely OLD strong, under the
assumed weak five.

If sixteen is strong, the
reference through one/six/sixteen
already has
\[
 \|x\|^2>.7600299945+(4-\sqrt6)^2/10
                >1.0004379945 ,
\]
using \(\sqrt6<2.44949\).
Only strong fifteen remains.

Every FINITE weak-five own gives
\(x_1>.63\).
Here \(\overline h=\sqrt{10}-.61<2.553\).
Native five uses the actual finite
\(H_4\), time \(A>2.768\) and
radical \(<.111\), yielding
\[
                     x_1>1.414(2.768-.111-.81(2.553))>.63 .
\]
Native seven uses the actual
negative core \(c_5+g n_5+n_1/2\),
with \(A>3.0688\), radical
\(<.2486,g<.901\), giving
\(x_1>1.414(3.0688-.2486-.901(2.553))>.63\).
Native eight is RETAINED and
uses that same actual core:
\(A>3.1435\), radical \(<.322\),
\(g<.9239\), giving
\(x_1>1.414(3.1435-.322-.9239(2.553))>.6543\).
For every finite label at least
nine, the actual pair projection
gives
\[
                   x_1>\sqrt2\{A-\sqrt{1-g^2}-g\overline h\}.
\]
The brace increases with \(g\):
\(A'>\sqrt{10}>\overline h\),
and the negative radical has
positive derivative.
At \(g=.939\), \(A>3.192\),
\(\sqrt{1-g^2}<.344\), giving
\(x_1>1.414(3.192-.344-.939(2.553))>.6373\).
The elementary native-nine lower
bracket follows by signs of
\(8g^3-6g-1\).
All these feet are future/all-wall
points constructed with original
roots one through six.

The positive reference through
\((1,.63),(6,B_0),(15,B_0+\sqrt{15}-\sqrt6)\)
has energy
\[
 .63^2+(B_0-.63)^2/5+(\sqrt{15}-\sqrt6)^2/9
                         >.7822920845+1.4234^2/9>1.0074 .
\]
Its heights lie below actual
strong fifteen, its slopes
\(.63,.27763,1/(\sqrt{15}+\sqrt6)\)
decrease.
Thus every finite weak-five own
is impossible.

For dotted \(1<g<6/5\), retain
the actual foot BEFORE dropping
the chamber difference
\(r=e_5-e_6\).
The future core
\(c_5+g^{-1}n_5\) gives
\[
 T=P_1+2gP_5-g^{-1}r,\qquad
        x\cdot T\ge K=\sqrt2\{A-\sqrt{1-g^{-2}}\}.
\]
The function
\(A-\sqrt{1-g^{-2}}-g\overline h\)
decreases: \(A'<13/4\), the
radical derivative exceeds one
and \(\overline h>5/2\).
At \(6/5\) it gives the strict
bound
\[
 K-2gf_5>
 \sqrt2\{\sqrt{163/10}-\sqrt{11}/6-(6/5)\overline h\}
                                      >.5965 .
\]
A rational certificate uses
\(\sqrt{163/10}>4.03732\),
\(\sqrt{11}/6<.552771\),
\(\sqrt{10}<3.162278\),
\(\sqrt2>1.41421\), whose lower
product is \(.596535556834\).

Set
\[
 F=1.80474,\quad B=3.44164,\quad X=.59647 .
\]
Positive square bounds give
\(1.80473<f_5<F\), \(f_{15}>B\).
Since \(F-f_5<.00001\), the
actual primitive pressure obeys
\(x\cdot T\ge K>X+2gF\).
Define exact positive rationals
\[
\begin{gathered}
 a=(7B-8F+X)/66=10.25003/66,\qquad
 t=(2F-3X+B)/22=5.26171/22,\\
 y=2t-a=21.32023/66,\qquad
 \mu=X-y=18.04679/66 .
\end{gathered}
\]
Let
\[
 W=(X,y,y,y,t,t,
            \underbrace{a,\ldots,a}_{9},0,0,0,0).
\]
Then \(W\cdot P_5=F\),
\(W\cdot P_{15}=B\), \(W\cdot r=0\).
For the ACTUAL unknown \(g\), put
\[
               \lambda=y-a-2g\mu<0,\qquad
               \rho=a+\mu/g-t>0 .
\]
Indeed \(y-a<.168\),
\(2g\mu>2(.2734)\), and
\(a>.1553,\mu>.2734,t<.2392,1/g>5/6\)
give \(\rho>.1439\).
There is the exact original-
coordinate identity
\[
                         W=\mu T+\lambda P_5+aP_{15}+\rho r .
\]
Its NEGATIVE coefficient is
intentional and legal:
the assumed old weakness is
the upper bound \(S_5\le f_5<F\).
The other signs use the actual
strong fifteen lower bound and
the original strict chamber
inequality \(x\cdot r>0\).
Consequently
\[
 x\cdot W>
 \mu(X+2gF)+\lambda F+aB
       =\mu X+(y-a)F+aB=W^2 .
\]
The last identity follows also
by pairing the displayed coordinate
decomposition with \(W\), using
\(W\cdot T=X+2gF\).
No pressure on a residual was
invented and no independently
optimized energy was added.

Finally
\(X>.5964,y>.3230,t>.2391,a>.1553\),
so
\[
 W^2>.5964^2+3(.3230)^2+2(.2391)^2+9(.1553)^2
                                 =1.00008039>1 .
\]
The one original-vector square
\(\|x\|^2\ge2x\cdot W-W^2>W^2\)
contradicts \(\|x\|<1\).
This proves the head-independent
OLD floor \eqref{app:exact19:eq:Doldfive}.
It asserts neither a \(.58\)
five floor nor centre/fork strength.


\subsection{Native-five \(D_{19}\): the fork bounds actually needed by the centre}
\label{app:exact19:Dfiveforkinputs}

Assume the isolated head has
native five, and put
\[
 a=\varphi,\quad g=a/2,\quad
 c_H=A_Hp-g\omega_1,\quad A_H^2=(a-1)/2 .
\]
Its actual finite \(H_4\) point,
proved in
Section~\ref{app:exact19:longfeet},
gives \(S_4\ge2\sqrt2A_H>1.572\).
We need two fork consequences,
rather than the full outgoing-
fork taxonomy.

First BOTH actual fork edge
cosines are at most ten.
Use the independent old strong-
five floor
\eqref{app:exact19:eq:Doldfive}
and an actual old strong
\(K=11\) or twelve.
Its curved reference \(w_0\)
has first-five mean \(b<.361\),
individual \(P_5\) price \(>.146\),
endpoint price \(>1/7\), and
\(q_{11}>.8472225,q_{12}>.8689525\).
Insert the head ONCE through
\(w_H=w_0+.64(P_4/4-P_5/5)\).
The price is \(.128\) at five.
The orthogonal residual square
is \(1/20\); its effective gain
is \(>1.28(.393-.361)-.4096/20=.02048\).

If an actual fork had \(C>10\),
its original-edge comparison would
give the signed pressure
\[
 x\cdot V>\sqrt{171/11}>3.942,\qquad
              V=P_{18}\pm e_{19},\quad V^2=19 .
\]
The residual \(V-P_K\) vanishes
on the entire head-paid support.
A legal \(1/7\) endpoint payment
gives the one-vector totals
\[
\begin{split}
 .8472225+.02048+2(3.942-2.886)/7-8/49&>1.00615,\\
 .8689525+.02048+2(3.942-3.034)/7-7/49&>1.00600 .
\end{split}
\]
Thus both actual fork cosines
are at most ten, without a
centre floor.

Second no ACTUAL wall whose
spatial part is \(d\omega_f\),
for either fork \(f\), can have
finite own \(d=\cos(\pi/m)\).
This statement does not require
the wall to be first on a fork
edge.
It also applies to a centre
wall of own zero with one
finite fork contact.

For such a wall,
\(A^2=(19/2)d^2-1\), and
\[
                     M=A_H\sqrt{(19/2)d^2-1}-(a/2)d .
\]
The head and wall are distinct
by their different old supports.
For \(1/2\le d\le1\), \(M\)
is positive and increasing:
at \(1/2\) its squared-product
difference is \((5a-6)/8>0\);
its derivative is positive since
\(A_H\sqrt{19/2}>a/2\).
Furthermore
\[
                    M<M(1)<.556(2.916)-.809
                              =.812296<.82<\cos(\pi/6).
\]
The mutual is therefore FINITE,
with exactly one of
\(M=1/2,\sqrt2/2,a/2\).
Its strict scalar range excluded
orthogonality, parallel and
dotted mutuals before algebra.

Only the elementary finite-angle
fact is used next:
\(t=2d=\zeta_{2m}+\zeta_{2m}^{-1}\)
is an algebraic integer for
EVERY Coxeter polytope.
The fields below are fields of
known real angle constants.
They are NOT identified with
a defining field of the reflection
group, and no dotted product
or invariant cycle is declared
integral.

Put \(L=\mathbb Q(a)\),
\(\mathcal O_L=\mathbb Z[a]\),
\(B=19-10a\), \(\mathrm N_L(B)=71\).
Squaring the positive mutual
equation gives exactly
\[
\begin{array}{c|c|c}
 M&\text{quadratic for integral }t&\text{discriminant}\\ \hline
 1/2&Bt^2-2t+2a-6&12(45-26a)\\
 \sqrt2/2&Bt^2-2\sqrt2t-4&8(39-20a)\\
 a/2&Bt^2-2at+2a-8&4(173-97a).
\end{array}
\]
The first and last discriminant
norms are \(144(179)\), \(16(3739)\).
They are nonsquares, since
\(13^2<179<14^2\),
\(61^2<3739<62^2\).
Thus the quadratics are
irreducible over \(L\).
The relative norm of integral
\(t\) must be integral in \(L\),
yet its rational absolute norm
would be \(20/71\), respectively
\(44/71\), impossible.

For the middle row work in the
genuine quadratic extension
\(L_2=L(\sqrt2)\).
An element \(D\in L\) can be
square there only if \(D\) or
\(D/2\) is square in \(L\):
expand \((r+s\sqrt2)^2\), whose
mixed coefficient forces \(rs=0\).
For \(D=8(39-20a)\), the two
norms are \(64(341),16(341)\),
both nonsquare as
\(18^2<341<19^2\).
This quadratic is irreducible
over \(L_2\).
Its integral root would have
relative norm \(-4/B\), rational
absolute norm \(256/71^2\),
again impossible.
Every finite own is excluded.

Own zero has negative time;
parallel own has the nonzero
feasible future null \(c+n_f\).
Consequently an isolated fork
wall, when encountered, has
dotted own.
No integrality of that dotted
own is used.
These are the only outgoing-
fork facts needed below.

\subsection{Native-five \(D_{19}\): one head-paid centre reference}
\label{app:exact19:Dfivereference}

The complete new bulk through
fourteen is already available.
The actual native-five mutual
calculation uses only the first
seven long roots, so the proof
in Section~\ref{app:exact19:Bhead}
applies here.
Explicitly, for a helper with
\(r<1.2\) and \(i\le7\),
\[
 A_H^2A_i^2-[a(r+1/2)]^2
 =[i(a-1)-a^2]r^2-2r-a/2
                                <.0568r-.809<0 .
\]
It would have positive actual
mutual with the head.
The remaining isolated dotted
five is excluded by the same
actual native-gap calculation:
weakness forces
\(1<r<\sqrt{5/2}\), and its
positive increasing head
magnitude is
\(<.556(4.899)-1.618(1.581)
=.165786<1/2\).
Thus five/six/seven are truly
\(.58\)-strong.

Choose an actual strong
\(K=13\) or fourteen.
All internal weak helper holes
cost \(<.000158\), using the
earliest permitted set \(8,10,12\).
The rational certificate
\eqref{app:exact19:eq:BE17},
restricted to its initial
increments, gives
\[
 E_{13}>.9056051,\quad E_{14}>.9241299,\quad
 f_{13}<3.1955,\quad f_{14}<3.3316,\quad b=f_5/5<.3652 .
\]
The reference \(w_0\) has \(P_5\)
price \(>.151\) and endpoint
price \(>1/8\).
Insert the head once in
\(w_H=w_0+.14(P_4-.8P_5)\).
Its debit is \(.112\) at five,
orthogonal residual square \(.8\),
and effective gain
\[
                   .28(1.572-.8(1.826))-.0196(.8)=.015456 .
\]
Hence the actual paid references
have
\[
                        L_{13}>.9209031,\qquad L_{14}>.9394279 .
\]
Their first four coordinates are
\(<.3932\), their sums remain
\(f_K\), and the individual end
price stays \(>1/8\).

For an actual pressure
\(x\cdot V\ge G\), \(V=1\)
through the chosen \(K\),
\(V^2\le R\), the residual
\(V-P_K\) vanishes on the
ENTIRE head-paid support.
One paid square gives
\[
 \|x\|^2>L_K+\frac{16G-R-16f_K+K}{64}.
\tag{D19fivepay}\label{app:exact19:eq:Dfivepay}
\]
Therefore the uniform sufficient
contradiction is
\[
                              16G-R>43.19021 .
\tag{D19fivebudget}\label{app:exact19:eq:Dfivebudget}
\]
The exact thresholds for the
two \(K\)'s are
\(43.1902016,43.1822144\).
An ACTUAL \(.50\)-strong centre
would have \(S_{17}>3.7695\);
\(V=P_{17},R=17\) already
violates this budget.
The following centre classification
may thus assume \(.50\) weakness
as a CONTRARY hypothesis.
No centre floor is being supplied
in advance.


\subsection{Native-five \(D_{19}\): every isolated centre fails}
\label{app:exact19:Dfiveisolatedcentre}

Own zero has negative time.
Own three has the actual four-
leaf affine star centred at
seventeen; own four the weighted
star with both forks; six the
\(6,3\) path through sixteen;
parallel the feasible positive
null \(c+n_{17}\).
Each null is future through its
positive cap coefficient and
nonpositive against all omitted
original walls.

At native five use the LEFT
finite \(H_4\) on the actual
cap/seventeen/sixteen/fifteen.
Its positive correction has
residual \(\delta=(5-3a)/8\);
the actual projected point,
normalized by \(\alpha=\sqrt2g\),
gives \(V=1\) through fourteen,
\(R<17\), and
\[
 G>\sqrt{17-1/(2g^2)}-\sqrt\delta/\alpha
                                      >4.029-.119>3.91 .
\]
This contradicts
\eqref{app:exact19:eq:Dfivebudget}.

At native seven/eight use
\(u=r^{-2}\) and the actual core
\(c+r^{-1}n_{17}+(2r)^{-1}n_{16}\).
Its cap/sixteen products vanish,
seventeen product is negative,
and its square is
\(-(1-3/(4r^2))\).
It is future and every other
original-wall product is
nonpositive.
Its normalized displacement is
\[
 \begin{gathered}
 V=P_{15}+(1-u/4)(e_{16}+e_{17})+(u/2)e_{18},
                 \quad R=17-u+3u^2/8<17,\\
 G=\sqrt{17-u/2}-\sqrt{u/2-3u^2/8}>4.047-.27=3.777 .
 \end{gathered}
\]
Indeed \(1.168<u<100/81\);
the second squared expression
decreases and at \(1.168\) is
\(.072416<.27^2\).
Thus \(16G-R>43.432\).

For EVERY finite own at least
nine, \(r>.939\) and
\(r^{-2}<1.1352\).
The actual finite pair with
seventeen has normalized
\(V=P_{16}+(e_{17}+e_{18})/2\),
\(R=16.5\).
Its TRUE projected-point demand
is larger than
\[
 \sqrt{17-1/(2r^2)}-\sqrt{1/(2r^2)-1/2}
                                      >4.053-.26=3.793 .
\]
The native brackets used here
are only ordinary real-angle
square/cubic bounds.
Every finite case violates the
common budget.

For dotted \(r>1\),
\(u=r^{-2}\in(0,1)\), the actual
core \(c+r^{-1}n_{17}\) gives
\[
\begin{gathered}
 V=P_{16}+(1-u/2)e_{17}+(u/2)e_{18},\quad
 R=17-u+u^2/2,\\
 G=\sqrt{17-u/2}-\sqrt{u(1-u)/2}.
\end{gathered}
\]
The complete continuous estimate
is \(16G-R>43.19278\).
To prove it, use the concave
root chord, \(\sqrt{17}>4.123105\)
and
\(8/(\sqrt{17}+\sqrt{33/2})<.9775\),
obtaining
\[
 16G-R>48.96968+.0225u-u^2/2
                                  -8\sqrt2\sqrt{u(1-u)} .
\]
For \(u\le.45\), \(u^2\le.45u\)
and monotonicity give the lower
bound
\(48.96968-.091125-\sqrt{31.68}
>43.249555\).
For \(.45\le u\le.55\),
use \(u^2\le u-.2475\) and
the circle maximum
\[
 Au+B\sqrt{u(1-u)}
                         \le(A+\sqrt{A^2+B^2})/2 .
\]
With \(A=.4775,B=8\sqrt2\) and
\(\sqrt{128.22800625}<11.3238\),
the lower bound is
\(48.96968+.12375-5.90065=43.19278\).
For \(u\ge.55\), use
\(u^2\le1.55u-.55\),
\(A=.7525\) and
\(\sqrt{128.56625625}<11.33872\),
giving \(43.19907\).
All radicals have positive
rational square checks.
Thus the ENTIRE dotted isolated
centre case is excluded.

\subsection{Native-five \(D_{19}\): the complete finite-white centre list}
\label{app:exact19:Dfivecentrelist}

Deleting seventeen leaves the
old \(A_{16}\) and both isolated
fork \(A_1\)'s.
The complete finite positive
white list is:
\[
\begin{array}{l}
 \text{isolation};\\
 \text{left-only }3@1,2,15,16\text{ or }4@1,16;\\
 \text{one fork with any finite native contact};\\
 \text{two forks }(3,3),(3,4),(3,5),\text{ either orientation};\\
 \text{left endpoint-three }j=1,16\text{ and one fork-three/four,
 or both fork-threes}.
\end{array}
\]
A second left contact makes a
cycle; a next-to-end attachment
plus another fork has excessive
exceptional arms; no exceptional
finite diagram has the required
full rank.
The whole-forest hypothesis
leaves no foreign component.

For direct \(.50\) tests use
\(T_{.50}(\sqrt{34})<265/138<1.921\).
For left-only three,
\[
          \rho=(\gamma+j/34)/
                          \sqrt{1-j(17-j)/34}.
\]
At \(j=2\), own zero has zero
time, three an affine cycle;
every own at least four is
directly strong.
At \(j=15\), own zero has the
actual positive affine kernel
\((1,1,2,2,2,1,1)\) on
\(c,14,15,16,17,18,19\);
every positive own is strong.
At \(j=16\), zero has kernel
\((1,1,2,2,1,1)\) on
\(c,15,16,17,18,19\), three
its triangle, four the weighted
star using both forks, six its
\(6,3\) path.
Finite owns at least nine and
dotted owns are strong;
five/seven/eight need the actual
point below.
At \(j=1\), zero has negative
time, three the actual long
cycle, four the weighted star
using sixteen and one fork,
six its \(6,3\) path.
Its remaining far-helper owns
are retained for a separate
calculation.

For left-only four,
\[
       \rho=(\gamma+j/(17\sqrt2))/
                               \sqrt{1-j+j^2/17}.
\]
Far own zero has zero time;
all other far cases and even
near own zero are directly
strong.

For one fork \(\beta\),
\[
                       \rho=(\gamma+\beta/2)/\sqrt{1-\beta^2}.
\]
Own zero is the actual finite-
fork wall already excluded by
Section~\ref{app:exact19:Dfiveforkinputs},
whose proof did not require it
to be first on a fork edge.
At own three, deleting the
contacted fork gives a positive
finite \(D_{19}\) replacement.
Actual positive inverse projection
makes it the reciprocal fork
endpoint, of cosine
\[
                            C_f=(17+19\beta)/2\ge53/4>10 .
\]
This contradicts the actual fork
bound.
The centre and fork opposites
are allowed to coincide; no
false distinctness was assumed.
Own four has the affine weighted
star using sixteen and the OTHER
fork; six its \(6,3\) path;
parallel its positive future
null.
For \(\gamma\ge\varphi/2\) and
\(\beta\ge\cos(\pi/5)\), the
ratio is \(>2.06\).
Thus only \(\beta=1/2,1/\sqrt2\)
needs further point estimates.

For two forks \(\beta,\eta\),
\[
             \rho=\frac{\gamma+(\beta+\eta)/2}
                                   {\sqrt{1-\beta^2-\eta^2}} .
\]
Own zero of \((3,3)\) is the
actual all-three square, three
an affine triangle, six/parallel
their actual null obstructions.
Own four/five will be treated
below; finite at least seven
and dotted own are strong.
Own zero of \((3,4)\) has an
actual finite \(B_{19}\)
reciprocal endpoint at the
three-fork, with
\(C_f=17/2+(19/4)\sqrt2>10\).
Positive owns there, and even
own zero in \((3,5)\), are
directly strong.

For left endpoint-three at
\(j=1,16\) with fork magnitudes
\(\beta,\eta\), the exact data
are
\[
 \delta=9/17-\beta^2-\eta^2,\qquad
               \rho=\frac{\gamma+j/34+(\beta+\eta)/2}{\sqrt\delta}.
\]
One fork-three at own zero
has the actual long/four-cycle.
Near positive own is strong.
Far own three has its fork
triangle, four the weighted
star using sixteen and the
OTHER fork; larger owns are
strong.
One fork-four or two fork-threes
has residual \(1/34\), strong
even at own zero.
Both signed fork orientations
and every mixed own-zero role
are now included.
Every positive affine/null
combination is future through
the cap and nonpositive against
ALL omitted original walls.


\subsection{Native-five \(D_{19}\): every remaining centre point exception}
\label{app:exact19:Dfivecentrepoints}

For near-left three and own
five/seven/eight,
\(v=\gamma\omega_{17}+\omega_{16}/2\),
\(A^2=34\gamma^2+32\gamma+7\).
Use the ACTUAL finite pair
\(c,n_{17}\), with residual
\(\delta=1-\gamma^2>0\) and
projected-point pressure
\(H=AB/(B+\sqrt\delta)\),
\(B^2=A^2+\delta\).
Its normalized displacement is
\[
 V=P_{16}+\frac{\gamma}{2\gamma+1}(e_{17}+e_{18}),
              \quad\alpha=\sqrt2(\gamma+1/2),\quad R<16.23 .
\]
At native five,
\(A^2=(49a+31)/2\),
\(B^2=(97a+65)/4\).
Positive-square bounds
\(A>7.425,B>7.448,\sqrt\delta<.588,\alpha<1.852\)
give
\[
                      G>7.425(7.448)/(8.036(1.852))>3.715 .
\]
The sharper \(R=17-a/2<16.2\)
gives \(16G-R>43.24\).
The weaker \(A-\sqrt\delta\)
pressure would be insufficient.
At native seven/eight,
\(A/\alpha>4.019\),
\(\sqrt\delta/\alpha<.221\);
hence \(G>3.798\), which closes
these rows as well.

For one/two fork helpers write
\[
 v=\gamma\omega_{17}+\beta\omega_f+\eta\omega_h,\quad
 \alpha=\sqrt2(\gamma+(\beta+\eta)/2),\quad
                    A^2=17\alpha^2-(1-\beta^2-\eta^2).
\]
At own five, the LEFT actual
finite \(H_4\) has correction
\(g(3n_{17}/2+n_{16}+n_{15}/2)\),
residual \((5-3a)/8\).
For a single helper
\(\beta=1/2,1/\sqrt2\), or
two helpers \(\beta=\eta=1/2\),
its actual projected point has
\(V=1\) through fourteen,
\(R<17\), and
\[
                         G>4.082-.091>3.99 .
\]
Use \(\alpha>1.497\),
\(A/\alpha>4.082\),
\(\sqrt{(5-3a)/8}/\alpha<.091\).
The contacted fork walls remain
outside the selected finite
piece, with their actual
products nonpositive.
Every own-five helper exception
is excluded.

For one fork-four and own
finite at least seven or dotted,
contrary \(.50\) weakness forces
\(.9<\gamma<1.01\).
The actual positive core
\[
                         c+\gamma^{-1}n_{17}+(2\gamma)^{-1}n_{16}
\]
has zero cap/sixteen products,
negative seventeen product,
square \(-k\),
\(k=1-3/(4\gamma^2)>0\).
It is future and every omitted
wall product is nonpositive.
With
\(\alpha=\sqrt2(\gamma+\sqrt2/4)\),
\(V\) is flat through fifteen,
\(R<17\), and
\[
              G=\sqrt{17-(1/2)/\alpha^2}-\sqrt k/\alpha
                                      >4.103-.291=3.812 .
\]
The complete seals are
\(\alpha>1.772,\alpha^2>3.141\),
\(k<.265,\sqrt k<.515\).

For one fork-three the entire
contrary range is
\(.9<\gamma<1.42\).
Use the SAME actual core, with
\(\alpha=\sqrt2(\gamma+1/4)\),
\(A^2=17\alpha^2-3/4\).
Its displacement is flat through
fifteen and
\[
 G=\sqrt{17-(3/4)/\alpha^2}-\sqrt k/\alpha>3.756,\qquad
                    R=17-(3/4+k)/\alpha^2<16.8 .
\]
The continuous proof is explicit.
For \(\gamma\le1.1\), the first
root is \(>4.088\), the second
\(<.324\).
The latter squared ratio increases
up to \(1.1\), its derivative
having sign
\(-2\gamma^3+3\gamma+3/8>0\).
For \(\gamma\ge1.1\), the first
root is \(>4.098\).
Uniformly \(\sqrt k/\alpha<.342\):
the inequality \(k/\alpha^2<7/60\)
is equivalent to
\[
                    112\gamma^4+56\gamma^3-473\gamma^2+360>0 .
\]
At \(t=\gamma-5/4\) this becomes
\(15/4-45t+787t^2+616t^3+112t^4\).
Because \(t\ge-.35\), it is at
least
\(15/4-45t+(2857/5)t^2>0\),
with negative discriminant.
For the norm,
\(1.75-.75/\gamma^2-.4(\gamma+.25)^2\)
is concave and positive at
both \(.9,1.42\).
Thus \(16G-R>43.296\) throughout,
including all dotted parameters.

The two-fork-three/own-four row
is a separate actual point.
Let \(s=\sqrt2,\gamma=s/2\) and
\[
 \xi=c+(4-2s)n_{17}+(2-s)n_{16}
                                 +(3-2s)(n_{18}+n_{19}).
\]
All coefficients are positive.
Its cap/seventeen/sixteen products
vanish, each fork product is
\(1/2-s<0\), and
\(\xi^2=11-8s<0\).
It is future and feasible
against every original wall.
For \(\alpha=1+s/2\),
\[
 V=P_{15}+(5-3s)(e_{16}+e_{17}+e_{18}),\quad
 R=144-90s<16.721,
\]
\[
 G=\sqrt{14+2s}-\sqrt{92s-130}
                                      >4.102-.329=3.773 .
\]
Hence \(16G-R>43.647\).
All non-far-helper centre cases
have now been excluded at their
actual changed threshold.

\subsection{Native-five \(D_{19}\): the only centre left is a dotted far helper}
\label{app:exact19:Dfivefarcentre}

For the far-left-three cap,
\(v=\gamma\omega_{17}+\omega_1/2\),
\(A^2=34\gamma^2+2\gamma-1/2\).
Own zero/three/four/six/parallel
already has its actual time,
affine or positive-null
obstruction.
The remaining finite-own
pressure table, normalized by
\(\sqrt2\gamma\), is
\[
\begin{array}{c|c|c|c}
 \text{own}&V=P_{16}+Y&
        \text{actual point/core}&G>\\ \hline
 5&a_3P_3+(e_{17}+e_{18})/2&
       \text{finite }H_4,\ \delta=2/3-\gamma^2&4.1158\\
 7,8&a_2P_2+(e_{17}+e_{18})/2&
       c+\gamma n_{17}+n_1/2&3.963\\
 \text{finite }\ge9&a_1P_1+(e_{17}+e_{18})/2&
       \text{finite pair }c,n_{17}&3.965 .
\end{array}
\]
Here
\(a_3=1/(6\gamma)<5/24\),
\(a_2=1/(4\gamma)<5/18\),
\(a_1=1/(2\gamma)<.533\).
All selected roots are actual.
Their positive corrections give
future points retaining every
unknown wall and both original
forks; the finite rows use
projected POINTS.

All finite rows satisfy
\(A/(\sqrt2\gamma)\ge\sqrt{17.75}>4.213\).
The first two normalized radical
subtractions are \(<.097125,.25\).
For finite at least nine,
\(\gamma\ge.95\) gives the crude
pair lower bound
\(>4.213-.233=3.98\).
For \(.939<\gamma\le.95\),
the crude lower bound is
\(>4.213-.26\), and the TRUE
finite pressure improves
\(A-\sqrt\delta\) by
\(\delta/(B+\sqrt\delta)\).
Since
\(\delta\ge.0975\),
\(\sqrt2\gamma<1.344\),
\(B<5.674,\sqrt\delta<.344\),
its normalized improvement is
\(>.012\).
Thus \(G>3.965\) in the whole
finite-nine-and-higher range.
Native eight was retained
geometrically, not removed by
a field condition.

The head overlap is NONZERO.
For \(Y=a_lP_l+(e_{17}+e_{18})/2\),
\[
 \begin{gathered}
 V-P_K=P_{16}-P_K+Y,\qquad
 w_H\cdot(V-P_K)=a_l(w_H\cdot P_l),\\
       \|V-P_K\|^2=16-K+l a_l^2+1/2 .
 \end{gathered}
\]
With the same legal \(1/8\)
endpoint price, the sufficient
barrier is therefore
\[
             16\{G-a_l(w_H\cdot P_l)\}
                        -(16+l a_l^2+1/2)>43.19021 .
\]
The true head-paid first coordinates
give \(w_H\cdot P_l<.3932l\).
The three rows have respective
lower bounds
\(45.29059,43.25856,43.30270\).
All violate the norm budget.
No aggregate lower bound was
subtracted to manufacture
a residual pressure.

For dotted \(1<\gamma\le10/9\),
let \(a_0=1/(2\gamma)\in[.45,.5)\).
The actual core
\(c+\gamma^{-1}n_{17}\) gives
\[
 \begin{gathered}
 V=P_{16}+a_0P_1+(1-2a_0^2)e_{17}+2a_0^2e_{18},\\
 G=\sqrt{17+2a_0-a_0^2}-a_0\sqrt{2-8a_0^2}.
 \end{gathered}
\]
The residual squared norm is
\(16-K+1-3a_0^2+8a_0^4\le16-K+.75\).
The function \(G-.3932a_0\)
strictly increases, since its
derivative is
\[
 \frac{1-a_0}{\sqrt{17+2a_0-a_0^2}}
       +\frac{16a_0^2-2}{\sqrt{2-8a_0^2}}-.3932>0 .
\]
At \(a_0=.45\),
\(\sqrt{17.6975}>4.2068\),
\(\sqrt{.38}<.6165\), giving
\(G-.3932a_0>3.752435\).
The paid barrier is
\(>16(3.752435)-16.75=43.28896\),
a contradiction on the full
small-dotted interval.

For \(\gamma\ge11/8\), its exact
edge ratio
\((\gamma+1/34)/\sqrt{9/17}>1.928\)
already gives the contradictory
\(.50\)-strong centre floor.
The ENTIRE remaining actual
centre has consequently the
profile
\[
 c=Ap-\gamma\omega_{17}-\omega_1/2,\qquad
 A^2=34\gamma^2+2\gamma-1/2,\qquad
                              10/9<\gamma<11/8 ,
\tag{D19fivecentre}\label{app:exact19:eq:Dfivecentre}
\]
with dotted own.
No outgoing fork equality,
isolation or symmetry was
assumed.


\subsection{Native-five \(D_{19}\): classify short fork edges at the actual adjacent vertex}
\label{app:exact19:Dfiveq}

Replace old seventeen by its
ACTUAL opposite
\eqref{app:exact19:eq:Dfivecentre}.
The other endpoint \(q\) of that
original edge has finite forest
\[
                   c\;--3--\;n_1\;--3--\cdots--3--\;n_{16},
                     \qquad n_f,\ n_h ,
\]
an \(A_{17}\) and two orthogonal
\(A_1\)'s.
This is the actual neighbouring
vertex, not a formal completion
or a vertex assumed farthest.
Set
\[
 B=17\gamma+1/2,\quad
 C=(34\gamma+1)/(3\sqrt2),\quad
 s=A\sqrt{17}/3,\quad
                   q=Cp-sP_{17}/\sqrt{17},\quad C^2-s^2=1 .
\]
Both forks are unchanged unit
normals tangent at \(q\).
Put \(\sigma=s/\sqrt{34}=A/(3\sqrt2)\).
The centre interval gives
\[
                 A^2<2129/32,\qquad \sigma<77/40=1.925,
                                  \qquad B>349/18 .
\]

We need the products of the new
path endpoint duals with the OLD
actual wall \(n_{17}\).
Let
\[
 w=(e_1-P_{17}/17)/\sqrt2,\quad
 \eta=c+w,\quad\eta^2=9/17,\quad
 d_{16}=\sqrt2(P_{16}-16P_{17}/17).
\]
These are tangent to \(q\).
The new path duals are
\(\Omega_c=(17/9)\eta\),
\(\Omega_{16}=d_{16}+\Omega_c/17\):
\(\eta\) is orthogonal to old
one through sixteen, and
\(d_{16}\cdot c=-1/17\).
Directly,
\[
                 \Omega_c\cdot n_{17}=-B/9,\qquad
                 \Omega_{16}\cdot n_{17}=-(144+B)/153 .
\tag{D19qduals}\label{app:exact19:eq:Dqduals}
\]
These are actual Lorentz tangent
identities; no projected mirror
has been presumed Coxeter.

Call a q-fork edge short if its
actual length satisfies
\(\cosh L<\sqrt{3/2}\).
Let \(r\) be its actual opposite,
with positive own magnitude
\(\beta\), full finite-white
cost \(Q<1\), and time \(a_r>0\).
The unit \(A_1\) inverse gives
\[
 a_r^2=\beta^2+Q-1,\qquad
               1-Q<\beta^2<(3/2)(1-Q).
\tag{D19qshort}\label{app:exact19:eq:Dqshort}
\]
Own zero cannot satisfy this.
Parallel own has the actual
positive future null \(r+n_f\)
and is impossible.
All finite and dotted owns
remain in the displayed strip.

After releasing that \(A_1\),
the endpoint retains the ENTIRE
new \(A_{17}\) and the other
\(A_1\).
Its complete positive finite
white list and costs are
\[
\begin{array}{c|c}
 \text{contact outside the released }A_1&Q\\ \hline
 \text{none}&0\\
 \text{only other }A_1,\text{ finite label }m\ge3&\cos^2(\pi/m)\\
 A_{17}\text{ end-three}&17/36\\
 A_{17}\text{ next-to-end-three}&8/9\\
 A_{17}\text{ end-four}&17/18\\
 A_{17}\text{ end-three plus other-three}&13/18\\
 A_{17}\text{ end-three plus other-four}&35/36 .
\end{array}
\]
A second path contact is a cycle.
Its rank-eighteen one-node
extensions are \(A,B,D\), with
no exceptional position.
Joining both retained pieces
allows only end-three followed
by other-three/four, giving
\(A_{19},B_{19}\).
A next-to-end attachment plus
another edge has excessive
exceptional arms, and an end-four
followed by another node has
a nonterminal double.
Thus every finite white type
is included.

Applying the exact short strip
leaves:
isolated dotted own
\(1<\beta^2<3/2\);
path end-three with own five/
six; or other-\(A_1\)-only
own-five/white-four,
own-four/white-five,
own-three/white-seven, or
white-three with finite own
at least seven or dotted,
\(3/4<\beta^2<9/8\).
Here are the complete label
checks for the other-\(A_1\)
case: white three/four/five/
six/seven/at-least-eight gives
respectively the own-square
strips
\[
 (.75,1.125),\quad(.5,.75),\quad
 ((3-a)/4,3(3-a)/8),\quad(.25,.375),
\]
a subinterval containing only
native own-three at white seven,
and upper bound \(<.25\) at
white at least eight.
The other cost rows are empty:
\(Q\ge8/9\) has own-square
upper bound \(<1/6\), and
\(Q=13/18\) has only possible
native three below its positive-
time lower bound.
The mixed \(35/36\) row is
smaller still.
This is local finite classification
at the ACTUAL \(q\), with no
farthest hypothesis there.

Every short opposite \(r\) is
distinct from OLD \(n_{17}\):
that wall has dotted contact
\(-\gamma\) with new \(c\), which
cannot occur in its retained
positive finite white endpoint.
Actual acuteness/native products
therefore apply to \(r,n_{17}\).

For path end-three, \(j=c\) or
sixteen, write
\(r=a_rq-\beta n_f-\Omega_j/2\).
Own is five/six, so
\(\beta>.809,a_r<.472\).
Using \eqref{app:exact19:eq:Dqduals},
\[
 r\cdot n_{17}
     =-a_r\sigma+\beta/2-(\Omega_j\cdot n_{17})/2
          >-.472(1.925)+.4045+.533=.0289>0 .
\]
For endpoint \(c\) the last
contribution exceeds one; for
endpoint sixteen it exceeds
\(.533\).
Both are impossible.

For a cap supported on the two
q-\(A_1\)'s,
\[
 r=a_rq-\beta n_f-\eta n_h,\qquad
           r\cdot n_{17}=-a_r\sigma+(\beta+\eta)/2 .
\]
Own-three/white-seven has
\(a_r<.249\), half spatial sum
\(>.7\), so the product is
\(>.7-.249(1.925)>.22\).
Own-five/white-four and its
reverse have
\(a_r^2=(a-1)/4<.3931^2\),
\((\beta+\eta)/2>.75805\).
The product is
\(>.75805-.3931(1.925)=.0013325>0\).
These are strict actual-wall
acuteness contradictions.

Only two short types remain:
\[
\begin{array}{c|c}
 I&r=a_rq-\beta n_f,\quad1<\beta^2<3/2\\
 J&r=a_rq-\beta n_f-n_h/2,\quad
      \text{finite own }\ge7\text{ or dotted},\
      .9<\beta<1.061,\quad a_r^2=\beta^2-3/4<3/8 .
\end{array}
\]
For type \(J\),
\(\rho=-r\cdot n_{17}
=a_r\sigma-(\beta+1/2)/2\ge0\).
But
\(\rho<.613(1.925)-.7=.480025<1/2\).
The ACTUAL native gap forces
\(\rho=0\).
Hence this wall has zero products
with ALL old one through seventeen;
its exact OLD expansion is
\[
 r=H_\beta p-\beta\omega_f-\omega_h/2,\qquad
                 H_\beta^2=(19/2)\beta^2+(17/2)\beta+11/8 .
\tag{D19qJ}\label{app:exact19:eq:DqJ}
\]
It need not be first on an old
fork edge; its old dual expansion
follows just from the actual
products and unit norm.

\subsection{Native-five \(D_{19}\): one q-fork edge is nonshort}
\label{app:exact19:Dfiveqnonshort}

Both q-fork edges cannot be short.
The surviving types cover both
same-wall and distinct-wall
possibilities.
An \(I\) wall has zero other
contact and cannot close the
other edge.
A \(J\) wall has own \(>.9\)
and other contact \(1/2\);
in the reverse direction it
does not have either surviving
short type.
Thus two short opposites in
the remaining list are distinct.

Two type-\(I\) walls have product
\(-a_ra_t\in(-1/2,0)\), because
both time squares are \(<1/2\).
This violates the native gap.
For \(I\) at one fork and \(J\)
at the other, the q-product is
\[
             -a_Ia_J+\beta_I/2
                     >1/2-\sqrt3/4>0 ,
\]
using \(a_I^2<1/2,a_J^2<3/8\)
and \(\beta_I>1\).

For two type-\(J\) walls, use
their OLD expansions
\eqref{app:exact19:eq:DqJ}.
Their negative mutual is
\[
 M(b,d)=H_bH_d-(17/2)bd-(19/4)(b+d)-17/8 .
\]
This function increases in both
independent parameters \(b,d>.9\).
Indeed \(H_b'>\sqrt{19/2}\), and
\[
                 (19/2)H_d^2-[(17/2)d+19/4]^2
                                      =18d^2-19/2>0 .
\]
Since BOTH actual parameters are
\(<1.061\), monotonicity gives
\[
 M(b,d)<M(1.061,1.061)
           =1.061^2-1.061-3/4=-.685279<0 .
\]
Their actual product would be
positive.
Equality of the two actual
parameters was not assumed;
the equal argument is only an
upper bound after monotonicity.

At least one actual q-fork edge
therefore has
\[
                              D=\cosh L\ge\sqrt{3/2}.
\tag{D19qnonshort}\label{app:exact19:eq:Dqnonshort}
\]
No ordinary \(A_1\) energy theorem
or farthest-forest classification
was applied at \(q\).

\subsection{Native-five \(D_{19}\): compare that endpoint with the ORIGINAL farthest vertex}
\label{app:exact19:Dfiveqfinish}

For the nonshort q-fork edge,
let its ACTUAL other endpoint be
\(r'=Dq-\sqrt{D^2-1}\,n_f\),
in whichever signed orientation
has
\eqref{app:exact19:eq:Dqnonshort}.
Put
\[
                \tau=(-z\cdot q)/E
                        =C-(s/\sqrt{17})S_{17}.
\]
The ORIGINAL farthest property
gives \(0<\tau\le1\) and
\[
                       \tau D-\sqrt{D^2-1}\,x\cdot n_f\le1 .
\]
We retain \(\tau\); it is not
replaced by one.
For fixed \(0<\tau\le1\),
\((D\tau-1)/\sqrt{D^2-1}\)
increases with \(D\):
its derivative is
\((D-\tau)/(D^2-1)^{3/2}>0\).
The lower endpoint
\(D_0=\sqrt{3/2}\) is therefore
legal, giving
\[
                               x_{18}\pm x_{19}\ge\sqrt6\,\tau-2 .
\]
Substitution yields the actual
SAME-vector pressure
\[
\begin{gathered}
 x\cdot V\ge G,\qquad
 V=P_{17}+t(e_{18}\pm e_{19}),\quad
 t=\sqrt{17}/(\sqrt6\,s),\\
                   G=\sqrt{17+6t^2}-2t,\qquad R=17+2t^2 .
\end{gathered}
\]
Every endpoint used is actual
and hence satisfies every
original and unknown wall.
There is no formal replacement
or independent energy at \(q\).

Since \(s^2=17A^2/9\) and
\(\gamma>10/9\),
\[
                  t^2=3/(2A^2)<243/7079<9/256,\qquad0<t<3/16 .
\]
The cleared comparison is
\(243(256)=62208<63711=9(7079)\).
The function
\(\sqrt{17+6t^2}-2t\)
decreases on this interval.
At its upper endpoint,
\[
 G>\sqrt{17+27/128}-3/8>4.1486-.375=3.7736,\qquad
 R<17+9/128=17.0703125 .
\]
The positive square check is
\(4.1486^2=17.21088196<17+27/128\).
Thus
\[
                              16G-R>43.3072875>43.19021 .
\]
The displacement is one through
the actual chosen \(K=13,14\).
Its residual is orthogonal
to the head-paid reference,
and its \(1/8\) individual
endpoint price is still available.
Equation~\eqref{app:exact19:eq:Dfivepay}
gives
\[
                 \|x\|^2>1+(43.3072875-43.19021)/64>1.0018 ,
\]
contradiction.

\begin{proposition}\label{app:exact19:prop:Dfive}
A whole \(D_{19}\) forest at the
original farthest vertex cannot
have an isolated native-five
long head.
\end{proposition}
\begin{proof}
The complete centre reduction
leaves exactly
\eqref{app:exact19:eq:Dfivecentre}.
At its actual adjacent vertex
the full local short-edge
classification and old-wall
native gap force one nonshort
fork edge.
Its actual endpoint violates
the single ORIGINAL-vector
head-paid norm budget, with
the exact distance slack retained.
No equality or classification
of the outgoing fork opposites
was required.
\end{proof}

\subsection{A reusable centre test for the remaining finite heads}
\label{app:exact19:Dfinitecentretest}

The centre geometry in the preceding argument is useful beyond the
native-five head. We separate it from the early payment that changes
with the head. Throughout this section the forest is the SAME whole
\(D_{19}\), with its signed last coordinate and actual centre seventeen.
No field restriction is used.

Suppose an actual strong \(K\in\{13,14\}\) has supplied a legally funded
vector \(w\), using only initial coordinates one through six to modify a
curved reference. Write \(L\) for its effective energy: if its legally
available pressure is \(x\cdot w\ge d\), then \(L=2d-\|w\|^2\).
In particular \(L\) need not equal \(\|w\|^2\).
Assume its INDIVIDUAL \(P_K\) coefficient is \(>1/8\), its demand there
is \(f_K=\sqrt K-.58/\sqrt2\), and all initial changes have zero
coordinate sum. Consequently \(w\cdot P_K=f_K\).

For any actual pressure \(x\cdot V\ge G\), put
\[
 B=w\cdot(V-P_K),\qquad R_*=K+\|V-P_K\|^2 .
\]
Consume \(1/8\) of the original \(P_K\) coefficient and use the ONE
vector \(w+(V-P_K)/8\). Its pressure follows from the remaining positive
anchor/foot decomposition, rather than subtraction of unrelated lower
bounds. Expanding its complete square gives
\[
 \|x\|^2\ge
 L+\frac{16(G-B)-R_*-16f_K+K}{64}.
\tag{D19generalpay}\label{app:exact19:eq:Dgeneralpay}
\]
If \(V\) is one through \(K\), its residual vanishes there and has no
overlap with the early zero-sum changes. The curved reference vanishes
after \(K\), so \(B=0\) and \(R_*=V^2\).
For a nonflat displacement we retain BOTH displayed terms.

\begin{lemma}\label{app:exact19:lem:Dcentretest}
Suppose the reference above makes \(16(G-B)-R_*>H\) sufficient to
contradict \(\|x\|^2<1\).
Either of the following sets of bounds suffices:
\[
\begin{array}{c|c|c|c}
 H\text{ at most}&w_1\text{ at most}&w\cdot P_2\text{ at most}
                                      &w\cdot P_3\text{ at most}\\ \hline
 42.278&.478&.956&1.291\\
 41.539&.54&.956&1.291 .
\end{array}
\]
Then the actual centre opposite is impossible.
\end{lemma}

\begin{proof}
Both ACTUAL fork cosines are less than ten. A fork of cosine at least
ten gives \(V=P_{18}\pm e_{19}\), \(G\ge\sqrt{171/11}>3.94\), \(R=19\),
hence \(16G-R>44.04\).
A \(.50\)-strong centre gives \(V=P_{17}\), \(G>3.7695\), \(R=17\),
hence \(16G-R>43.312\). Both clear either threshold.
We may classify the centre under its actual \(.50\)-weak hypothesis.

The complete positive finite-white list and all direct, affine and
reciprocal calculations are explicitly proved in
Section~\ref{app:exact19:Dfivecentrelist}: isolation; left-only three at
one, two, fifteen, sixteen or four at one/sixteen; one arbitrary finite
fork; fork pairs \((3,3),(3,4),(3,5)\); and left endpoint-three joined
to one fork-three/four or both fork-threes.
Each formula depends on its ACTUAL old contacts, not the head.
Own-three at one fork has reciprocal cosine
\((17+19\beta)/2\ge53/4\); own-zero at a \((3,4)\) pair has reciprocal
cosine \(17/2+(19/4)\sqrt2>10\).
The positive inverse proves these are actual endpoints even when a
centre and fork opposite coincide.
The affine kernels there preserve every original wall term by term.
Thus all own-zero/three/four/six/parallel, mixed-left and direct-ratio
decisions remain valid EXCEPT own-zero with one finite fork, which was
separately removed for native five.

We prove that case without any head-specific obstruction.
Let its fork be \(f\), the other fork \(h\), and
\(c=Ap-\beta\omega_f\), \(A^2=(19/2)\beta^2-1\).
Fork labels three, four and six have the positive actual nulls
\[
\begin{split}
 &2c+4n_f+6n_{17}+5n_{16}+4n_{15}
                         +3n_{14}+2n_{13}+n_{12}+3n_h,\\
 &\sqrt2c+2n_f+2n_{17}+n_{16}+n_h,\\
 &\sqrt3c+2n_f+n_{17}.
\end{split}
\]
They are future through \(c\), nonpositive against every omitted wall,
so compactness excludes them. At label five the ACTUAL \(H_4\) on
\(c,n_f,n_{17},n_{16}\) has positive correction
\(\beta(3n_f/2+n_{17}+n_{16}/2)\), residual \((5-3\varphi)/8\).
The finite projected POINT gives
\[
 V=P_{15}+(e_{16}+e_{17}+e_{18}\pm e_{19})/2,\quad R=16,\quad
 G>\sqrt{3+8\varphi}-\varphi^{-3}>3.75 .
\]
Indeed \(1.618<\varphi<1.619\) gives first root \(>3.992\) and
\(\varphi^{-3}<.238\). Thus \(16G-R>44\).
Both signed orientations are included.

For every finite fork label at least seven put \(v=\beta^{-2}\).
The positive actual core
\[
                  c+\beta^{-1}n_f+2(\beta^{-1}-\beta)n_{17}
\]
has zero products with \(c,n_f\), negative products with
\(n_{17},n_{16},n_h\), and square
\(- (1-\beta^2)(4\beta^2-3)/\beta^2<0\).
All unknown-wall products are nonpositive and its time is positive.
Its future normalization gives
\[
\begin{gathered}
 V=P_{16}+(3-2v)e_{17}+(v-1)e_{18}\pm(v-1)e_{19},\\
 G=\sqrt{19-2v}-\sqrt{-8+14v-6v^2},\quad
 R=27-16v+6v^2,\quad 1<v<100/81 .
\end{gathered}
\]
The line \(4.364-.242v\) lies below the first root: its endpoint square
checks are \(4.122^2<17\) and \(329.284^2<1339(81)\), and concavity
puts the chord below the root. With \(z=v-7/6\), the other root is
\(\sqrt{1/6-6z^2}\le1/\sqrt6-3\sqrt6z^2\), and
\(R=16.5-2z+6z^2\). Hence
\[
 16G-R>42.2746-1.872z+111.57z^2>42.2646 .
\tag{D19forkcontinuous}\label{app:exact19:eq:Dforkcontinuous}
\]
The completion uses \(1.872^2<.04(111.57)\).
This bound alone is too weak for the first threshold; we do not round
it up. Native seven has \(.064<z<.067\), giving \(>42.60\); every
native label at least nine has \(z<-.033\), giving \(>42.45\).
For native eight, \(v=4-2\sqrt2\) gives
\[
\begin{split}
 G&=\sqrt{11+4\sqrt2}-\sqrt{68\sqrt2-96}>3.6732094,\\
 R&=107-64\sqrt2<16.490333,\qquad
                          16G-R>42.2810174>42.278 .
\end{split}
\]
The positive square seals are
\(4.0812809^2<11+4(1.414213562)\),
\(.4080715^2>68(1.414213563)-96\).
These cases exhaust all finite labels without arithmeticity or the
native-five finite-fork obstruction.

The remaining point exceptions are already constructed, with all
original-wall checks, in Sections~\ref{app:exact19:Dfiveisolatedcentre}
and~\ref{app:exact19:Dfivecentrepoints}.
We use their point identities and scalar inequalities, not their
head-paid conclusions. Their complete weaker bounds are
\[
\begin{array}{l|r}
 \text{actual centre profile}&16G-R\text{ greater than}\\ \hline
 \text{isolated five}&45.56\\
 \text{isolated seven/eight}&43.432\\
 \text{isolated finite at least nine}&44.188\\
 \text{isolated dotted}&42.328\\
 \text{near-left-three, own five}&43.24\\
 \text{near-left-three, own seven/eight}&44.538\\
 \text{one/two fork-three/four, own five}&46.84\\
 \text{one fork-four, larger finite/dotted own}&43.992\\
 \text{one fork-three, larger finite/dotted own}&43.296\\
 \text{two fork-threes, own four}&43.647 .
\end{array}
\]
These vectors are flat through fourteen, hence have no early head cross.
For isolated dotted, for example, the actual core gives
\[
 G=\sqrt{17-u/2}-\sqrt{u(1-u)/2}>4.062-.354=3.708,\quad
 R=17-u+u^2/2\le17,\quad 0<u<1 .
\]
Near-left own five uses the finite pair's projected POINT
\(AB/(B+\sqrt{1-\gamma^2})\), not the smaller unprojected
\(A-\sqrt{1-\gamma^2}\).
The one-fork continuous intervals \(.9<\gamma<1.42\) and
\(.9<\gamma<1.01\) retain the polynomial/core proofs of
Section~\ref{app:exact19:Dfivecentrepoints}, including every dotted
parameter left by weakness.
All corrections are positive actual-root combinations or positive
finite inverse projections, so neither signed fork nor any unknown wall
is discarded.

Only far-left-three remains.
The finite point identities of
Section~\ref{app:exact19:Dfivefarcentre} are local to
\(v=\gamma\omega_{17}+\omega_1/2\), not to the head:
\[
\begin{array}{c|c|c}
 \text{own}&V=P_{16}+aP_l+(e_{17}+e_{18})/2&G\text{ greater than}\\ \hline
 5&l=3,\ a<5/24&4.1158\\
 7,8&l=2,\ a<5/18&3.963\\
 \text{finite }\ge9&l=1,\ a<.533&3.953 .
\end{array}
\]
They are the actual \(H_4\) point, the negative future core
\(c+\gamma n_{17}+n_1/2\), and the positive pair projection.
Their normalized losses are \(<.097125,<.25,<.26\) from a first term
\(>\sqrt{17.75}>4.213\).
For these NONFLAT displacements retain
\[
                    B=a\,w\cdot P_l,\qquad R_*=16+la^2+1/2 .
\]
Under the first coordinate row the three tests are greater than
\(44.91,42.50,42.387527\); the last is exactly the safe expression
\(16(3.953-.533(.478))-(16.5+.533^2)\).
Under the second coordinate row the first two bounds still apply,
and the last becomes
\[
             16(3.953-.533(.54))-(16.5+.533^2)=41.858791>41.539 .
\]
Every finite far own is removed with its actual head cross.

All dotted \(\gamma\ge11/8\) are directly \(.50\)-strong because
\((\gamma+1/34)/\sqrt{9/17}>1.928\).
For \(1<\gamma\le10/9\), put \(a=(2\gamma)^{-1}\in[.45,.5)\).
The actual core \(c+\gamma^{-1}n_{17}\) gives
\[
\begin{split}
 V&=P_{16}+aP_1+(1-2a^2)e_{17}+2a^2e_{18},\\
 G&=\sqrt{17+2a-a^2}-a\sqrt{2-8a^2},\qquad
 B=aw_1,\quad R_*\le16.75 .
\end{split}
\]
The function \(G-aw_1\) increases even for \(w_1=.54\): its derivative is
\[
 \frac{1-a}{\sqrt{17+2a-a^2}}
           +\frac{16a^2-2}{\sqrt{2-8a^2}}-w_1>0 .
\]
At \(a=.45\), the positive square checks
\(\sqrt{17.6975}>4.2068,\sqrt{.38}<.6165\) give \(G>3.929375\).
The two coordinate rows yield left sides \(>42.6784\) and \(>42.232\),
respectively. Both suffice.

The ONLY centre left is the literal far-three helper
\[
 c=Ap-\gamma\omega_{17}-\omega_1/2,\qquad
 A^2=34\gamma^2+2\gamma-1/2,\qquad 10/9<\gamma<11/8 .
\]
Sections~\ref{app:exact19:Dfiveq}--\ref{app:exact19:Dfiveqfinish}
prove a local geometric statement from precisely this interval.
Their derivation uses neither a native-five head nor head energy:
the actual adjacent vertex has forest \(A_{17}+A_1+A_1\); its two short
fork edges are classified with their complete finite-white costs,
old \(n_{17}\) acuteness and the actual native gap.
At least one is nonshort.
Comparing its ACTUAL endpoint with the ORIGINAL farthest vertex retains
\(0<\tau\le1\) and gives a displacement flat through seventeen, with
\[
                  G>3.7736,\quad R<17.0703125,\quad
                                16G-R>43.3072875 .
\]
This clears both thresholds in
\eqref{app:exact19:eq:Dgeneralpay}.
The adjacent vertex was never declared farthest.
One original \(x\) pays every pressure, including all early and nonflat
crosses, and the full centre list is exhausted.
\end{proof}

\subsection{Changing the first knot without losing a weak-six omission}
\label{app:exact19:Dmixedfloors}

Put \(f_i=\sqrt i-.58/\sqrt2\), \(\epsilon=.03/\sqrt2\), and
\[
 F=f_5-\epsilon=\sqrt5-.61/\sqrt2,\qquad b=F/5 .
\]
The independent OLD theorem supplies the INDIVIDUAL floor \(S_5>F\).
Useful exact scalar bounds are
\[
\begin{gathered}
 1.80473<F<1.804733,\quad .360946<b<.3609466,\quad f_5/5<.36519,\\
 E_{13}>.9056051,\quad E_{14}>.9241299,\quad
 f_6>2.0393,\quad f_{13}<3.1955,\quad f_{14}<3.3316 .
\end{gathered}
\]
The \(E_K\) are the already proved curved-reference scalar values
in \eqref{app:exact19:eq:BE17}, not numerical enumeration outputs.

Start with the full \(.58\) reference and lower its knot at five by
\(\epsilon\), leaving the actual strong later knots.
The energy change is
\(-2\epsilon\beta_5+(6/5)\epsilon^2\), where
\[
 \beta_5=(6/5)\sqrt5-\sqrt6-.116/\sqrt2<.15177 .
\]
The bounds \(\sqrt5<2.236068,\sqrt6>2.449489,1/\sqrt2>.7071\)
give this seal, and \(\epsilon<.021214\) gives total loss
\[
             2\epsilon\beta_5-(6/5)\epsilon^2
                      <.042428(.15177)-.00054=.00589929756<.0059 .
\]
The new first-five mean is \(b\); the first subsequent slope is at most
\(\sqrt6-\sqrt5+\epsilon<.234637\).
Thus its original \(P_5\) coefficient is \(>.126\).
If six is genuinely strong it is \(>.126309\), while the INDIVIDUAL
\(P_6\) coefficient is \(>.0375\), as used below.

Actual \(.58\)-weak knots at six through fourteen are pure far-three
helpers and nonadjacent. Their ordinary disjoint averaging losses sum
to \(<.000408\). If six is weak, the interval five--seven changes its
loss to \((2\sqrt6-\sqrt5-\sqrt7+\epsilon)^2/2\).
Its EXTRA over the ordinary omission is
\[
 \epsilon(2\sqrt6-\sqrt5-\sqrt7)+\epsilon^2/2
          <.021214(.0172)+.000225=.0005898808 .
\]
Only this interval touches the changed knot. Hence the COMPLETE
omission loss, including weak six, is \(<.0010\).
The mixed reference therefore has
\[
                  q_{13}>.8987051,\qquad q_{14}>.9172299,
                                      \qquad \text{terminal price}>1/8 .
\tag{D19mixedenergy}\label{app:exact19:eq:Dmixedenergy}
\]
If six is strong all omissions begin at seven; the sharper complete
bound is
\[
         1/6912+1/16384+1/32000+1/55296<.000256 .
\]
Omitting later knots only increases the individual capacities just
identified. This calculation makes no assertion that five or six is
new-strong unless its separate geometric argument proves that fact.

\subsection{Native seven: exclude a weak six before making the reference}
\label{app:exact19:Dsevenearly}

Let \(g_0=\cos(\pi/7)\), \(S=g_0^2\), and
\[
 c_H=A_Hp-g_0\omega_1,\qquad A_H^2=2S-1,\qquad u=g_0^{-2}.
\]
The polynomial \(8g_0^3-4g_0^2-4g_0+1=0\), increasing above \(.9\),
gives
\[
 .90096886<g_0<.90096888,\quad 1.2319140<u<1.2319143,\quad
                           .8116<S<.8119,\quad A_H<.7897 .
\]
Expand at \(.901\): value, derivative and quadratic coefficient are
\(.000257608,8.275224,17.624\).
Increments \(-.00003114,-.00003112\) give opposite signs, and
reciprocal-square multiplications give the \(u\) interval.
These are native-angle scalars, not a field census.

The positive core of Section~\ref{app:exact19:longfeet},
\(c_H+g_0^{-1}n_1+2(g_0^{-1}-g_0)n_2\), gives
\[
\begin{split}
 Z&=(1-u/2)(e_1+e_2)+(u-1)e_3,\quad\sum Z_i=1,\\
 x\cdot Z&\ge H=\sqrt{1-u/2}-\sqrt{(7u-3u^2-4)/2}>.4318779,\\
 R&=Z^2=3-4u+3u^2/2<.348763 .
\end{split}
\]
Its selected products are zero or negative, its square negative, and
positive actual-root coefficients check EVERY other wall.
The seals use first root \(>.6197119\), second \(<.187834\), and
\(R(1.2319140)=.348762155094\).

If six were \(.58\)-weak, its actual bulk profile would be
\[
 c_6=Ap-\gamma\omega_6-\omega_1/2,\quad
 A^2=12\gamma^2+2\gamma-1/2,\quad \varphi/2\le\gamma<6/5 .
\]
It is distinct from \(c_H\) by its contact at six.
Its mutual magnitude is \(M_6=A_HA-g_0(2\gamma+1)\).
The squared positive-term difference
\((20S-12)\gamma^2-2\gamma+1/2-2S\) is positive and increasing:
at \(\gamma=.809\) it is greater than
\[
                    4.232(.809)^2-1.618-1.1238=.027963592 .
\]
Also \((12\gamma+1)^2=12A^2+7\) proves \(M_6\) increases, and
\[
                      M_6<.7897(4.38)-.9009(3.4)=.395826<1/2 .
\]
The actual native alphabet has no magnitude strictly between zero and
one half. This excludes ALL finite and dotted weak-six parameters.
Thus six is genuinely strong before any early reference is built;
the complete omission loss now starts at seven and is \(<.000256\).

A weak-five far helper has
\(A^2=10\gamma^2+2\gamma-1/2\) and
\(M_5=A_HA-g_0(2\gamma+1)\).
The derivative is positive, by the squared comparison
\((16S-10)A^2+6A_H^2>0\); and
\[
                         M_5<.7897(4.038)-.9009(3.4)=.1257486<1/2 .
\]
Acuteness and the native gap force \(M_5=0\), hence
\((16S-10)\gamma^2-2\gamma+1/2-2S=0\).
This polynomial increases above \(.809\), is negative at one, and
positive at \(1.04\), so \(1<\gamma<1.04\).
The actual core and weak UPPER bound at five, together with \(x_5>x_6\),
give
\[
 x_1>\sqrt2\{A-\sqrt{1-\gamma^{-2}}-\gamma(\sqrt{10}-.58)\}.
\]
The braces decrease: \(A'<13/4\), the radical derivative is greater
than one, and \(\sqrt{10}-.58>2.5\).
At \(1.04\), \(A>3.52\), the radical is \(<.275\), and
\(\sqrt{10}-.58<2.58228\), yielding \(x_1>.79114420896>.79\).
The legal reference through \((1,.79),(6,f_6),(K,f_K)\) has energy
\[
             .79^2+(2.0393-.79)^2/5+1.156^2/7>1.127 .
\]
Its slopes decrease, and \(K=14\) only raises the final chord energy.
This contradicts the SAME norm. Thus weak five can ONLY be the actual
isolated dotted cap; it has not been silently discarded.

For that cap put \(v=\gamma_5^{-2}\in(0,1)\).
Its actual local first-six-root pressure, proved in the head reduction,
is
\[
 Z_5=P_5+(v/2)(e_6-e_5),\qquad x\cdot Z_5\ge T_5>1.8185 .
\]
Weakness forces \(v>.499\).
For \(v\le.499\) the same core gives loss at most
\[
 \frac{v}{\sqrt{10}+\sqrt{10-v}}+\sqrt{v(1-v)}
       <.16015(.499)+.4999991=.57991395<.58 .
\]
Both upper terms increase there; the denominator seal uses
\(\sqrt{10}>3.1622,\sqrt{9.5}>3.0822\), and
\(.4999991^2>.499(.501)\). Such a cap would be strong.

\subsection{Native seven: pay the head and dotted-five foot jointly}
\label{app:exact19:Dsevenreference}

If five is \(.58\)-strong, take the ordinary curved reference, with
first-five mean \(b=f_5/5<.36519\), omission loss \(<.000256\),
\(P_5\) price \(>.15\) and \(P_K\) price \(>1/8\).
Use \(w=w_0+.45(Z-P_5/5)\), paying \(.09\) at five.
The residual is orthogonal to \(w_0\), and its full effective gain is
\[
                  .9(H-b)-.2025(R-.2)>.0298946025 .
\]
Thus \(L_{13}>.9352437025,L_{14}>.9537685025\).

If five is the retained weak isolated dotted cap, use the OLD floor
\(F\), the ACTUAL strong \(f_6\), and every later actual strong knot.
Section~\ref{app:exact19:Dmixedfloors} gives loss \(<.0059+.000256\),
individual \(P_5\) price \(>.126309\), and \(P_6\) price \(>.0375\).
Pay \(.094\) at five and set \(w_H=w_0+.47(Z-P_5/5)\).
The effective head gain
\(.94(H-F/5)-.2209(R-.2)>.0338136753\)
gives intermediate energies \(.9332627753,.9517875753\).
The REMAINING \(P_5\) price is \(>.0323\), and exactly
\[
 w_H\cdot P_5=F,\quad w_H\cdot P_6=f_6,\quad
                T_5-w_H\cdot Z_5>.013767+.01615v .
\]

Project the second foot onto ONLY the legitimate \(P_5,P_6\) space:
\[
 \Pi Z_5=(1-3v/5)P_5+(v/2)P_6,\qquad
                         Q=\|Z_5-\Pi Z_5\|^2=v^2/5 .
\]
With \(\nu=.046\), its INDIVIDUAL costs are
\(\nu(1-3v/5)<.0322276<.0323\) and \(\nu v/2<.023<.0375\).
No head-induced \(P_3\) coefficient is spent as an old anchor.
The residual is \(v/10\) on coordinates one through four, \(-2v/5\)
at five, zero at six. Its head cross is
\(w_H\cdot(Z_5-\Pi Z_5)=.047v\), not zero.
The ONE completed square includes it, giving extra effective gain
\[
\begin{split}
 2\nu(T_5-w_H\cdot Z_5)-\nu^2Q
  &>.001266564+.0014858v-.0004232v^2\\
  &>.0019026009768>.0019026 .
\end{split}
\]
The polynomial increases on \([.499,1]\).
Final energies are therefore
\(L_{13}>.9351653753,L_{14}>.9536901753\).
The foot is used once, not as an independent norm.

Both final references satisfy
\[
 w_1<.453,\quad w\cdot P_2<.906,\quad w\cdot P_3<1.291,
                                       \qquad \text{remaining }P_K>1/8 .
\]
In the dotted branch the last residual adds at most \(.0046\) to each
of the first four coordinates; the preceding bounds were
\(.4475,.895,1.2709\). The strong-five branch has
\(.4481,.8962,1.2756\).
All changes have zero sum within the first six.
The exact sufficient thresholds in
\eqref{app:exact19:eq:Dgeneralpay}, for \(K=13,14\), are
\[
\begin{array}{c|cc}
 \text{strong-five}&42.27240304&42.26441584\\
 \text{dotted-five}&42.2774159808&42.2694287808 .
\end{array}
\]
The common \(42.278\) and first coordinate row of
Lemma~\ref{app:exact19:lem:Dcentretest} apply.
This excludes the ENTIRE native-seven head branch in both actual states
at five.

\subsection{Native eight: the strong and weak six cases give different prices}
\label{app:exact19:Deightreference}

Let \(s=\sqrt2\), \(g_0=\cos(\pi/8)\), and
\[
 c_H=A_Hp-g_0\omega_1,\quad A_H^2=s/2,\quad
                         u=g_0^{-2}=4-2s .
\]
The same positive actual head core gives
\[
\begin{split}
 Z&=(s-1)(e_1+e_2)+(3-2s)e_3,\quad\sum Z_i=1,\\
 x\cdot Z&\ge H=\sqrt{s-1}-\sqrt{17s-24}>.43955,\\
 R&=23-16s<.372584 .
\end{split}
\]
All-wall feasibility and the future negative square follow from the
local core proof in Section~\ref{app:exact19:longfeet}.
The numerical seals use \(1.414213562<s<1.414213563\), with the
first root \(>.64359\) and second \(<.20404\).
Use ONLY the independent OLD five floor \(F\), and split six by its
ACTUAL \(.58\) status.

If six is strong, the mixed curved reference has omissions beginning
at seven, total \(<.000256\), and \(P_5\) price \(>.126\).
Its energies are \(>.8994491,.9179739\) for \(K=13,14\).
Pay \(.09\) at five and insert \(.45(Z-P_5/5)\).
The orthogonal residual gain is
\[
             .9(.43955-.3609466)-.2025(.372584-.2)=.0357948 .
\]
The final reference therefore satisfies
\[
\begin{gathered}
 L_{13}>.9352439,\quad L_{14}>.9537687,\\
 w_1<.45735,\quad w\cdot P_2<.9147,\quad w\cdot P_3<1.2629 .
\tag{D19eightstrong}\label{app:exact19:eq:Deightstrong}
\end{gathered}
Its remaining \(P_K\) price is \(>1/8\).

If six is weak, its ACTUAL wall is
\[
 c_6=Ap-\gamma\omega_6-\omega_1/2,\quad
 A^2=12\gamma^2+2\gamma-1/2,\quad \varphi/2\le\gamma<6/5 .
\]
Its distinctness from the head is explicit in its six contact.
The mutual magnitude \(M=A_HA-g_0(2\gamma+1)\) is positive and
increasing: the squared positive-term difference is
\((5s-2)\gamma^2-2\gamma-(1+s)/2\), positive already at \(.8\);
the derivative comparison is \((5s-2)A^2+7s/2>0\).
At \(6/5\),
\[
                          M<.8409(4.38)-.9238(3.4)=.542222<.55 .
\]
The actual native gap therefore forces \(M=1/2\).
Since \(M(1)<1/2\) and \(M(1.165)>.502559707\), we have
\(1<\gamma<1.165\).

The actual core and the weak UPPER bound \(S_6\le f_6\), with
\(x_6>x_7\), now give
\[
 x_1>\sqrt2\{A-\sqrt{1-\gamma^{-2}}-\gamma(\sqrt{12}-.58)\}.
\]
The braces decrease: \(A'<15/4\), the radical derivative is \(>1\),
and \(\sqrt{12}-.58>2.8\).
At \(1.165\), use \(A>4.2563\), the radical \(<.5131\), and
\(\sqrt{12}-.58<2.88411\), obtaining \(x_1>.5419382>.54\).
The actual helper nonadjacency makes seven strong.
No strong-six floor was inferred.

Take the legal reference through
\((1,.54),(5,F),(7,f_7)\) and the later actual strong knots.
Its initial slopes satisfy
\[
 m=(F-.54)/4>.3161825,\quad n=(f_7-F)/2<.215485,\quad d_7<.183 .
\]
Thus its INDIVIDUAL prices at \(P_1,P_5\) are \(>.2238,>.1006\).
Relative to the mixed reference with six omitted, the first-coordinate
gain is \((5/4)(.54-b)^2>.04005125\).
Keeping the enlarged total omission budget \(<.0010\), its energies are
\[
                         q_{13}>.93875635,\quad q_{14}>.95728115 .
\]
This accounts for weak six before any head pressure is inserted.

The correct projection is onto BOTH available anchors \(P_1,P_5\):
\[
 \Pi Z=(Z_1-t)P_1+tP_5,\quad t=(2-s)/4,\quad
                          Q=\|Z-\Pi Z\|^2=37/2-13s<.115224 .
\]
Use \(\mu=.3\).
Its separate costs are
\(\mu(Z_1-t)<.080331<.2238\) and \(\mu t<.043935<.1006\).
The old reference is constant only on coordinates two through five;
we do not project onto a fictitious first-five constant block.
Its actual pairing is
\[
                     w_0\cdot Z=.54(s-1)+m(2-s)<.409 ,
\]
so the full head gain is \(>.6(.03055)-.09(.115224)=.00795984\).
The final weak-six reference satisfies
\[
\begin{gathered}
 L_{13}>.94671619,\quad L_{14}>.96524099,\\
 w_1=.54,\quad w\cdot P_2<.9366,\quad w\cdot P_3<1.2603 .
\tag{D19eightweak}\label{app:exact19:eq:Deightweak}
\end{gathered}
It retains the \(P_K\) coefficient \(>1/8\).

In the strong-six case the exact sufficient thresholds in
\eqref{app:exact19:eq:Dgeneralpay} are
\(42.2723904,42.2644032\), so \(H=42.273\) suffices.
In the weak-six case they are \(41.53816384,41.53017664\), so
\(H=41.539\) suffices.
The first and second coordinate rows of
Lemma~\ref{app:exact19:lem:Dcentretest}, respectively, apply.
In particular the weak-six case keeps its larger first coordinate in
every nonflat centre pressure; it never borrows the strong-six bound.
This excludes the ENTIRE native-eight head branch without a
new-strong-five assumption or an independent head energy.

\subsection{Native nine: one mixed reference covers every early omission}
\label{app:exact19:Dninereference}

Put \(g=\cos(\pi/9)\), \(c_H=A_Hp-g\omega_1\),
\(A_H^2=2g^2-1\), and \(u=g^{-2}\).
The polynomial \(8g^3-6g-1=0\), increasing above \(.9\), gives
\[
                    .9396<g<.9398,\qquad 1.1322<u<1.1327 .
\]
The cubic numerators at \(2349/2500,4699/5000\) are
\(-21983608,203952792\), with denominators
\(15625000000,125000000000\).
The intended root is above \(.9\) by the elementary cosine estimate.
Reciprocal checks
\(.9398^2(1.1322)=.999986258088<1\) and
\(.9396^2(1.1327)=1.000002110832>1\)
give the second interval.

The positive head core again supplies
\[
\begin{gathered}
 Z=(1-u/2)(e_1+e_2)+(u-1)e_3,\quad
 H=\sqrt{1-u/2}-\sqrt{(7u-3u^2-4)/2}>.45865,\\
 R=3-4u+3u^2/2<.3941 .
\end{gathered}
\]
Here the first root is \(>.6585\), the second \(<.19985\);
the second radicand is increasing and \(R\) decreasing on the interval.
Their endpoint values \(.039936065,.39401526\) give the stated seals.
The normalized core is future and satisfies EVERY original wall by
the same actual positive-combination proof as before.

Use the mixed reference of
Section~\ref{app:exact19:Dmixedfloors}, including the enlarged weak-six
omission. It uses only the OLD five floor, has
\(q_{13}>.8987051,q_{14}>.9172299\), individual \(P_5\) price \(>.126\),
and terminal price \(>1/8\).
Pay \(.1\) at five and insert \(.5(Z-P_5/5)\).
The full effective gain is
\[
                  (H-b)-.25(R-.2)>.049125 .
\]
Hence
\[
\begin{gathered}
 L_{13}>.9478301,\quad L_{14}>.9663549,\\
 w_1=w_2<.478,\quad w\cdot P_2<.956,\quad w\cdot P_3<1.284 .
\end{gathered}
\]
The coordinate formulas are \(w_1=b+.4-u/4\) and
\(w\cdot P_3=3b+.2\); all other omissions retain the same constant
first-five block before this head change.
The terminal coefficient is unchanged.
The exact sufficient thresholds are \(41.4668736,41.4588864\);
the common \(H=41.467\) satisfies the first row of
Lemma~\ref{app:exact19:lem:Dcentretest}.
Its full centre inventory, fork bound and point tests apply here.
They leave only the literal dotted far helper before its final
adjacent-vertex step. For native nine the following direct estimate
finishes the entire dotted interval with this same head-paid reference.
We do not need the adjacent-vertex step in this branch.
For \(1<\gamma<11/8\) put \(a=(2\gamma)^{-1}\in(4/11,1/2)\).
Its actual core has
\[
\begin{split}
 V&=P_{16}+aP_1+(1-2a^2)e_{17}+2a^2e_{18},\\
 G&=\sqrt{17+2a-a^2}-a\sqrt{2-8a^2},\\
 B&=aw_1,\qquad R_*\le16.75 .
\end{split}
\]
The last bound follows exactly from
\(.75-(1-3a^2+8a^4)=(.25-a^2)(8a^2-1)\ge0\).
The first root is \(>4.1946\), since its square is at least \(2129/121\).
The remaining loss plus head cross is uniformly \(<.55\):
\[
\begin{array}{c|c}
 a\text{ interval}&a\sqrt{2-8a^2}+.478a\text{ less than}\\ \hline
 [4/11,.375]&.354+.478(.375)=.53325\\
 [.375,.4]&.375(.936)+.478(.4)=.5422\\
 [.4,.425]&.4(.849)+.478(.425)=.54275\\
 [.425,.45]&.425(.745)+.478(.45)=.531725\\
 [.45,.5]&.45(.617)+.478(.5)=.51665 .
\end{array}
\]
The radical term decreases for \(a\ge4/11>1/\sqrt8\), so each row
uses its left endpoint, and the linear term its right endpoint.
The first row uses the global bound \(1/(2\sqrt2)<.354\);
the others are positive square checks for
\(\sqrt{.875},\sqrt{.72},\sqrt{.555},\sqrt{.38}\).
Consequently the correct generalized test is
\[
                      16(G-aw_1)-16.75>41.5636>41.467 .
\]
It uses \(R_*=K+\|V-P_K\|^2\), not the unadjusted \(V^2\), and retains
the actual nonzero head cross.
Every dotted parameter is covered, without a head/helper norm being
added independently.

\subsection{Every finite long-head label at least ten}
\label{app:exact19:Dlargefinite}

We use OLD loss \(.61\) throughout this branch.
Put \(f_i=\sqrt i-.61/\sqrt2\), \(b=f_5/5<.361\).
The independent OLD theorem gives \(S_5>f_5\).
For an isolated head with \(g=\cos(\pi/m)\), \(m\ge10\), put
\[
 c_H=A_Hp-g\omega_1,\quad A_H^2=2g^2-1,\quad
                          u=g^{-2},\quad 1<u<1.106 .
\]
The positive core \(c_H+g^{-1}n_1+2(g^{-1}-g)n_2\) has zero products
with the cap and one, negative products with two/three, positive time,
and negative square \(-[7-4g^2-3g^{-2}]\).
Every other original-wall product is nonpositive term by term.
Its normalized pressure is
\[
\begin{split}
 x\cdot Z&\ge h(u)=\{\sqrt{2-u}-\sqrt{P(u)}\}/\sqrt2,\quad
                               P(u)=(u-1)(4-3u),\\
 Z&=(1-u/2)(e_1+e_2)+(u-1)e_3,\quad\sum Z_i=1,\quad
                               R=3-4u+3u^2/2 .
\end{split}
\]
For any actual OLD reference constant \(b\) through five insert
\(.6Y_H\), where \(Y_H=Z-P_5/5\).
Its INDIVIDUAL five price is \(.12\), smaller than the available
\(b-d_5>.3609-.214>.146\).
Orthogonality gives effective gain greater than
\[
                   F(u)=1.2(h(u)-.361)-.36(R-.2).
\tag{D19largegain}\label{app:exact19:eq:Dlargegain}
\]

The function \(F\) decreases on \(1<u\le1.106\).
Here \(P'\ge.364,P<.073\), so the two negative terms in
\[
 h'=-\frac1{2\sqrt2\sqrt{2-u}}-\frac{P'}{2\sqrt2\sqrt P}
\]
have magnitudes \(>.35,>.45\).
Since \(-R'<1\), \(F'<-.96+.36<0\).
At \(u=1.106\), positive squares give
\(\sqrt{2-u}>.9455,\sqrt P<.26888\):
\(P=.072292<.26888^2\), and
\(.4784(1.41422)=.676562848<.9455-.26888\).
Thus \(h>.4784\), \(R-.2=.210854\), and
\[
                                F(u)>.06497256 .
\tag{D19largecommongain}\label{app:exact19:eq:Dlargecommongain}
\]
For every \(m\ge11\) we have \(u<1.087\).
Since \(\pi/11<2/7\), the alternating cosine lower bound gives
\[
       \cos(\pi/11)>1-2/49+2/7203-4/5294205>.9593 .
\]
The polynomial decreases on \([0,2/7]\);
\(.9593^2(1.087)=1.00031880463>1\), and cosine monotonicity covers all
larger labels.
At \(1.087\), \(\sqrt{2-u}>.9555,\sqrt P<.25357\), because
\(P=.064293<.25357^2\) and
\(.496(1.415)=.70184<.9555-.25357\).
Hence \(h>.496\), \(R-.2=.2243535\), and
\[
                           F(u)>.08123274\quad(m\ge11).
\tag{D19largerfinitegain}\label{app:exact19:eq:Dlargerfinitegain}
\]

\paragraph{Late coverage under this actual head.}
The OLD rank-valid bulk supplies an independent strong \(K=11\) or
twelve, with ordinary curved energies
\(q_{11}>.8472225,q_{12}>.8689525\), endpoint price \(>1/7\).
The ONE head payment gives
\[
                         L_{11}>.91219506,\quad L_{12}>.93392506 .
\]
We establish late coverage with these references, not with a late
conclusion from an earlier assumed-weak-five branch.

The complete rank19 finite-white classification and all-wall
projection identities at fifteen/sixteen are printed in
Section~\ref{app:exact19:Doldfive}.
Before its numerical comparisons, those local proofs use only the
retained \(A_{14}+D_4\), respectively \(A_{15}+D_3\), and actual released
root. Its weak-five assumption supplied the earlier reference, not
these identities. We now apply them with the new head-paid reference.
For clarity the full coverage is as follows.

Isolated and left-only rows are \(.61\)-strong except pure far-left
three helpers. The four isolated loss bounds at \(q^2\ge30\), the
near/next-to-end ratios and actual own-zero/three/four/six/parallel
principals in that section are unchanged; native-five uses its ACTUAL
long \(i+1\) dual, at most seventeen.
At right \(D_4\), long-leaf-three has residual \(1/2\);
own-zero is the affine double fork, own-three the triangle, larger
owns have ratio at least \(1+1/\sqrt2\).
The \(D_3\) middle-three has the same decisions.
End-four/own-zero has the printed weighted affine kernel, every
positive own is strong, and end-five is finite \(H_4\), directly strong
even at own-zero.
Both-sided bridges are endpoint-three; own-zero has the long cycle or
middle star, while positive own has numerator \(>3/4\), residual
at most \(5/32\).
Two right contacts make a cycle; there is no foreign component or
exceptional full-rank bridge.

For right fork-three at fifteen or end-three at sixteen, own-zero has
actual affine \(E_8\) arms five/two/one.
Own-three gives the FIVE-cycle at fifteen and FOUR-cycle at sixteen;
four gives the affine star with squared weights \(1/2,1/4,1/4\);
six/parallel have their actual nulls.
At \(D_4\), all remaining finite owns at least seven and dotted owns
are directly strong, with ratio \(>1.626>25/16\).
At \(D_3\), weakness forces a remaining finite \(.9<\gamma<1\);
dotted own is directly strong.
There are just two point exceptions:
\[
\begin{array}{c|c|c|c}
 \text{row}&\text{actual construction}&G>&V^2<\\ \hline
 \text{own five/right three at fifteen or sixteen}
            &\text{positive LEFT }H_4\text{ point}&3.75&16\\
 \text{larger finite own/right end-three at sixteen}
            &c+\gamma^{-1}n_{16}+(2\gamma)^{-1}n_{15}&3.662&16 .
\end{array}
\]
Section~\ref{app:exact19:Doldfive} prints their full point/core
formulas, finite inverse positivity and EVERY-wall check.
They are flat through at least twelve, so have zero head-residual
cross. With the ACTUAL separate endpoint data
\(f_{11}<2.886,f_{12}<3.034\), the legal \(1/7\) payment gives
\[
\begin{split}
 .91219506+2(3.662-2.886)/7-5/49&>1.0318,\\
 .93392506+2(3.662-3.034)/7-4/49&>1.0317 .
\end{split}
\]
The \(H_4\) demand is larger.
Every old-weak late survivor is therefore a literal pure far-left-three
helper, with own five, finite at least seven or dotted, parameter \(<6/5\).

Two adjacent such helpers have mutual \(>-.44>-1/2\), hence are
orthogonal by the actual native gap.
For fifteen/sixteen the retained \(A_{14}\) and ENTIRE right \(D_3\)
give \(I-(14/30)\mathbf1\mathbf1^T\), positive with positive inverse;
fourteen/fifteen retain \(A_{13}+D_4\), cost \(13/28\).
The actual mixed-point calculation printed in
Section~\ref{app:exact19:Doldfive} has prefix inverse \(2\min(i,j)\),
weighted loss \(<7/16+1/(4\sqrt{10})<.61\), and preserves BOTH forks
and all unknown walls.
Weak helpers are thus nonadjacent across the bulk/late join; an ACTUAL
\(J\in\{15,16\}\) is OLD strong under the finite-head hypothesis.

\paragraph{Labels at least eleven.}
Use the ordinary OLD reference ending at that actual \(J\).
At most five nonadjacent omissions cost \(<5/3200\), giving energy
\[
                        >E_{15}-5/3200>.92564-.0015625=.9240775 .
\]
Its five price is still \(>.146\).
Insert the head foot ONCE.
Equation~\eqref{app:exact19:eq:Dlargerfinitegain} yields
\[
                         \|x\|^2>.9240775+.08123274=1.00531024>1 .
\]
The earlier references established coverage; their energies are not
added to this final one.

\paragraph{Native ten.}
Here \(u=2-(2/5)\sqrt5\in(1.1,1.106)\), since
\(447/200<\sqrt5<9/4\).
If sixteen is OLD strong, its ordinary reference has
\(q_{16}>E_{16}-5/3200>.9401975\); the common head gain gives
\(\|x\|^2>1.00517006\).
Otherwise sixteen is the actual pure weak helper and fifteen strong.
Take the reference ending fifteen and insert the single head correction.
It has effective energy \(>.98905006\), remaining \(P_{15}\) coefficient
\(>d_{14}>8/61>.13\), and ACTUAL coordinate bounds
\[
                   w_1<.511,\quad w\cdot P_2<1.022,\quad
                                             w\cdot P_3<1.323 .
\]
They follow from \(b+.6(.8-u/2),2b+.6(1.6-u),3b+.24\).

Write the actual weak helper
\[
 c=Ap-\gamma\omega_{16}-\omega_1/2,\qquad
                              A^2=32\gamma^2+2\gamma-1/2 .
\]
Its finite feet and negative cores give \(x\cdot V\ge G\),
\(V=P_{15}+Y\), with COMPLETE data
\[
\begin{array}{c|c|c|c|c}
 \text{own}&Y&G>&Y^2<&w\cdot Y<\\ \hline
 5&aP_3+(e_{16}+e_{17})/2,\ a<5/24&3.9954&.631&.275625\\
 7,8&aP_2+(e_{16}+e_{17})/2,\ a<5/18&3.8426&.655&5.11/18\\
 \text{finite }\ge9&aP_1+(e_{16}+e_{17})/2,\ a<.533
                                       &3.8326&.785&.272363 .
\end{array}
\]
These are respectively the actual positive \(H_4\) point, future core
\(c+\gamma n_{16}+n_1/2\), and finite pair point.
Their first normalized term is at least \(\sqrt{16.75}>4.0926\);
the losses are respectively \(<.097125,.25,.26\).
For the last, \(\gamma>.939\) and
\(.939^2(1.1352)=1.0009296792>1\).
The displayed norms and overlaps follow by squaring \(Y\) and using
the ACTUAL coordinate bounds, retaining every head/helper cross.
All finite projections are points with positive coefficients in actual
normals, so EVERY original wall is preserved.

For the entire dotted helper range,
\(a=(2\gamma)^{-1}\in(5/12,1/2)\), and
\[
\begin{split}
 Y&=aP_1+(1-2a^2)e_{16}+2a^2e_{17},\quad Y^2\le.75,\\
 G&=\sqrt{16+2a-a^2}-a\sqrt{2-8a^2},\quad
                                              G-w\cdot Y>3.5429 .
\end{split}
\]
Indeed \(.75-Y^2=(.25-a^2)(8a^2-1)\ge0\).
The function \(G-.511a\) increases, since the second term in \(G'\)
is \((16a^2-2)/\sqrt{2-8a^2}>7/9>.511\), while the first is positive.
At \(a=5/12\),
\[
 G-.511a=\sqrt{2399}/12-5\sqrt{22}/72-.511(5/12)
                    >4.0816-.325724-.212917=3.542959 .
\]
Here \(4.0816^2<2399/144\) and \(\sqrt{22}<4.69042\).
This bounds the ACTUAL overlap, without a separately invented pressure
on \(Y\).

Pay \(\nu=.13\) from the INDIVIDUAL \(P_{15}\) coefficient and use
the ONE vector \(w+\nu Y\).
Its full square, including the already paid head, gives
\[
              \|x\|^2>.98905006+.26(G-f_{15}-w\cdot Y)-.0169Y^2 .
\]
Since \(f_{15}<3.442\), the finite/dotted deficits are respectively
greater than \(.2777,.1167,.1182,.1009\); their extra gains exceed
\(.0615381,.0192725,.0174655,.013559\).
Every case satisfies
\(\|x\|^2>.98905006+.013559=1.00260906>1\).
Own-three/four/six and parallel were already excluded by actual
affine/null principals, so this helper list is complete.
Native ten and every finite \(m\ge10\) are excluded, without borrowing
late floors from a different hypothesis.

\subsection{A dotted head: one exact parabola on its whole parameter interval}
\label{app:exact19:Ddottedhead}

Let the isolated actual head be
\[
 c_H=A_Hp-g\omega_1,\quad g>1,\quad A_H^2=2g^2-1,\quad
                                       u=g^{-2}\in(0,1).
\]
The ACTUAL positive core \(c_H+g^{-1}n_1\) has zero cap product,
negative products with one/two, negative square \(u-1\), positive time.
Positivity of its coefficients checks EVERY other original wall.
Its future normalization gives
\[
\begin{split}
 x\cdot Z&\ge H=\sqrt{1-u/2}-\sqrt{u(1-u)/2},\\
 Z&=(1-u/2)e_1+(u/2)e_2,\quad\sum Z_i=1,\qquad R=1-u+u^2/2 .
\end{split}
\]
There is no second-neighbour correction of either sign.
We treat the WHOLE dotted interval, not a selected compact subinterval.

Put \(s=\sqrt6\), \(z=u-2/3\), and
\(q=s/3-sz/8-9z^2/100\).
Here \(-2/3<z<1/3\), and \(q\ge s/4-.04>0\).
The exact identity
\[
 1-u/2-q^2
   =z^2\{3s/50-3/32-(9s/400)z-81z^2/10000\}
\]
is nonnegative: its bracket is at least
\(.045s-.09735>0\), using \(|z|\le2/3\), \(s>2.4494\).
Thus the first root is at least \(q\).
For \(r=\sqrt{u(1-u)/2}\), the elementary square
\((r-1/3)^2\ge0\) gives \(r\le1/3-z/4-3z^2/4\).
Consequently
\[
 H\ge(s-1)/3+(2-s)z/8+.66z^2,\qquad
 R=5/9-z/3+z^2/2 .
\tag{D19headparabola}\label{app:exact19:eq:Dheadparabola}
\]
This is a global square comparison, not a sampled minimum.

Return to mixed floors \(f_i=\sqrt i-.58/\sqrt2\),
\(F=\sqrt5-.61/\sqrt2\), \(b=F/5\).
Section~\ref{app:exact19:Dmixedfloors} includes ALL actual omissions,
including enlarged weak six, and supplies
\[
 q_{13}>.8987051,\quad q_{14}>.9172299,\quad
                  P_5\text{ price}>.126,\quad P_K\text{ price}>1/8 .
\]
Weak isolated dotted five is retained through its OLD individual floor;
neither five nor six has been upgraded to \(.58\)-strong.
Use ONE vector \(w=w_0+.34(Z-P_5/5)\), consuming only \(.068\)
of the actual five coefficient.
The effective head gain is
\[
                       J=.68(H-b)-.1156(R-.2).
\tag{D19dottedgain}\label{app:exact19:eq:Ddottedgain}
\]
Using \eqref{app:exact19:eq:Dheadparabola}, \(b<.361\), and
\(2.4494897<s<2.4494898\), the constant/linear/quadratic terms give
\[
                    J>.0419687+A_Jz+.391z^2,\quad 0<A_J<.0004 .
\]
The exact linear coefficient is \(A_J=.208533333\ldots-.085s\).
Completing the square gives
\[
                                  J>.04196 ,
\tag{D19dotteduniform}\label{app:exact19:eq:Ddotteduniform}
\]
since \(.0004^2<.00000011(4(.391))\).
The ACTUAL coordinates retain their correlation with \(u\):
\[
\begin{gathered}
 w_1=b+.34(.8-u/2),\quad .462946<w_1<.633,\\
 w\cdot P_2=2b+.204<.926,\qquad w\cdot P_3=3b+.136<1.219 .
\tag{D19dottedcoordinates}\label{app:exact19:eq:Ddottedcoordinates}
\end{gathered}
Every change has zero sum in the first five and preserves the INDIVIDUAL
\(P_K\) coefficient. No artificial old \(P_2/P_3\) pressure is paid.

\subsection{The dotted-head centre: retain correlation in the nonflat row}
\label{app:exact19:Ddottedcentre}

The general square \eqref{app:exact19:eq:Dgeneralpay} applies with
the actual \(L_K>q_K+J\).
The two threshold upper bounds are
\[
                     44.6108736-64J,\qquad44.6028864-64J .
\]
By \eqref{app:exact19:eq:Ddotteduniform}, the common sufficient test is
\[
                              16(G-B)-R_*>41.926 .
\tag{D19dottedtest}\label{app:exact19:eq:Ddottedtest}
\]
An actual fork of cosine at least ten has test \(>44.04\);
a \(.50\)-strong centre has test \(>43.312\).
Thus both fork cosines are less than ten and the centre is \(.50\)-weak.

The complete finite-white classification, actual affine/nulls,
reciprocal endpoint computations and flat point proofs of
Lemma~\ref{app:exact19:lem:Dcentretest} apply before its nonflat
coordinate tests. They use the head only through the paid threshold.
Every flat point bound is at least \(42.328\).
Own-zero/single-fork with finite label at least seven has the COMPLETE
continuous estimate \eqref{app:exact19:eq:Dforkcontinuous},
\(>42.2646\), sufficient now without discrete refinements;
three/four/six have the actual nulls and five the positive \(H_4\) point.
Thus every isolated, near, fork-only, mixed bridge and reciprocal
exception is excluded. The actual signed forks and unknown-wall checks
are those printed in that shared proof; no native-five-only fork theorem
is used.

The remaining far-left-three wall has
\[
 c=Ap-\gamma\omega_{17}-\omega_1/2,\quad
                                  A^2=34\gamma^2+2\gamma-1/2 .
\]
Its complete finite feet are those constructed in the shared proof:
\[
 V=P_{16}+aP_l+(e_{17}+e_{18})/2,\quad
                       B=a\,w\cdot P_l,\quad R_*=16+la^2+1/2 .
\]
Own-five and seven/eight have respectively
\((l,a,G)=(3,<5/24,>4.1158)\) and \((2,<5/18,>3.963)\).
The ACTUAL coordinate bounds \eqref{app:exact19:eq:Ddottedcoordinates}
give tests \(>45.15,42.638\), sufficient in
\eqref{app:exact19:eq:Ddottedtest}.

For finite own at least nine, \(l=1,a<.533,G>3.953\).
A separately maximized head coordinate would be unnecessarily weak.
Retain the ACTUAL correlation: its generalized test is greater than
\[
 S_9=16(3.953-.533w_1)-(16.5+.533^2)=46.463911-8.528w_1 .
\]
Subtract the LARGER unrounded threshold above and substitute the
actual \(J,w_1\), using only \(b<.361\):
\[
\begin{split}
 S_9-(44.6108736-64J)&>D(u),\\
 D(u)&=43.52H-7.3984R+1.44976u-17.7762266 .
\end{split}
\tag{D19correlated}\label{app:exact19:eq:Dcorrelated}
\]
This is the complete nonzero head/centre cross calculation.
The global parabola \eqref{app:exact19:eq:Dheadparabola} gives
\[
 D(u)>.1073+A_Dz+25.024z^2,\quad
 A_D=14.795893333\ldots-5.44s,\quad0<A_D<1.471 .
\]
The same two positive square bounds for \(\sqrt6\) give the constant
and linear seals; the quadratic coefficient is exactly \(25.024\).
Therefore
\[
                D(u)>.1073-\frac{A_D^2}{4(25.024)}
                              >.1073-.022=.0853 .
\]
The exact last check is
\(1.471^2=2.163841<.022(100.096)=2.202112\).
This ENTIRE finite-at-least-nine row violates the actual completed
square; it need not clear the coarser standalone test41.926.

For a dotted far helper with \(1<\gamma\le10/9\), put
\(a=(2\gamma)^{-1}\in[.45,.5)\).
Its same actual core gives
\[
\begin{split}
 V&=P_{16}+aP_1+(1-2a^2)e_{17}+2a^2e_{18},\\
 G&=\sqrt{17+2a-a^2}-a\sqrt{2-8a^2},\quad B=aw_1,\quad R_*\le16.75 .
\end{split}
\]
Since \(w_1<.633\), \(G-aw_1\) increases on the whole interval:
the first derivative term is positive, and
\((16a^2-2)/\sqrt{2-8a^2}-w_1\ge1.24/.617-.633>0\).
At \(.45\), \(G>3.929375\), as proved in the shared centre lemma.
The generalized test is thus greater than
\[
                                S_{\rm dot}=46.12-7.2w_1 .
\]
Pointwise it dominates \(S_9\), since
\[
 S_{\rm dot}-S_9=-.343911+1.328w_1
                              >-.343911+1.328(.462946)>0 .
\]
The correlated contradiction therefore covers every small dotted far
helper without substituting unrelated maxima.

Every \(\gamma\ge11/8\) is directly \(.50\)-strong.
The ONLY centre left is the literal interval \(10/9<\gamma<11/8\).
The head-independent actual-q proof in
Sections~\ref{app:exact19:Dfiveq}--\ref{app:exact19:Dfiveqfinish}
gives a flat-through-seventeen pressure with
\(16G-R>43.3072875>41.926\).
Its \(1/8\) terminal payment is still available and its residual has
zero cross with the head-paid support.
The original slack \(0<\tau\le1\), both actual fork orientations and
actual other endpoint are retained; the adjacent \(q\) is never called
farthest.
Equation~\eqref{app:exact19:eq:Dgeneralpay} contradicts the ONE original
norm. Every \(g>1\) dotted-head parameter is excluded.

\subsection{The whole classical alternatives at rank nineteen}
\label{app:exact19:wholeclassicalconclusion}

\begin{proposition}\label{app:exact19:prop:wholeD19}
At the original farthest vertex of a compact hyperbolic Coxeter
\(19\)-polytope the finite forest cannot be a whole \(D_{19}\).
\end{proposition}
\begin{proof}
Section~\ref{app:exact19:Dhead} first excludes every nonisolated head,
then isolated own-zero/three/four/six and parallel.
Its remaining head list is disjoint and exhaustive: native five,
seven, eight, nine, finite label at least ten, or dotted.
The independent OLD five floor in
Section~\ref{app:exact19:Doldfive} relies on none of these later exclusions.

Proposition~\ref{app:exact19:prop:Dfive} excludes native five.
Sections~\ref{app:exact19:Dsevenearly}--\ref{app:exact19:Dsevenreference}
exclude native seven in both actual states at five.
Section~\ref{app:exact19:Deightreference} excludes native eight in both
actual states at six.
Section~\ref{app:exact19:Dninereference} excludes native nine while
retaining every weak-six omission.
Section~\ref{app:exact19:Dlargefinite} proves its own fresh late
coverage and excludes the entire finite-label tail.
Sections~\ref{app:exact19:Ddottedhead}--\ref{app:exact19:Ddottedcentre}
exclude the entire continuous dotted interval with correlated payments.
Each branch uses actual walls at this SAME rank, one ORIGINAL farthest
vector, every original-wall check and legal individual prices.
No face is presumed Coxeter, no adjacent vertex presumed farthest,
and no arithmeticity or stronger-bound census enters.
The complete head partition is exhausted.
\end{proof}

Together with Propositions~\ref{app:exact19:prop:wholeA19}
and~\ref{app:exact19:prop:wholeB19}, this excludes all three whole
classical alternatives at nineteen.
Proper forests with spectators and the necessary forest cover have
separate proofs; this appendix does not replace those obligations.
Its conclusions enter that SAME-rank composition, not a rank-decrement
principle.










